\PassOptionsToPackage{unicode}{hyperref}
\PassOptionsToPackage{hyphens}{url}
\PassOptionsToPackage{dvipsnames,svgnames,x11names}{xcolor}
\documentclass[10pt,letterpaper]{article}

\usepackage{xcolor}
\usepackage[margin=2.0cm]{geometry}
\usepackage{lmodern}
\usepackage{amsmath,amssymb,amsfonts,mathtools}
\usepackage{textcomp}
\usepackage{microtype}
\usepackage{parskip}
\usepackage{longtable,booktabs,array,calc}
\usepackage{caption}
\captionsetup[table]{skip=6pt,font=small,labelfont=bf}
\captionsetup[figure]{font=small,labelfont=bf}
\newcounter{none}
\usepackage{etoolbox}
\usepackage{graphicx}
\usepackage{fancyvrb}
\newcommand{\VerbBar}{|}
\newcommand{\VERB}{\Verb[commandchars=\\\{\}]}
\DefineVerbatimEnvironment{Highlighting}{Verbatim}{commandchars=\\\{\},fontsize=\footnotesize}
\RecustomVerbatimEnvironment{verbatim}{Verbatim}{fontsize=\footnotesize}
\newenvironment{Shaded}{}{}
\newcommand{\AlertTok}[1]{\textcolor[rgb]{0.80,0.00,0.00}{\textbf{#1}}}
\newcommand{\AnnotationTok}[1]{\textcolor[rgb]{0.38,0.63,0.69}{\textbf{\textit{#1}}}}
\newcommand{\AttributeTok}[1]{\textcolor[rgb]{0.49,0.56,0.16}{#1}}
\newcommand{\BaseNTok}[1]{\textcolor[rgb]{0.25,0.63,0.44}{#1}}
\newcommand{\BuiltInTok}[1]{\textcolor[rgb]{0.00,0.50,0.00}{#1}}
\newcommand{\CharTok}[1]{\textcolor[rgb]{0.25,0.44,0.63}{#1}}
\newcommand{\CommentTok}[1]{\textcolor[rgb]{0.38,0.63,0.69}{\textit{#1}}}
\newcommand{\CommentVarTok}[1]{\textcolor[rgb]{0.38,0.63,0.69}{\textbf{\textit{#1}}}}
\newcommand{\ConstantTok}[1]{\textcolor[rgb]{0.53,0.00,0.00}{#1}}
\newcommand{\ControlFlowTok}[1]{\textcolor[rgb]{0.00,0.44,0.13}{\textbf{#1}}}
\newcommand{\DataTypeTok}[1]{\textcolor[rgb]{0.56,0.13,0.00}{#1}}
\newcommand{\DecValTok}[1]{\textcolor[rgb]{0.25,0.63,0.44}{#1}}
\newcommand{\DocumentationTok}[1]{\textcolor[rgb]{0.73,0.13,0.13}{\textit{#1}}}
\newcommand{\ErrorTok}[1]{\textcolor[rgb]{0.80,0.00,0.00}{\textbf{#1}}}
\newcommand{\ExtensionTok}[1]{#1}
\newcommand{\FloatTok}[1]{\textcolor[rgb]{0.25,0.63,0.44}{#1}}
\newcommand{\FunctionTok}[1]{\textcolor[rgb]{0.02,0.16,0.49}{#1}}
\newcommand{\ImportTok}[1]{\textcolor[rgb]{0.00,0.50,0.00}{\textbf{#1}}}
\newcommand{\InformationTok}[1]{\textcolor[rgb]{0.38,0.63,0.69}{\textbf{\textit{#1}}}}
\newcommand{\KeywordTok}[1]{\textcolor[rgb]{0.00,0.44,0.13}{\textbf{#1}}}
\newcommand{\NormalTok}[1]{#1}
\newcommand{\OperatorTok}[1]{\textcolor[rgb]{0.40,0.40,0.40}{#1}}
\newcommand{\OtherTok}[1]{\textcolor[rgb]{0.00,0.44,0.13}{#1}}
\newcommand{\PreprocessorTok}[1]{\textcolor[rgb]{0.74,0.48,0.00}{#1}}
\newcommand{\RegionMarkerTok}[1]{#1}
\newcommand{\SpecialCharTok}[1]{\textcolor[rgb]{0.25,0.44,0.63}{#1}}
\newcommand{\SpecialStringTok}[1]{\textcolor[rgb]{0.73,0.40,0.53}{#1}}
\newcommand{\StringTok}[1]{\textcolor[rgb]{0.25,0.44,0.63}{#1}}
\newcommand{\VariableTok}[1]{\textcolor[rgb]{0.10,0.09,0.49}{#1}}
\newcommand{\VerbatimStringTok}[1]{\textcolor[rgb]{0.25,0.44,0.63}{#1}}
\newcommand{\WarningTok}[1]{\textcolor[rgb]{0.38,0.63,0.69}{\textbf{\textit{#1}}}}

\usepackage{newunicodechar}
\newunicodechar{µ}{\ifmmode\mu\else\textmu\fi}
\newunicodechar{°}{\ifmmode^\circ\else\textdegree\fi}
\newunicodechar{·}{\ifmmode\cdot\else\textperiodcentered\fi}
\newunicodechar{²}{\ensuremath{^{2}}}
\newunicodechar{³}{\ensuremath{^{3}}}
\newunicodechar{¹}{\ensuremath{^{1}}}
\newunicodechar{⁰}{\ensuremath{^{0}}}
\newunicodechar{⁴}{\ensuremath{^{4}}}
\newunicodechar{⁵}{\ensuremath{^{5}}}
\newunicodechar{⁶}{\ensuremath{^{6}}}
\newunicodechar{⁷}{\ensuremath{^{7}}}
\newunicodechar{⁸}{\ensuremath{^{8}}}
\newunicodechar{⁹}{\ensuremath{^{9}}}
\newunicodechar{⁺}{\ensuremath{^{+}}}
\newunicodechar{⁻}{\ensuremath{^{-}}}
\newunicodechar{₀}{\ensuremath{_{0}}}
\newunicodechar{₁}{\ensuremath{_{1}}}
\newunicodechar{₂}{\ensuremath{_{2}}}
\newunicodechar{₃}{\ensuremath{_{3}}}
\newunicodechar{₄}{\ensuremath{_{4}}}
\newunicodechar{ₙ}{\ensuremath{_{n}}}
\newunicodechar{ν}{\ensuremath{\nu}}
\newunicodechar{Δ}{\ensuremath{\Delta}}
\newunicodechar{Ω}{\ensuremath{\Omega}}
\newunicodechar{α}{\ensuremath{\alpha}}
\newunicodechar{β}{\ensuremath{\beta}}
\newunicodechar{γ}{\ensuremath{\gamma}}
\newunicodechar{δ}{\ensuremath{\delta}}
\newunicodechar{ε}{\ensuremath{\varepsilon}}
\newunicodechar{θ}{\ensuremath{\theta}}
\newunicodechar{κ}{\ensuremath{\kappa}}
\newunicodechar{λ}{\ensuremath{\lambda}}
\newunicodechar{π}{\ensuremath{\pi}}
\newunicodechar{ρ}{\ensuremath{\rho}}
\newunicodechar{σ}{\ensuremath{\sigma}}
\newunicodechar{τ}{\ensuremath{\tau}}
\newunicodechar{→}{\ensuremath{\rightarrow}}
\newunicodechar{∎}{\ensuremath{\blacksquare}}
\newunicodechar{−}{\ensuremath{-}}
\newunicodechar{∧}{\ensuremath{\wedge}}
\newunicodechar{∫}{\ensuremath{\int}}
\newunicodechar{≈}{\ensuremath{\approx}}
\newunicodechar{≤}{\ensuremath{\leq}}
\newunicodechar{≥}{\ensuremath{\geq}}
\newunicodechar{⊆}{\ensuremath{\subseteq}}
\newunicodechar{∞}{\ensuremath{\infty}}
\newunicodechar{š}{\ifmmode\check{s}\else\v{s}\fi}
\newunicodechar{✓}{\ensuremath{\checkmark}}
\newunicodechar{≠}{\ensuremath{\neq}}
\newunicodechar{↔}{\ensuremath{\leftrightarrow}}
\newunicodechar{≪}{\ensuremath{\ll}}
\newunicodechar{↛}{\ensuremath{\nrightarrow}}
\newunicodechar{⇒}{\ensuremath{\Rightarrow}}
\newunicodechar{∪}{\ensuremath{\cup}}
\newunicodechar{⟨}{\ensuremath{\langle}}
\newunicodechar{⟩}{\ensuremath{\rangle}}
\newunicodechar{・}{\textperiodcentered}

\usepackage{xurl}
\usepackage{bookmark}
\usepackage{hyperref}

\setcounter{secnumdepth}{-\maxdimen}
\providecommand{\tightlist}{\setlength{\itemsep}{1pt}\setlength{\parskip}{0pt}}

\definecolor{openieAmber}{HTML}{D97706}
\definecolor{openieAmberBg}{HTML}{FFFBEB}
\definecolor{openieAmberBorder}{HTML}{F59E0B}
\definecolor{openieSlate}{HTML}{334155}
\definecolor{openieDark}{HTML}{0F172A}
\definecolor{openieCodeBg}{HTML}{F8FAFC}
\definecolor{openieCodeBorder}{HTML}{CBD5E1}
\definecolor{openieBlue}{HTML}{0057B8}

\hypersetup{
  pdftitle={RCE-JEPA: World Coding, with JEPA as Premise},
  pdfauthor={David Charlot, Open Interface Engineering},
  colorlinks=true,
  linkcolor=openieBlue,
  citecolor=openieBlue,
  urlcolor=openieBlue,
  pdfcreator={XeTeX / Tectonic — Open Interface Engineering}
}

\urlstyle{same}
\setlength{\emergencystretch}{3em}
\hfuzz=6pt
\sloppy

\newcommand{\openiefig}[3]{%
  \IfFileExists{#1}{%
    \includegraphics[width=#2\linewidth,height=0.70\textheight,keepaspectratio]{#1}%
  }{%
    \fcolorbox{openieCodeBorder}{openieCodeBg}{%
      \parbox{0.92\linewidth}{\footnotesize\ttfamily Figure asset embedded in archival PDF; interactive/vector source: \url{#3}}%
    }%
  }%
}

\title{\vspace{-1.2em}\textbf{RCE-JEPA: World Coding, with JEPA as Premise}}
\author{\textbf{David Charlot, Open Interface Engineering}\\[0.15em]
{\small Research paper · open topic $\cdot$ Open Interface Engineering $\cdot$ \href{https://research.openie.dev}{research.openie.dev}}}
\date{\small Revision \texttt{5c50a85} (23 November 2026) $\cdot$ Web: \href{https://research.openie.dev/papers/rce-jepa/}{research.openie.dev/papers/rce-jepa/} $\cdot$ Archival LaTeX: \href{https://research.openie.dev/pdfs/latex/rce-jepa.pdf}{/pdfs/latex/rce-jepa.pdf} (\href{https://research.openie.dev/pdfs/latex/rce-jepa.tex}{.tex})}

\begin{document}
\maketitle
\vspace{-0.6em}

% OpenIE Measurement Honesty Callout Box
\noindent\fcolorbox{openieAmberBorder}{openieAmberBg}{%
  \parbox{\dimexpr\linewidth-2\fboxsep-2\fboxrule\relax}{%
    \small\color{openieDark}%
    \textbf{Measurement.} Research paper, open topic, not a journal final. \textbf{\texttt{H2}} (\texttt{0.04904 J / 84.25 J = 0.000582 <= 0.05}, \texttt{1,718$\times$} cheaper), \textbf{\texttt{H15}} (single-shot and bounded 5-turn \textbf{\texttt{H15-Iterative}} closing \textbf{\texttt{40 / 40}} residual tasks with \texttt{0} contradictions plus warm \texttt{DNULLMM1} Gear-0~v2), \textbf{\texttt{FormJEPA-v1}} (\texttt{8.42M} trainable parameters), \textbf{\texttt{FormEBM-v2}} (\texttt{5.85M} parameters), and span-normalized \textbf{\texttt{FormEBM-v3}} (\texttt{6.05M} parameters; \textbf{\texttt{H5 AUROC = 0.9437}}, \textbf{\texttt{H11 = 94.0\%}} strict 5-step monotonicity, OOD \texttt{H15} inversion resolved \texttt{0.2585 $\to$ 0.5278}) on \texttt{1,387} disjoint Forms, plus the 151-task \textbf{\texttt{rce repair} rung-\texttt{f}} audit (\texttt{holds: false}, \texttt{observe} mode) report Apple M5 Max and Linux \texttt{x86-64} \texttt{RAPL} on-chip \texttt{reported\_j} counters; cloud server joules are \texttt{not\_available}. No hardware rail or wall meter was attached (\texttt{measured\_j = None}, \texttt{board\_metered = false}). No fabricated citations.
  }%
}

\vspace{6pt}
% OpenIE Amber Left-Rule Summary Quote
\noindent{\color{openieAmber}\vrule width 3pt}\hspace{8pt}%
\parbox{\dimexpr\linewidth-11pt\relax}{%
  \small\itshape\color{openieSlate}%
  RCE-JEPA (RCE is the Reality Code Engine; JEPA is a Joint Embedding Predictive Architecture) keeps its name so the work stays findable. The subject is World Coding. JEPA is the premise: a learned model predicts a compact summary of a program instead of writing the next word-piece of text. A World of Code is the space of canonical program Forms plus the exact transition and acceptance rules (the Laws of Coding). The proposer slot is any learned model. The checker decides. Hypotheses H1 to H15 are stated in the paper. Measured results now cover H1 (median description-length ratio 9.34 over 185 hand-written programs), H2 (checker-to-model energy ratio 0.000582, or 1,718 times cheaper on Apple M5 Max in on-chip reported\_j), H5 and H11 on unseen disjoint Form shapes via FormEBM-v2 (5.85M parameters: area under the receiver operating characteristic curve (AUROC) 0.8922 >= 0.85, 84.05\% solvable top-1 accuracy, 2.39 times fewer checker calls, and 93.0\% >= 80\% strict 5-step cumulative defect monotonicity at Spearman 0.9931 and 47.09 mJ per candidate core reported\_j) and span-normalized decomposed FormEBM-v3 (6.05M parameters: unseen-shape H5 AUROC 0.9437 >= 0.85, 83.44\% defect-head and 80.98\% composite solvable top-1 accuracy, H11 strict 5-step monotonicity 94.0\% on E\_depth at Spearman 0.9966 and 92.0\% on E\_total at Spearman 0.9846, and zero-shot out-of-distribution H15 multi-proposer pooled AUROC lifted from 0.2585 to 0.5278, 0.6018 on large-span valid versus short invalid candidates, at 232.45 mJ per candidate core reported\_j), H7 fuzzed arm (0 invariant violations among 1,275 mutants), H8 (0 replay mismatches on 10,500 streams across six targets), H14 (deterministic ladder closing 111 of 151 tasks, 73.5\%, at zero model calls, leaving 40 residual tasks), H15 single-shot and bounded checker-feedback H15-Iterative across those 40 residual tasks (17/40 to 34/40 single-shot; 22/40 at k=1 and 27/40 at k=16 for 8.42M FormJEPA-v1; 40/40 closed on Claude Opus 5.5 and Claude Sonnet 5.5 with 0 contradictions under bounded checker feedback plus public standard-library headers, 39/40 on xAI Grok 4.20 Reasoning; 40/40 Gear-0 replay at 0 calls and 0.325 microseconds per lookup; and 200/200 Zipf(0.85) 3-Gear stream admissions with 16.67 times fewer frontier episodes), and the 151-task rce repair rung f promotion audit (3/119 = 2.5\% > 2.0\%, 0 contradictions, keeping rung f in observe mode). Full intent-to-Form synthesis (H3), cross-module signature-and-trace out-of-distribution energy calibration beyond 0.5278, and external wall-meter energy (measured\_j) remain open.
}
\vspace{6pt}

{
\small
\setcounter{tocdepth}{2}
\tableofcontents
}
\vspace{4pt}
\hrule
\vspace{8pt}

\textbf{Status.} This paper states a thesis, fifteen hypotheses (H1 to
H15), six falsifiers and a programme of work, accompanied by interactive
instruments at
\href{https://research.openie.dev/living/rce-jepa/}{/living/rce-jepa/}.
Its measured results now cover the Phase 0 deterministic floors (the
Form corpus and H1, the adversarial fuzzed arm of H7, the six-target
determinism streams of H8, and the deterministic repair ladder of H14),
the local Apple M5 Max checker-versus-model \texttt{reported\_\allowbreak{}j} ratio
\textbf{H2}
(\texttt{C\_check\ /\ E\_model\ =\ 0.000582\ \textless{}=\ 0.05},
\texttt{1,718} times cheaper), single-shot and bounded checker-feedback
\textbf{H15} across the \textbf{40 open residual tasks} (\texttt{17/40}
to \texttt{34/40} single-shot; \textbf{\texttt{40/40}
(\texttt{100.0\%})} under bounded \texttt{\textless{}=\ 5}-turn checker
feedback plus public standard-library signature headers in
\texttt{H15-\allowbreak{}Iterative} on \texttt{Claude\ Opus\ 5.\allowbreak{}5} and
\texttt{Claude\ Sonnet\ 5.\allowbreak{}5} with \texttt{0} contradictions,
\texttt{39/40} on \texttt{xAI\ Grok\ 4.\allowbreak{}20\ Reasoning}, and
\textbf{\texttt{40/40} zero-call replay} at \texttt{0.325\ µs} /
\texttt{1.77\ µJ} from memory-mapped Gear-0), the purpose-trained latent
Form JEPA repair proposer \textbf{\texttt{FormJEPA-\allowbreak{}v1}} (\texttt{8.42M}
trainable parameters over disjoint Form S-expressions, admitting
\texttt{177/200} held-out synthetic defects and \texttt{27/40} residual
tasks at \texttt{k\ =\ 16}, cutting recurring \texttt{Zipf(\allowbreak{}0.\allowbreak{}85)} Gear-2
episodes from \texttt{200} to \texttt{12}, \texttt{16.67×}, with
\texttt{200/200} stream admissions when paired with
\texttt{H15-\allowbreak{}Iterative}), the contrastive multi-step margin energy-based
model (EBM) scorer \textbf{\texttt{FormEBM-\allowbreak{}v2}} (\texttt{5.85M}
trainable parameters on the identical strictly disjoint
\texttt{1,387}-Form split, lifting unseen-shape \textbf{H5} area under
the receiver operating characteristic curve (AUROC) from \texttt{0.618}
to \textbf{\texttt{0.8922\ \textgreater{}=\ 0.85}} with \texttt{84.05\%}
solvable top-1 accuracy and \texttt{2.39×} fewer checker calls, and
lifting unseen-chain \textbf{H11} strict 5-step edit monotonicity from
\texttt{36.0\%} to \textbf{\texttt{93.0\%\ \textgreater{}=\ 80\%}} at
Spearman \texttt{0.9931} and \texttt{47.09\ mJ}/\allowbreak{}candidate core
\texttt{reported\_\allowbreak{}j}), the span-normalized two-head decomposed scorer
\textbf{\texttt{FormEBM-\allowbreak{}v3}} (\texttt{6.05M} trainable parameters on the
identical split, raising unseen-shape \textbf{H5} AUROC to
\textbf{\texttt{0.9437\ \textgreater{}=\ 0.85}} with \texttt{83.44\%}
defect-head and \texttt{80.98\%} composite solvable top-1 accuracy,
raising unseen-chain \textbf{H11} strict 5-step monotonicity to
\textbf{\texttt{94.0\%}} on \texttt{E\_depth} at Spearman
\texttt{0.9966} and \textbf{\texttt{92.0\%}} on \texttt{E\_total} at
Spearman \texttt{0.9846}, and eliminating Wave-4 cross-task edit-span
inversion on zero-shot out-of-distribution (OOD) \textbf{H15} by lifting
multi-proposer pooled AUROC from \texttt{0.2585} to
\textbf{\texttt{0.5278}} and large-span valid versus short invalid AUROC
to \textbf{\texttt{0.6018}} at \texttt{232.\allowbreak{}45\ mJ}/candidate core
\texttt{reported\_\allowbreak{}j}), and the pre-registered 151-task
\textbf{\texttt{rce\ repair} rung \texttt{f}} promotion audit
(\texttt{3/119\ =\ 2.5\%\ \textgreater{}\ 2.0\%} stochastic share,
\texttt{3/3} passing \texttt{W-MUT-1}, \texttt{0} contradictions,
confirming rung \texttt{f} remains in \texttt{observe} mode). Honest
disclosure: full intent-to-Form synthesis (H3), end-to-end joule parity
(H4), cross-module signature-and-trace OOD energy calibration beyond
\texttt{0.5278} across multi-module git-history rewrites, reference-app
encoding (H6), cost-table prediction (H9), quotient data-saving (H10),
decision-head calibration (H12), discrete-versus-continuous long-horizon
planning (H13), diffusion/evolutionary H15 arms, and external wall-meter
\texttt{measured\_\allowbreak{}j} remain open. JEPA is the premise: predict a summary
of meaning, not the surface text. The subject is World Coding, any
learned proposer over Forms, with the Laws of Coding as the acceptance
rule. The deciding half (commit gate, exact replay, projection, and
integer cost units) is described in
\href{https://research.openie.dev/papers/rce/}{the RCE paper}; this
paper is about the proposer: a learned model that steers search over
Forms while the engine renders every target. Any model in the loop,
hosted or local, only proposes, and security never depends on it.


\begin{center}
\openiefig{fig/rce-jepa-fig01-pipeline.png}{0.94}{https://research.openie.dev/figures/rce-jepa/pipeline.svg}
\end{center}


\emph{Figure 1. The proposed pipeline. Any learned proposer maps stated
intent toward a candidate program. JEPA, which predicts in a latent
space, is the instance drawn. A decoder turns that proposal into
candidate programs in the fixed shape and ranks them by an energy score
(lower means a better fit). A deterministic checker accepts each
candidate or reports that it requires change, and returns the broken
rule for re-ranking. Everything inside the dashed box is learned and
untrusted; only accepted programs reach the engine. Local Apple M5 Max
counters are recorded strictly as \texttt{reported\_\allowbreak{}j}
(\texttt{measured\_\allowbreak{}j\ =\ None}, \texttt{board\_\allowbreak{}metered\ =\ false}), and
remote cloud server energy is recorded as \texttt{not\_\allowbreak{}available}. .}

\textbf{In plain words}

\begin{enumerate}
\def\labelenumi{\arabic{enumi}.}
\tightlist
\item
  The subject is World Coding. A World of Code is the space of program
  Forms, plus exact rules for the next state and for what is accepted
  (the Laws of Coding). JEPA is the premise, not the whole method:
  predict a summary of meaning instead of the next word of code. AI
  (artificial intelligence) here means any learned proposer, not one
  architecture.
\item
  The model proposes and the checker decides. Trace-trained decoders,
  diffusion, reinforcement learning inside a proof checker, evolutionary
  search, library learning and constrained decoding can all sit in that
  slot. The gate does not change when the proposer does. Safety does not
  depend on the proposer being right.
\item
  Measured floors and residual benchmarks stay strictly labeled. Over
  185 hand-written programs the median description-length ratio is 9.34
  (H1); on Apple M5 Max at commit \texttt{1762fb4b}, checking a Form
  costs \texttt{49.04\ mJ} against \texttt{84.25\ J} per local 35B
  candidate generation (\texttt{1,718×} cheaper, supporting H2); fuzzed
  checks found 0 broken invariants in 1,275 accepted mutants (H7);
  replay matched on 10,500 streams across six targets (H8); and the
  deterministic ladder closed 111 of 151 tasks (73.5\%) with zero model
  calls, leaving 40 real-history tasks (H14). Across those 40 residual
  tasks (H15), single-shot proposers admit \texttt{17/40} to
  \texttt{34/40}, bounded checker feedback with public standard-library
  headers (\texttt{H15-\allowbreak{}Iterative}, zero access to the original Form)
  closes \textbf{\texttt{40/40} (\texttt{100.0\%})} on
  \texttt{Claude\ Opus\ 5.\allowbreak{}5} and \texttt{Claude\ Sonnet\ 5.\allowbreak{}5} with
  \texttt{0} contradictions (\texttt{39/40} on
  \texttt{xAI\ Grok\ 4.\allowbreak{}20\ Reasoning}) and crystallizes
  \textbf{\texttt{40/40} into Gear 0 (\texttt{DNULLMM1})} for zero-call
  replay, purpose-trained \texttt{FormJEPA-\allowbreak{}v1} (\texttt{8.42M}
  parameters) admits \texttt{27/40} at \texttt{k\ =\ 16} while
  contrastive \texttt{FormEBM-\allowbreak{}v2} (\texttt{5.85M} parameters: AUROC
  \texttt{0.8922}, \texttt{93.0\%} monotonicity) and span-normalized
  \texttt{FormEBM-\allowbreak{}v3} (\texttt{6.05M} parameters: unseen-shape AUROC
  \texttt{0.9437} on H5, \texttt{94.0\%} \texttt{E\_depth} /
  \texttt{92.0\%} \texttt{E\_total} monotonicity on H11, and zero-shot
  H15 pooled AUROC lifted from \texttt{0.2585} to \texttt{0.5278}) rank
  Form candidates without invoking the checker, and the \textbf{3-Gear
  Cascade} admits \texttt{200/200} recurring \texttt{Zipf(\allowbreak{}0.\allowbreak{}85)}
  arrivals with \texttt{16.67×} fewer frontier episodes (\texttt{12} vs
  \texttt{200}). \texttt{rce\ repair} rung \texttt{f} stays in
  \texttt{observe} mode
  (\texttt{3/119\ =\ 2.5\%\ \textgreater{}\ 2.0\%}, \texttt{0}
  contradictions), and no external wall-meter reading
  (\texttt{measured\_\allowbreak{}j}) is claimed.
\end{enumerate}

\section{Principles}\label{principles}

\begin{itemize}
\tightlist
\item
  The model only proposes.
\item
  The checker decides.
\item
  Security never depends on the model.
\item
  No model runs inside the Reality Code Engine's own toolchain loop.
\item
  Every accepted program form is replayable bit for bit.
\item
  Any decision model, hosted or local, is only a proposer outside the
  gate.
\item
  Deterministic first. The first rule of automation is having a manual
  process that works: a ladder of exact, deterministic proposers runs
  before any model, and the model acts only on what that ladder leaves
  open (stated in the proposer section below).
\end{itemize}

\section{Terms and acronyms used on this
page}\label{terms-and-acronyms-used-on-this-page}

Listed here so that each is defined before it is used. Acronyms that
appear only inside reference entries (venues, report series) are not
listed.

\begin{itemize}
\tightlist
\item
  AGX: NVIDIA's Jetson AGX line of embedded computers
\item
  AILANG: a programming language for AI-written code (a project name)
\item
  AIME: American Invitational Mathematics Examination
\item
  AI: artificial intelligence
\item
  AIXI: Hutter's theoretical model of a universally intelligent agent
\item
  AMD: Advanced Micro Devices
\item
  ANDF: Architecture Neutral Distribution Format
\item
  API: application programming interface
\item
  APU: accelerated processing unit, combining CPU cores and an
  integrated GPU on one package
\item
  ARC: Abstraction and Reasoning Corpus, the task set used by Chollet's
  measure and by the ARC Prize
\item
  AST: abstract syntax tree
\item
  AUROC: area under the receiver operating characteristic curve
\item
  BLAKE3: a cryptographic hash function
\item
  BLEU: bilingual evaluation understudy, a text-overlap score
\item
  CALM: consistency as logical monotonicity
\item
  CEGIS: counterexample-guided inductive synthesis
\item
  CPU: central processing unit
\item
  CUDA: NVIDIA's GPU programming platform
\item
  CWM: Code World Model, an open-weights code model trained in part on
  execution traces
\item
  Dafny: a programming language in which specifications and proofs sit
  beside the code
\item
  DBOS: a durable-execution system (a project name)
\item
  DIR: a directory path given to a command
\item
  DNULLMM1: open-addressed binary memory-mapped hash-table store format
  (64-byte header plus 48-byte slots)
\item
  EB-JEPA: an energy-based JEPA library (a project name)
\item
  EBM: energy-based model
\item
  ECE: expected calibration error
\item
  e-graph: an equality graph, one structure holding many equivalent
  programs
\item
  EMA: exponential moving average
\item
  EWD: the numbered manuscript series of Edsger W. Dijkstra
\item
  ESBMC: efficient SMT-based bounded model checker
\item
  FAIR: Meta's Fundamental AI Research group
\item
  FPR: false positive rate
\item
  GELU: Gaussian error linear unit activation function
\item
  GPT: generative pre-trained transformer, OpenAI's model family
\item
  GPU: graphics processing unit
\item
  GRPO: the reinforcement-learning update named coupled-GRPO in the
  DiffuCoder report
\item
  GSM8K: a grade-school math word-problem set used by LLM-JEPA
\item
  HVM2: the Higher Order Company's interaction-combinator model of
  computation
\item
  IBM: International Business Machines
\item
  ICOT: Institute for New Generation Computer Technology (Japan)
\item
  IEEE: Institute of Electrical and Electronics Engineers
\item
  IQR: interquartile range
\item
  JACM: Journal of the ACM (Association for Computing Machinery)
\item
  LCM: Large Concept Model, a model that predicts in a
  sentence-embedding space rather than in tokens
\item
  LHS: left-hand side of an example event trace
\item
  LLM: large language model
\item
  MBPP: the short Python problem set named in the type-constrained
  decoding report
\item
  MDE: model-driven engineering
\item
  METR: the group that reported the early-2025 developer productivity
  trial (Becker et al.)
\item
  MLA: multi-head latent attention
\item
  MLIPS: million logical inferences per second
\item
  MLIR: multi-level intermediate representation
\item
  MLP: multi-layer perceptron
\item
  MoE: mixture of experts
\item
  MPC: model predictive control, planning by repeated short rollouts
\item
  NAIL: a canonical-form programming language (a project name)
\item
  NVIDIA: a maker of GPUs
\item
  OMG: Object Management Group
\item
  OOD: out-of-distribution
\item
  PAR: pattern-based automatic program repair
\item
  PID: process identifier in the operating system
\item
  POSIX: Portable Operating System Interface
\item
  PTX: NVIDIA's parallel thread execution instruction set
\item
  RAND RM: a RAND Corporation research memorandum
\item
  RAPL: Intel's running average power limit energy counters
\item
  README: a repository's front-page documentation file
\item
  RF: label of an outside-proposal rung-f rule in the RCE specification
\item
  RFC: request for comments, an Internet standards document
\item
  RISC-V: an open reduced instruction set computer architecture
\item
  RJ-F: label of a falsifier of the RCE-JEPA premise
\item
  RL: reinforcement learning
\item
  RMS: root mean square
\item
  ROC: receiver operating characteristic
\item
  ROCm: Radeon Open Compute platform
\item
  SAT: Boolean satisfiability
\item
  SATB: Satiation Benchmark telemetry and harness discipline
\item
  SFT: supervised fine-tuning
\item
  SHA-256: secure hash algorithm with a 256-bit digest
\item
  SMC: System Management Controller
\item
  SoC: system-on-chip
\item
  SONAR: the sentence-embedding space used by Large Concept Models,
  covering 200 languages
\item
  STAN-CS: Stanford computer science technical report series
\item
  STOP: self-taught optimizer
\item
  SUT: system under test
\item
  SWE-bench: a benchmark of real repository issues. SWE-bench Verified
  is a reviewed subset on which several code reports give a
  \href{mailto:pass@1}{\nolinkurl{pass@1}} figure
\item
  TOPLAS: ACM Transactions on Programming Languages and Systems
\item
  TPR: true positive rate
\item
  UNCOL: universal computer-oriented language
\item
  USD: United States dollars
\item
  VAX: Digital Equipment Corporation's virtual address extension
  computers
\item
  Verus: a verifier for a subset of Rust
\item
  VICReg: variance-invariance-covariance regularization
\item
  XML: extensible markup language
\end{itemize}

\section{How this page is organized}\label{how-this-page-is-organized}

The page is meant to be read top to bottom. After the hypotheses and
falsifiers at a glance, it has these parts, in this order:

\begin{enumerate}
\def\labelenumi{\arabic{enumi}.}
\tightlist
\item
  What already exists, and what this page is about (the proposer
  section)
\item
  The thesis, in five parts: Part 1, computation as a finite information
  system; Part 2, compression into Form and template; Part 3,
  kinematics; Part 4, a world model whose physics is known exactly; Part
  5, consequences as testable predictions
\item
  Convergence by commitment
\item
  Section 1. History: seventy years of ``say what you want, get the
  program''
\item
  Section 2. Why the earlier attempts stalled
\item
  Section 3. Three lanes toward a World of Code: training, hybrids, and
  reconceptions of computation
\item
  Section 4. Why now
\item
  Section 5. JEPA as premise, and where it stops
\item
  Section 6. The proposer slot is architecture-neutral
\item
  Section 7. A World of Code
\item
  Section 8. One concrete instance: the Reality Code Engine, with the
  hypotheses table and the falsifiers
\item
  Status and references
\end{enumerate}

\section{Hypotheses at a glance}\label{hypotheses-at-a-glance}

The page tests fifteen hypotheses, H1 to H15, listed here in full before
any of them is discussed. Each is stated with its metric and thresholds
in one table in
\href{https://research.openie.dev/papers/rce-jepa/\#hypotheses}{Section
8}, followed by the falsifiers RJ-F1 to RJ-F6.

\begin{enumerate}
\def\labelenumi{\arabic{enumi}.}
\tightlist
\item
  H1. Form space is much smaller than source space for the same
  programs. (Supported on 185 hand-written programs: median rho 9.34.)
\item
  H2. The checker is cheap relative to generation. (Supported on Apple
  M5 Max in \texttt{reported\_\allowbreak{}j}:
  \texttt{C\_check\ /\ E\_model\ =\ 0.000582\ \textless{}=\ 0.05},
  \texttt{1,718×} cheaper.)
\item
  H3. Prediction in representation space reaches useful acceptance.
\item
  H4. Energy per accepted program falls.
\item
  H5. Energy ranking concentrates mass on valid Forms. (Contrastive
  \texttt{FormEBM-\allowbreak{}v2} reaches AUROC
  \texttt{0.8922\ \textgreater{}=\ 0.85} and \texttt{84.05\%} solvable
  top-1 accuracy, and span-normalized \texttt{FormEBM-\allowbreak{}v3} raises
  unseen-shape AUROC to \texttt{0.9437\ \textgreater{}=\ 0.85} with
  \texttt{83.44\%} defect-head and \texttt{80.98\%} composite solvable
  top-1 accuracy vs \texttt{0.618} for uncalibrated
  \texttt{FormJEPA-\allowbreak{}v1}; zero-shot OOD H15 pooled AUROC rises from
  \texttt{0.2585} to \texttt{0.5278}, while cross-module signature/trace
  calibration beyond \texttt{0.5278} remains open.)
\item
  H6. The reference apps are encodable as Forms.
\item
  H7. Security does not depend on the model. (Fuzzed arm measured: 0
  violations in 1,275 mutants.)
\item
  H8. Every accepted Form replays deterministically. (Supported on
  10,500 streams across six targets: 0 mismatches.)
\item
  H9. Cost is predicted before execution.
\item
  H10. Learning on the quotient (canonical Forms) saves training data.
\item
  H11. Energy is monotone in edit distance. (Contrastive
  \texttt{FormEBM-\allowbreak{}v2} reaches \texttt{93.0\%\ \textgreater{}=\ 80\%}
  strict 5-step monotonicity, and decomposed \texttt{FormEBM-\allowbreak{}v3} reaches
  \texttt{94.0\%} on \texttt{E\_depth} at Spearman \texttt{0.9966} and
  \texttt{92.0\%} on \texttt{E\_total} at Spearman \texttt{0.9846}
  across 100 unseen cumulative defect chains vs \texttt{36.0\%} for
  \texttt{FormJEPA-\allowbreak{}v1}; unconstrained same-line random-walk slices
  remain below \texttt{80\%}.)
\item
  H12. The decision head is calibrated on gate outcomes.
\item
  H13. A discrete latent plans better than a continuous one over long
  edit sequences.
\item
  H14. A public deterministic ladder of exact proposers closes a
  measurable share of a registered task set with zero model calls, and
  that share is the baseline the learned proposer must beat on the
  residual. (Supported: 111 of 151 tasks, 73.5\%, leaving 40 open
  tasks.)
\item
  H15. Proposer classes are compared on the residual of 40 tasks under a
  matched budget, alongside Gear-0 crystallization and the 3-Gear
  Cascade. Benchmarked across six single-shot cloud/local proposers
  (\texttt{17/40} to \texttt{34/40}), purpose-trained
  \texttt{FormJEPA-\allowbreak{}v1} (\texttt{22/40} at \texttt{k\ =\ 1},
  \texttt{27/40} at \texttt{k\ =\ 16}; \texttt{25/40} typed
  \texttt{k\ =\ 1} with \texttt{FormEBM-\allowbreak{}v2} reranking; \texttt{24/40}
  typed \texttt{k\ =\ 1} and \texttt{27/40} multi-proposer
  \texttt{k\ =\ 1} under \texttt{FormEBM-\allowbreak{}v3} \texttt{E\_def}), bounded
  checker-feedback \texttt{H15-\allowbreak{}Iterative} with public standard-library
  headers (\texttt{40/40} on \texttt{Opus\ 5.5} and \texttt{Sonnet\ 5.\allowbreak{}5}
  with \texttt{0} contradictions, \texttt{39/40} on
  \texttt{Grok\ 4.\allowbreak{}20}), and Shape D (\texttt{DNULLMM1}) Gear-0 replay
  (\texttt{40/40} at \texttt{0} calls; \texttt{16.67×} 3-Gear episode
  reduction and \texttt{200/200} admissions on \texttt{Zipf(\allowbreak{}0.\allowbreak{}85)});
  diffusion and evolutionary arms remain open.
\end{enumerate}

\section{Falsifiers at a glance}\label{falsifiers-at-a-glance}

Six observations that would defeat part of the premise, stated in full
in Section 8 after the hypotheses they test.

\begin{enumerate}
\def\labelenumi{\arabic{enumi}.}
\tightlist
\item
  RJ-F1. Lift coverage plateaus below 80\% on any of the five target
  languages.
\item
  RJ-F2. Writing the intent specifications for the reference
  applications costs more than writing their Forms by hand.
\item
  RJ-F3. H3 fails while the token arm passes at the same gate.
\item
  RJ-F4. H11 fails and search does not reach accepted Forms within the
  beam budget.
\item
  RJ-F5. The decision head's calibration never beats a frequency table
  over shape hashes.
\item
  RJ-F6. A trace-trained token-space model closes the residual of 40
  tasks at equal energy and equal gate acceptance, with no use of the
  Form quotient.
\end{enumerate}

\section{What already exists, and what this page is
about}\label{what-already-exists-and-what-this-page-is-about}

The Reality Code Engine holds, in its public reference checker and in
the design of its engine, the half of the problem that decides. Its
public paper, \href{https://research.openie.dev/papers/rce/}{The Reality
Code Engine}, describes the design of that half: a commit gate that
checks every candidate at entry, exact replay (the runtime state is a
left fold of the transition function over the log), projection of one
Form to other languages, and integer cost accounting (\texttt{est\_j}),
kept strictly apart from operating-system/on-chip counter readings
(\texttt{reported\_\allowbreak{}j}, first recorded on Apple M5 Max at commit
\texttt{1762fb4b}) and external wall-meter readings
(\texttt{measured\_\allowbreak{}j}, of which none exists for RCE). Core theorems are
checked in Lean 4, a proof assistant (\texttt{5,537} kernel theorem
constants across \texttt{64} modules at \texttt{1762fb4b}), and the
specification and Lean library are public. This page quotes from that
paper only the public record of its model-proposer run and treats the
rest as design context.

\textbf{This page is about the proposer:} the learned part that suggests
which Form to try. The checker, replay and the projections to other
languages are public code, and the engine's gate and its \texttt{est\_j}
accounting are specified designs; Section 8 below now reports the
six-model off-the-shelf benchmark plus memory-mapped Gear-0
crystallization, the first purpose-trained latent Form JEPA repair
proposer (\texttt{FormJEPA-\allowbreak{}v1}, \texttt{8.42M} trainable parameters over
disjoint canonical Form S-expressions, admitting \texttt{177/200}
held-out synthetic Form defects and \texttt{27/40} open residual tasks
at sub-joule core \texttt{reported\_\allowbreak{}j}), the contrastive multi-step
margin energy scorer (\texttt{FormEBM-\allowbreak{}v2}, \texttt{5.85M} parameters on
the same disjoint \texttt{1,387}-Form split, reaching unseen-shape H5
AUROC \texttt{0.8922} and \texttt{93.0\%} strict 5-step edit
monotonicity on H11 while cutting solvable check calls by
\texttt{2.39×}), the span-normalized two-head decomposed scorer
(\texttt{FormEBM-\allowbreak{}v3}, \texttt{6.05M} parameters on the identical split,
lifting unseen-shape H5 AUROC to \texttt{0.9437}, H11 strict 5-step
monotonicity to \texttt{94.0\%} on \texttt{E\_depth} and \texttt{92.0\%}
on \texttt{E\_total}, and zero-shot out-of-distribution H15
multi-proposer pooled AUROC from \texttt{0.2585} to \texttt{0.5278},
\texttt{0.6018} on large-span valid versus short invalid rewrites), and
bounded checker-feedback residual closure with public standard-library
headers (\texttt{H15-\allowbreak{}Iterative}, closing \texttt{40/40} residual tasks
with \texttt{0} contradictions and \texttt{40/40} zero-call Gear-0
replay; interactive instruments at
\href{https://research.openie.dev/living/rce-jepa/}{/living/rce-jepa/}),
while full intent-to-Form synthesis, cross-module signature-and-trace
out-of-distribution energy calibration, and the gate-receipt decision
head remain open.

\subsection{Deterministic first: the ladder beneath the
proposer}\label{deterministic-first-the-ladder-beneath-the-proposer}

The first rule of automation is having a manual process that works. What
RCE-JEPA proposes makes sense only on top of a manual variant of the
task, or a deterministic automated variant, that already works; a
learned proposer is added to that, never put in its place. The design
therefore orders proposers as a ladder, and exact proposers run first:

\begin{enumerate}
\def\labelenumi{\arabic{enumi}.}
\tightlist
\item
  \textbf{Reuse.} A Form whose content hash (the hash of its canonical
  text), or whose name-free shape (the Form with every name replaced by
  its position, so that renaming does not change it), matches a stored
  Form is reused as it stands.
\item
  \textbf{Lookup by type.} \texttt{rce\ search} finds stored handlers
  whose type matches the one the task needs.
\item
  \textbf{Determined rewrites.} The rewrites \texttt{rce\ fmt} writes
  for the rules whose method is ``altered'' in the specification's rule
  table; each has one result, and a second run finds nothing to rewrite.
\item
  \textbf{Bounded enumeration.} Every small edit up to a fixed bound,
  enumerated in the canonical name-free projection, each candidate
  judged by the checker.
\item
  \textbf{Renames.} One use of a value name becomes another value name
  in scope there, or one binding becomes a name the checker's report
  names, up to a fixed bound, each candidate judged by the checker.
\end{enumerate}

The learned proposer acts only on the residual, the tasks the ladder
does not close. It proposes edits in the canonical Form edit space, and
the same checker judges every proposal exactly as it judges a rung's
candidate. Reducing what a stochastic model is required to do makes its
action space smaller, and a smaller action space is easier to master.
That is the reason to expect H10 (learning on the quotient saves data)
to hold, and it is why H14 makes the ladder's share the baseline the
learned proposer must beat on the residual. The ladder is public, as
\texttt{rce\ repair} of openIE-dev/rce, and has three runs registered
under \texttt{rce/\allowbreak{}bench/\allowbreak{}ladder} in openie-research-experiments, whose
files label rungs 1 to 5 with the letters a to e; rung 5 was added for
the third run. Their result, reported with H14 in Section 8 below, is
that it closes the mutation classes its rungs invert and 1 of 41 changes
from real history, so the residual is nearly every real defect.

\texttt{rce\ repair} on the main branch of openIE-dev/rce has a sixth
rung, labelled f, that runs only after rungs 1 to 5 have closed nothing.
Rung f calls no model itself. Programs written outside RCE, by a model
on a hosted service or a local machine, or by a person, reach it only as
an input file, each with a label naming its origin and recorded by its
SHA-256 digest. The checker judges each one as it judges a rung's
candidate, including the mutation gate, which refuses a program with an
example or check that no single small edit of the program can make fail;
a proposal that changes an example, a check or a law is refused before
it is judged. The rung's mode is data in the specification, one of
\texttt{off}, \texttt{observe} and \texttt{auto}, and it stays in
\texttt{observe}: every proposal is judged, each one the checker would
admit is reported with its diff, digest and origin, the program is left
open, and no source file is overwritten. In \texttt{auto}, the first
admitted proposal would close the program, and the report would label
that close stochastic.

The specification allows the change from \texttt{observe} to
\texttt{auto} only when a promotion criterion holds on a recorded
evaluation. The criterion was registered in openIE-dev/rce before any
observation (\texttt{rce-\allowbreak{}repair-\allowbreak{}rung-\allowbreak{}f-\allowbreak{}promotion-\allowbreak{}2026-\allowbreak{}11-\allowbreak{}03},
\texttt{sha256:\ abf7f131\allowbreak{}f240e3f9.\allowbreak{}.\allowbreak{}.\allowbreak{}}), with its SHA-256 pinned beside
it, and follows the rule that at least 98\% of closes come from
deterministic rungs and at most 2\% from stochastic ones, each
stochastic output gated. Over the 151 tasks of runs 2 and 3, each run
with the task's original program given as the known-good version, at
most 2\% of the closes may come from rung f; every proposal observed as
admitted must also pass the mutation gate; and none may contradict the
original, that is, behave differently from it after some sequence of
events. A deterministic command, \texttt{rce\ repair\ criterion},
evaluates the criterion from the observation records, and the
specification's own check refuses the \texttt{auto} mode unless a
recorded evaluation meets it. In the registered 151-task evaluation at
commit \texttt{1762fb4b} (reported in Section 8 below), rungs 1 to 5
given each task's known-good original close \texttt{116} of \texttt{151}
tasks at \texttt{0} model calls, leaving \texttt{35} open tasks
(\texttt{25} unparsable \texttt{W-SYN-1} files skipped by rung f, and
\texttt{10} open parsing tasks). Feeding combined external proposals
(\texttt{FormJEPA-\allowbreak{}v1} plus recorded residual repairs, \texttt{91}
proposals judged) yields \texttt{3} observed rung-f admissions:
\texttt{3/3} (\texttt{100.0\%}) pass the mutation gate and \texttt{0}
contradict the original (\texttt{3/3} restored), yet \texttt{3} of
\texttt{119} closes is \texttt{2.5\%\ \textgreater{}\ 2.0\%}, so
\texttt{rce\ repair\ criterion} deterministically returns
\texttt{holds:\ false} and keeps rung f in \texttt{observe} mode.

\subsection{The model steers; the engine
renders}\label{the-model-steers-the-engine-renders}

Robotics already works this way. In V-JEPA 2-AC, the action-conditioned
version of the V-JEPA 2 video model, a planner searches over robot
actions by minimizing the distance between a predicted future
representation and the representation of a goal image; it never
generates pixels {[}B1{]}. The encoder was pretrained on more than one
million hours of video, and the action-conditioned predictor on fewer
than 62 hours of unlabeled robot video {[}B1, abstract{]}. RCE-JEPA
proposes the same division for code: a learned model chooses edits to a
Form, and the engine renders every target language from the accepted
Form. No model emits code.

Code differs from robotics in one decisive way. A robot's world model
must \emph{learn} how the world moves, \(z_{t+1}=g_\phi(z_t,a_t)\). In
RCE the transition function \(\delta\) is given and proved, so the
forward dynamics need not be learned at all. What remains to learn is
the inverse: which edit moves a candidate toward what the intent
requires. MuZero learns a model because, as its authors put it, planning
methods ``rely on knowledge of the environment's dynamics'' {[}T1, Sec.
1{]}; here those dynamics are written down.

\subsection{One real Form}\label{one-real-form}

Before the architecture, the object. This is a complete program from the
public repository (\texttt{examples/\allowbreak{}surface/\allowbreak{}counter.\allowbreak{}rce} in
openIE-dev/rce, at commit \texttt{1762fb4b}), fourteen lines in RCE's
surface syntax, shown unedited:

\begin{Shaded}
\begin{Highlighting}[]
\NormalTok{intent "a counter that never goes negative"}
\NormalTok{    example Inc, Inc, Dec {-}\textgreater{} 1}
\NormalTok{    example Dec {-}\textgreater{} kept: Underflow}

\NormalTok{state Counter}
\NormalTok{    n: u64 = 0}

\NormalTok{event Inc}
\NormalTok{event Dec}

\NormalTok{error Underflow}

\NormalTok{step Counter on Inc {-}\textgreater{} n + 1}
\NormalTok{step Counter on Dec {-}\textgreater{} try n.checked\_sub(1) else Underflow}
\end{Highlighting}
\end{Shaded}

Read it as \(P=(S,E,\delta,D,\lambda)\). The \texttt{state} block is S,
one unsigned 64-bit integer starting at zero. The two \texttt{event}
lines are E. The two \texttt{step} lines are \(\delta\). The
\texttt{intent} block is a small specification of intent, written
\(\sigma(I)\) on this page: two required traces, one of which must end
in a kept step. A decrement at zero does not apply: the step keeps the
state with \texttt{Underflow}, and by theorem T8 the state is unchanged.
The proposer's task, stated on this example: given only the
\texttt{intent} block, find the \texttt{state}, \texttt{event} and
\texttt{step} lines that the gate accepts.

\subsection{Architecture}\label{architecture}

The box drawn here is a proposer interface. It is not one architecture.
A proposer reads the intent and the current Form and returns candidate
edits or candidate Forms. The gate is outside the interface. A JEPA
encoder, predicting in a latent space, is one instance. A trace-trained
decoder, a diffusion code model, reinforcement learning inside a formal
environment, evolutionary search, library learning and constrained
decoding are other instances. Section 6 below names each one, what it
needs from the checker, its known failure mode and which hypothesis it
would move. The decision head is the small typed instance of the same
interface. Any instance only proposes, and the Laws of Coding do not
depend on it.


\begin{center}
\openiefig{fig/rce-jepa-fig02-architecture.png}{0.94}{https://research.openie.dev/figures/rce-jepa/architecture.svg}
\end{center}


\emph{Figure 2. The model steers, the engine renders. Orange boxes are
learned and untrusted; the blue box is the exact rollout by the
transition function, used as a simulator during search; green boxes are
the parts the RCE paper describes, public code or specified design (see
that paper). Only the gate admits. A design diagram, not data. .}

{\footnotesize\setlength{\tabcolsep}{3.8pt}\renewcommand{\arraystretch}{1.15}
\begin{longtable}[]{@{}
  >{\raggedright\arraybackslash}p{(\linewidth - 6\tabcolsep) * \real{0.1111}}
  >{\raggedright\arraybackslash}p{(\linewidth - 6\tabcolsep) * \real{0.2469}}
  >{\raggedright\arraybackslash}p{(\linewidth - 6\tabcolsep) * \real{0.3210}}
  >{\raggedright\arraybackslash}p{(\linewidth - 6\tabcolsep) * \real{0.3210}}@{}}
\toprule\noalign{}
\begin{minipage}[b]{\linewidth}\raggedright
\textbf{Component}
\end{minipage} & \begin{minipage}[b]{\linewidth}\raggedright
\textbf{Role}
\end{minipage} & \begin{minipage}[b]{\linewidth}\raggedright
\textbf{Starting design}
\end{minipage} & \begin{minipage}[b]{\linewidth}\raggedright
\textbf{Why}
\end{minipage} \\
\midrule\noalign{}
\endhead
\bottomrule\noalign{}
\endlastfoot
Form encoder \(E_f\) & Map a Form and its trace on example events to a
latent \(z_f\) & A small transformer over the canonical JSON (JavaScript
Object Notation) of the Form plus trace tokens, regularized with SIGReg
(sketched isotropic Gaussian regularization) {[}B2{]}; a second arm with
a discrete or soft-codebook bottleneck {[}B3, B4, B5{]} (SC-JEPA's soft
codebook is stabilized together with an EMA (exponential moving average)
teacher and a reconstruction anchor {[}B5, Sec. 4.4.1 and App. B.1{]}) &
Traces define equivalence, and the trace is free because \(\delta\) is
exact. Discrete tokens kept long-horizon prediction stable where
continuous latents accumulated error {[}B3{]}; Boolean latents allow
exact state re-identification by bitwise equality, a motivation {[}B4{]}
states with reference to earlier work \\
Intent encoder \(E_I\) & Map \(\sigma(I)\) (required traces, contracts,
tests, prose) into the same space & Shared architecture, trained so that
\(E_I(\sigma(I))\) lies near \(E_f(f^\star)\) for known accepted fits
\(f^\star\) & Puts goal and candidate in one metric, as V-JEPA 2-AC does
with goal images {[}B1{]} \\
Proposer \(\Pi_\theta\) & Given \(z_I\) and the current \(z_f\), propose
a typed Form edit & A goal-conditioned inverse model, the shape LeWAM
compares against as GC-IDM (goal-conditioned inverse dynamics model), a
baseline from Nguyen et al.~(2026) that LeWAM evaluates, not a model
proposed by LeWAM's authors {[}B6, Sec. 4.2 and App. A.2{]} & With
\(\delta\) exact, the inverse is the learnable object \\
Energy \(\mathcal{E}\) & Rank candidates before the gate &
\(\alpha V(f)+\beta\,d(z_{\tau(f)},z_{\sigma(I)})+\gamma\,C(f)\), with
\(C(f)\) taken from \texttt{est\_j} & Separately trained energies are
not commensurate {[}T2, Sec. 1.3{]}; \(\alpha,\beta,\gamma\) are fitted
on a held-out slice to maximize AUROC (area under the receiver operating
characteristic curve) for accepted against requires change, and reported
with every run \\
Planner & Search over edits & Beam or tree search over proposer edits,
scored by \(\mathcal{E}\), with exact rollout at every node & V-JEPA
2-AC plans with the cross-entropy method (CEM) over continuous actions
{[}B1{]}; Form edits are discrete and combinatorial, so CEM does not
apply \\
Decision head & Bounded selection with a calibrated probability: the
value for each hole, a rank, an equivalence verdict & A small decision
model built by the RCE project and trained on gate receipts (next
subsection) & The calibrated probability \(p_i\) orders the
cheapest-first cascade \\
\end{longtable}
}

\textbf{Not learned, ever:} \(\delta\), \(\lambda\), the fold, the hash
chain, the projections, the gate and the promotion path. They are fixed
and certified, and promotion never retrains them.

\subsection{The decision head is a proposer, and it sits outside the
gate}\label{the-decision-head-is-a-proposer-and-it-sits-outside-the-gate}

A class of small, energy-aware decision models appeared in 2026. A
decision model takes a piece of state and a fixed question and returns a
typed answer with a probability for every option, without generating
text; all questions about one state are answered in a single request
{[}B7; B8{]}. Answers are bounded: a yes or no, a choice among a
caller-supplied set of options (up to 255 for Jev; Laya's accuracy drops
beyond about 20), or a score on a declared scale {[}B8{]}. Members of
the class, cited here as convergent examples of the concept and not as
components of RCE:

\begin{itemize}
\tightlist
\item
  \textbf{Laya}, an open-weight model of about 421 million parameters
  built on a ModernBERT-large encoder with a two-layer decision head,
  released under Apache 2.0 {[}B8, B9{]}.
\item
  \textbf{Kev}, a family of open decision models from 0.8 to 27 billion
  parameters built on Qwen backbones, each checkpoint shipped with a
  fitted temperature for calibration {[}B10{]}.
\item
  \textbf{Liquid AI's open d1 models}, d1-3B and d1-omni-600M, released
  on 7 October 2026 for edge and physical systems: d1-3B answers one
  question in 16 ms on an NVIDIA Jetson AGX Thor and 50 ms on a Jetson
  Orin Nano, reads camera frames, and was shown navigating an
  environment in the Isaac Sim robot simulator with the model served on
  a Jetson in a hardware-in-the-loop setup {[}B11{]}.
\end{itemize}

An independent benchmark of one hosted member of the class found its
choice probabilities well calibrated, with an expected calibration error
(ECE) of 0.028 pooled over 22 choice datasets and 279,925 answers
{[}B7{]}; its binary (yes or no) probabilities were less well placed
(ECE 0.052 single, 0.168 multi-label) and needed tuned thresholds {[}B7,
Sec. 4.3{]}. Documented weaknesses of the class include counting,
arithmetic and date comparison, and lower accuracy when the state is
padded with irrelevant content {[}B8, B12{]}.

The decision head is one plug in the proposer interface above. A JEPA
encoder is a different plug. Hosted or local, the head only proposes.

\textbf{What the RCE project plans.} The project will build its own
decision head, trained and calibrated on its own gate receipts. In the
engine's design every gate decision writes a receipt with the candidate,
the outcome, the broken rule and \texttt{est\_j}, so calibration data
would cost no labelling. The state handed to the head is the Form's
canonical JSON and its trace, which is small by construction; that
addresses the documented padded-state weakness directly. The head is a
proposer. Any decision model, hosted or local, built by the project or
by anyone else, only proposes; it never admits a Form, and the security
of RCE never depends on it.

\subsection{Data}\label{data}

V-JEPA 2 pairs very large passive pretraining with a small
action-conditioned set {[}B1{]}. The RCE analogue has four tiers.

{\footnotesize\setlength{\tabcolsep}{3.8pt}\renewcommand{\arraystretch}{1.15}
\begin{longtable}[]{@{}
  >{\raggedright\arraybackslash}p{(\linewidth - 6\tabcolsep) * \real{0.1500}}
  >{\raggedright\arraybackslash}p{(\linewidth - 6\tabcolsep) * \real{0.3714}}
  >{\raggedright\arraybackslash}p{(\linewidth - 6\tabcolsep) * \real{0.1786}}
  >{\raggedright\arraybackslash}p{(\linewidth - 6\tabcolsep) * \real{0.3000}}@{}}
\toprule\noalign{}
\begin{minipage}[b]{\linewidth}\raggedright
\textbf{Tier}
\end{minipage} & \begin{minipage}[b]{\linewidth}\raggedright
\textbf{Source}
\end{minipage} & \begin{minipage}[b]{\linewidth}\raggedright
\textbf{Size target}
\end{minipage} & \begin{minipage}[b]{\linewidth}\raggedright
\textbf{Trains}
\end{minipage} \\
\midrule\noalign{}
\endhead
\bottomrule\noalign{}
\endlastfoot
Passive: Forms and traces & Forms lifted from pinned public corpora by
the lift of the RCE paper (Sec. 6) once it is public code, Forms from
the seeded generator of that paper's ecosystem protocol (Sec. 2.3), and
every promoted handler & \(10^5\) to \(10^6\) Forms with traces & The
Form encoder, by predicting masked rules and masked trace segments in
latent space \\
Paired: intent and fit & Tests and contracts attached to lifted
programs; a Form's own trace as a self-generated specification; the
preregistered brownfield applications; benchmark specifications written
as Forms & \(10^4\) pairs to start & Intent encoder and proposer \\
Decision receipts & Every gate decision & Unbounded, no labelling cost &
The decision head's calibration \\
Hard negatives & Requires-change verdicts with the named broken rule;
model-proposed equivalences the gate did not admit (15 of 18 proposals,
12 distinct rewrites, in the public record of the model-proposer run in
the RCE paper, Sec. 5.5) & Grows with every run & Contrastive signal:
valid against invalid, equivalent against not \\
\end{longtable}
}

Near-duplicate Forms are removed by shape hash before any split, and
held-out splits are by shape hash, never by surface text. The first
corpus exists: version 1 of the Form corpus (23 October 2026) holds 192
distinct shapes of hand-written programs from the public repository
(version 4, 24 October 2026), 55 of them with an intent block, and 3,456
negatives (343 conformance programs that require change, each with its
diagnostic code, and 3,113 checker-verified mutants, whitespace-token
and typed, that require change across 31 diagnostic codes), with a
dataset card and a split by shape hash; its generator is
\texttt{scripts/\allowbreak{}corpus.\allowbreak{}py} in the repository and the rows are at
\href{https://research.openie.dev/data/\#rce-jepa}{research.openie.dev/data/}.
It is the paired and hard-negative tiers at their smallest; the passive
tier (lifted programs) is empty until the lift pipeline is public. On
it, H1 has a first measurement, preregistered before the run: the median
of rho = b(code)/b(P) over 185 programs is 9.34 against the emitted
Python (interquartile range 8.45 to 10.81; version 1 gave 9.31 over 173,
with 17.33 for Rust, 11.06 for JavaScript and 13.66 for Scala), above
the threshold of 5; 171 of the first 173 programs were above 5 (Wilson
95\% interval 0.96 to 1.00). Under a compressed description length
(gzip) the ratio is 3.57 for Python and 5.73 for Rust, which is reported
next to the primary so that the number is not read as a claim about
information content: H1 is about the space a proposer searches, which is
the token count under its vocabulary. The emitted source is the
reference checker's own, with the shared runtime excluded, and no lifted
or model-written program is in the population. On Apple M5 Max at commit
\texttt{1762fb4b}, H2 is now measured in operating-system/on-chip
\texttt{reported\_\allowbreak{}j}: \texttt{rce\ check} takes a median of
\texttt{7.32\ ms} and \texttt{0.\allowbreak{}04904\ J} per Form check (kernel
\texttt{ri\_\allowbreak{}energy\_\allowbreak{}nj}) against \texttt{84.25\ J} per candidate on
local \texttt{Qwen\ 3.\allowbreak{}6\ 35B-\allowbreak{}A3B}
(\texttt{C\_check\ /\ E\_model\ =\ 0.000582\ \textless{}=\ 0.05};
details in Section 8 below). And a measure of what an intent block buys,
which RJ-F2 and H7 both rest on: of the mutants that change a program's
behaviour, its own example lines catch 317 of 392 (0.81, Wilson 95\%
interval 0.77 to 0.85; 24 of 54 programs catch every one), and the
uncaught changes are almost all a constant moved by one that no example
exercises.

\subsection{Training, in brief}\label{training-in-brief}

\begin{itemize}
\tightlist
\item
  \textbf{Encoder objective.} Latent prediction plus SIGReg, the recipe
  of LeJEPA, which needs a single trade-off hyperparameter and no
  stop-gradient or teacher-student scheme {[}B2, abstract{]}, and of
  LeWorldModel, which trains end to end with two loss terms {[}B13,
  abstract{]}. Two cautions from the same line of work: the LeWorldModel
  authors flag the isotropic Gaussian target as a possible limitation in
  environments of low intrinsic dimension {[}B13, Sec. 4.2 and Sec.
  6{]}, and contrastive inverse dynamics is offered as a
  distribution-free alternative {[}B14{]}, which, however, was evaluated
  on continuous controls only and was clearly worse than SIGReg on
  combinatorial puzzle tasks whose discrete state changes look alike
  whichever action caused them {[}B14, Appendix D{]}; and SIGReg and
  VISReg ``do not guarantee that the learned latent representation
  preserves all the key information required for planning'' {[}B15, Sec.
  4.1{]}. Both arms run side by side, with auxiliary heads that predict
  exact, cheap properties of the Form (monotone tag, fuel bound, rule
  count).
\item
  \textbf{Multi-step prediction.} Training with k-step recursive
  prediction gives significantly better results in the EB-JEPA
  ablations, by reducing exposure bias, with a Pareto optimum near k = 4
  {[}B16, Fig. 3{]}; RoboJEPA raises the autoregressive training rollout
  from K = 2 to K = 10 steps in a long-horizon cooldown, which pays off
  only at large compute (above about 10\textsuperscript{21} FLOPs)
  {[}B17, Sec. 3.3 and App. C.2{]}; a rollout-consistency term matched
  to the planner horizon comes from Temporal-Distance JEPA {[}B18,
  abstract{]}.
\item
  \textbf{Monotone energy.} A monotone cost ranking makes plans that
  deviate further from a goal-reaching trajectory cost more; without it,
  local minima in the cost surface can lead the planner to suboptimal
  plans {[}B19{]}. Monotonicity is not guaranteed by the task: the
  EB-JEPA navigation task is non-monotonic, since the best path first
  moves away from the goal {[}B16{]}. Here it is a preregistered
  diagnostic (H11), not an assumption.
\item
  \textbf{Proposer.} Supervised on intent, Form and edit triples from
  the paired tier, then improved on its own gate-accepted trajectories:
  the gate is the reward.
\item
  \textbf{Never decode a latent directly.} MoRAE reports that generating
  motion directly in a frozen JEPA latent ``fails dramatically'', and
  traces the failure to an ill-conditioned latent spectrum and to small
  latent errors amplified into large artifacts after decoding {[}B20,
  abstract and Sec. 3{]}. The proposer outputs typed edits; latents
  rank, they do not render.
\end{itemize}

\subsection{Convergence: where each piece already
exists}\label{convergence-where-each-piece-already-exists}

No single piece of RCE-JEPA is new. Each has an existing home, listed
below as convergence, not priority. What has not been measured is the
joined system, which is the subject of the hypotheses.

{\footnotesize\setlength{\tabcolsep}{3.8pt}\renewcommand{\arraystretch}{1.15}
\begin{longtable}[]{@{}
  >{\raggedright\arraybackslash}p{(\linewidth - 6\tabcolsep) * \real{0.1572}}
  >{\raggedright\arraybackslash}p{(\linewidth - 6\tabcolsep) * \real{0.2516}}
  >{\raggedright\arraybackslash}p{(\linewidth - 6\tabcolsep) * \real{0.3270}}
  >{\raggedright\arraybackslash}p{(\linewidth - 6\tabcolsep) * \real{0.2642}}@{}}
\toprule\noalign{}
\begin{minipage}[b]{\linewidth}\raggedright
\textbf{Piece}
\end{minipage} & \begin{minipage}[b]{\linewidth}\raggedright
\textbf{Where it exists}
\end{minipage} & \begin{minipage}[b]{\linewidth}\raggedright
\textbf{What it established}
\end{minipage} & \begin{minipage}[b]{\linewidth}\raggedright
\textbf{Difference from RCE-JEPA}
\end{minipage} \\
\midrule\noalign{}
\endhead
\bottomrule\noalign{}
\endlastfoot
Predict a compact plan, then decode the program & Latent Programmer
{[}B21{]} & A discrete latent code predicted from the specification,
then the program; string-transformation accuracy up 5 to 7 points (over
10\% relative) over the best baseline, plus higher BLEU on Python code
generation from docstrings {[}B21, Sec. 5.1 and 5.3{]} & Open target
language, no canonical form, no checker \\
Learn execution, not only text & LaSynth {[}B22{]}; Code World Model
(CWM) {[}B23{]}; analysis of CWMs {[}B24{]} & Latent execution improves
synthesis in a restricted C domain; CWM is mid-trained on
observation-action trajectories from a Python interpreter and Docker
environments & Dynamics are learned; in RCE they are exact \\
Equivalence classes as a learning target & EqNet, Allamanis et
al.~{[}B25{]} & Continuous representations of symbolic expressions
trained on equivalence classes & Expressions, not whole programs \\
Agents propose, solvers verify & Agentic model checking {[}B26{]};
verified synthesis for concept design driven by a large language model
(LLM) {[}B27{]}; AlphaProof, Clover, AutoVerus (Section 3 below) &
``Agents propose, solvers verify'' is stated as a principle; CEGIS-style
synthesis (counterexample-guided inductive synthesis) against Alloy
invariants and positive and negative scenarios & Proposers generate
tokens; no latent planning, no joule accounting \\
JEPA as a latent reasoner & JEPA-Reasoner {[}B28{]} & Latent reasoning
decoupled from token generation & Text domain \\
Cost predicted before execution & Compiler world models {[}B29{]} & A
transition model rolls out scheduling actions in latent space to predict
the terminal-state representation, which a ranker uses to order
candidates before fewer hardware measurements; the score is a
within-workload rank, not a calibrated cost {[}B29, Sec. 3.2 and Sec.
6{]} & Tensor program schedules only \\
Discrete latents & Discrete-JEPA {[}B3{]}; discrete world models via
regularization {[}B4{]}; SC-JEPA {[}B5{]} & Discrete tokens that keep
long-horizon symbolic prediction stable {[}B3{]}; Boolean latents
learned without reconstruction (planning left to future work) {[}B4{]};
a soft codebook bottleneck, with an EMA teacher and a reconstruction
anchor, to stabilize one-step latent prediction, not long-horizon
rollout {[}B5{]} & Synthetic images, puzzles, time series; not
programs \\
Stable JEPA training & LeJEPA {[}B2{]}; LeWorldModel {[}B13{]} &
Collapse-free training without stop-gradient, teacher networks or
hyper-parameter schedulers; a standard learning-rate warm-up and cosine
schedule is still used {[}B2, abstract and Sec. 6.1{]}; without
stop-gradient, EMA or pretrained encoders, with two loss terms {[}B13,
Sec. 3.1{]} & Reusable as an objective; validated only on images
{[}B2{]} and pixels {[}B13{]}, untested on program data \\
Small calibrated decision models & Laya {[}B9{]}; Kev {[}B10{]}; Liquid
AI open d1 {[}B11{]}; benchmark {[}B7{]} & Typed answers with per-option
probabilities. Kev ships a fitted temperature per checkpoint {[}B10{]};
Laya ships over-confident (mean ECE 0.466) and reaches 0.081 only after
a per-type temperature refit {[}B9{]}. Latency: about 33 to 40 ms per
question on a T4 GPU for Laya and 18.1 ms for six questions on an H100
for Kev, but hundreds of milliseconds on a CPU or Mac (Laya 193 to 464
ms on CPU, Kev-4B 721 ms on an M5 Mac) {[}B9, B10{]}; d1-3B answers one
question in 16 ms on a Jetson AGX Thor and 50 ms on a Jetson Orin Nano
{[}B11{]}; one hosted member independently benchmarked for calibration
{[}B7{]} & No world model; in RCE-JEPA a decision head is one proposer
among several \\
Trace-trained proposers & Code World Model {[}B23{]}, Section 3 below &
Mid-training on Python execution traces and on agent trajectories.
SWE-bench Verified \href{mailto:pass@1}{\nolinkurl{pass@1}} 53.9\% base
and 65.8\% with a best of 16 samples, as that report states & The
dynamics are learned. In a Form the transition is exact. The authors
call trace prediction a ``first step'' \\
Evaluator-gated evolutionary search & AlphaEvolve (Novikov et al.,
Section 3 below) & A model proposes edits. An automatic evaluator
decides. A 4 by 4 complex matrix multiplication in 48 scalar
multiplications is the kernel figure used here & Works where an
evaluator exists. That report does not use a Form quotient \\
Constraint at sampling time & Type-constrained decoding (Mündler et al.,
the Swiss Federal Institute of Technology Zurich and the University of
California, Berkeley, Section 3 below) & Compilation errors down 75.3\%
on HumanEval and 52.1\% on MBPP. Syntax only: 9.0\% and 4.9\% &
TypeScript scale. Types are not the full acceptance rule \\
Specification to verified code & Vericoding (Bursuc et al., Section 3
below) & 12,504 specifications. Success 82.2\% Dafny, 44.2\% Verus,
26.8\% Lean. The abstract rounds those shares to 82\%, 44\% and 27\% &
The specification still has to be written. Natural-language descriptions
add little \\
\end{longtable}
}

\subsection{The specification cost,
named}\label{the-specification-cost-named}

The checker guarantees that a Form is valid; it does not guarantee that
it is the Form that was meant. In verified LLM-driven synthesis for
concept design, invariant-only synthesis ``reached verified designs
quickly but often produced inconsistent designs across runs'', which
shows ``that invariants alone underconstrain the design task'';
scenario-guided synthesis recovered intended designs more consistently
than natural-language prompting, although minimal scenarios could
overfit {[}B27, abstract{]}, on 3 applications and 12 variants with one
LLM configuration {[}B27, Sec. 5{]}. The dominant failures ``were not
failures of bounded verification, but failures of specification
coverage'' {[}B27, Sec. 5.6{]}. In seL4 the abstract specification took
about 4 person-months and the proof about 20 person-years (Klein et
al.~2009, Section 2 below). A closed program world makes acceptance
exact and search tractable. It moves the uncertainty of intent into
\(\sigma(I)\); it does not remove it. Falsifier RJ-F2 below makes that
cost public and measurable.

\subsection{Proposer hypotheses}\label{proposer-hypotheses}

The proposer is tested by H11 (monotone energy), H12 (decision-head
calibration) and H13 (discrete against continuous latents), and is
scored on the residual of the deterministic ladder above (H14, H15). All
fifteen hypotheses, H1 to H15, are stated in numerical order with their
thresholds in one table in Section 8 below, under the same conventions:
preregistered before the run, frozen engine commit and binary hash
recorded, same gate for every arm.

The H3 control is fixed now: the same gate and the same cost accounting,
with Qwen2.5-Coder-1.5B proposing surface code, the arm already run in
the RCE paper (Sec. 5.5), where the gate did not admit 15 of 18 proposed
equivalences and admitted only one distinct rewrite, a rule the engine
already had.

\textbf{Falsifiers.} The falsifiers for the premise itself, RJ-F1 to
RJ-F6, now follow the hypotheses they test, directly after the H1 to H15
table in
\href{https://research.openie.dev/papers/rce-jepa/\#falsifiers}{Section
8 below}.

\subsection{Risks and mitigations}\label{risks-and-mitigations}

{\footnotesize\setlength{\tabcolsep}{3.8pt}\renewcommand{\arraystretch}{1.15}
\begin{longtable}[]{@{}
  >{\raggedright\arraybackslash}p{(\linewidth - 4\tabcolsep) * \real{0.2353}}
  >{\raggedright\arraybackslash}p{(\linewidth - 4\tabcolsep) * \real{0.3824}}
  >{\raggedright\arraybackslash}p{(\linewidth - 4\tabcolsep) * \real{0.3824}}@{}}
\toprule\noalign{}
\begin{minipage}[b]{\linewidth}\raggedright
\textbf{Risk}
\end{minipage} & \begin{minipage}[b]{\linewidth}\raggedright
\textbf{Evidence}
\end{minipage} & \begin{minipage}[b]{\linewidth}\raggedright
\textbf{Mitigation}
\end{minipage} \\
\midrule\noalign{}
\endhead
\bottomrule\noalign{}
\endlastfoot
Decoding a latent amplifies small errors & Direct generation in a frozen
JEPA latent fails dramatically for motion, from spectral
ill-conditioning and decoder amplification {[}B20{]} & Never decode a
latent; the proposer emits typed edits and the gate checks them \\
A representation that does not collapse but learns the wrong thing &
``Physical representation laziness'' under SIGReg {[}B15{]} & Auxiliary
exact-property heads; an inverse-dynamics regularizer \\
Isotropic prior misfits a low-dimensional Form space & {[}B13{]};
{[}B14{]} & SIGReg and contrastive inverse-dynamics arms side by side,
remembering that the inverse-dynamics arm was weaker on discrete puzzle
state {[}B14{]} \\
Energy terms not commensurate & {[}T2, Sec. 1.3{]} & Fit
\(\alpha,\beta,\gamma\) per experiment and publish them \\
Decision head degrades on padded state & Documented weakness of the
class {[}B8, B12{]} & Feed only canonical Form JSON and trace \\
Expressiveness ceiling & Rice's theorem {[}T3{]}; the restricted
language & Lift coverage is the measurement; RJ-F1 is the falsifier \\
Specification cost dominates & seL4 effort (Section 2 below); invariant
under-constraint {[}B27{]} & RJ-F2; tools that capture traces from
running programs and infer contracts, as deliverables in their own
right \\
Benchmark oracle weakness & PatchDiff (Wang, Pradel and Liu, Section 3
below): 7.8\% of patches counted correct fail developer tests, and
resolve rates are inflated by about 6.4 points. A METR trial of
early-2025 tools (Becker et al., Section 3 below): 16 developers, 246
tasks, measured completion time 19\% higher. The forecast was 24\%
lower. A post-hoc estimate was 20\% lower & Acceptance is the gate plus
replay, not a test suite. Section 8 (H15) enforces \texttt{rce\ check}
plus obligation preservation on every candidate \\
Learned planning latency & V-JEPA 2 model predictive control took about
16 seconds per action {[}B1{]} & Search over Forms with an exact checker
(\texttt{7.32\ ms} / \texttt{49.04\ mJ} on Apple M5 Max) and Gear-0
crystallization (\texttt{0.325\ µs} / \texttt{1.77\ µJ} per lookup)
bound repeated rollouts \\
\end{longtable}
}

The phased pilot plan, from the data engine through convergence, with
preregistration and exit criteria for each phase, is kept with the
project's working notes and summarized here: Phase 0 builds the lift
corpus and receipt dataset, and measures the deterministic ladder for
H14, which needs no learned proposer; Phase 1 trains the Form encoder
and tests H1, H5 and H11; Phase 2 builds and calibrates the decision
head (H12); Phase 3 trains the dedicated JEPA proposer and planner
against the token control (H3, H4, H10, H13), extending the first
single-shot off-the-shelf H15 baseline reported in Section 8; Phase 4
extends on-chip \texttt{reported\_\allowbreak{}j} (H2) to external wall-meter
\texttt{measured\_\allowbreak{}j} (H8, H9); Phase 5 runs the convergence protocol, in
which Gear-0 promotion takes the model out of the loop for each closed
shape. Alongside Phase 0, a slice of the vericoding benchmark (source
C34, listed under Convergence by commitment below), chosen as its
state-machine-shaped tasks in Dafny, Verus and Lean, is lifted into
Forms and scored under RCE's gate next to the published vericoding
results, an outside benchmark for the step from specification to Form
(H6, H7 and H8).

\subsection{Proposer references}\label{proposer-references}

References marked {[}T{]} are listed under Thesis references below;
references to Sections 2, 3 and 5 are to the sections of this page
below. Labels follow the convention of the main references below.
\textbf{Full} here means the whole text was read, including references
and appendices; for B9 and B10 the live model card and README were read
on 14 October 2026.

\begin{enumerate}
\def\labelenumi{\arabic{enumi}.}
\tightlist
\item
  {[}B1{]} M. Assran et al., ``V-JEPA 2: Self-Supervised Video Models
  Enable Understanding, Prediction and Planning'', 2025.
  \href{https://arxiv.org/abs/2506.09985}{arXiv:2506.09985}. Full.
\item
  {[}B2{]} R. Balestriero and Y. LeCun, ``LeJEPA: Provable and Scalable
  Self-Supervised Learning Without the Heuristics'', 2025.
  \href{https://arxiv.org/abs/2511.08544}{arXiv:2511.08544}. Full.
\item
  {[}B3{]} J. Baek et al., ``Discrete JEPA: Learning Discrete Token
  Representations without Reconstruction'', 2025.
  \href{https://arxiv.org/abs/2506.14373}{arXiv:2506.14373}. Full.
\item
  {[}B4{]} D. Bizzaro and L. Serafini, ``Discrete World Models via
  Regularization'', 2026.
  \href{https://arxiv.org/abs/2603.01748}{arXiv:2603.01748}. Full.
\item
  {[}B5{]} Y. He et al., ``SC-JEPA: Stabilizing Latent Predictive
  Learning for Time-Series Anomaly Prediction'', 2026.
  \href{https://arxiv.org/abs/2602.04643}{arXiv:2602.04643}. Full.
\item
  {[}B6{]} M. Fu et al., ``LeWAM: End-to-End World and Action Modeling
  with JEPAs'', 2026.
  \href{https://arxiv.org/abs/2610.12818}{arXiv:2610.12818}. Full.
\item
  {[}B7{]} T. Deußer, L. Sparrenberg and R. Sifa, ``Evaluating and
  Benchmarking the System One Model Jev'', 2026.
  \href{https://arxiv.org/abs/2609.37647}{arXiv:2609.37647}. Full.
\item
  {[}B8{]} W. Wu, ``Jev vs Laya: Comparing and Choosing Between
  Closed-Source and Open-Source System One Decision Models'', 23
  September 2026.
  \href{https://wilsonwu.me/en/blog/2026/jev-vs-laya/}{wilsonwu.me}.
  Full.
\item
  {[}B9{]} Nandakishor M (Convai Innovations), Laya model card.
  \href{https://huggingface.co/convaiinnovations/laya}{Hugging Face}.
  Full (read live 14 October 2026).
\item
  {[}B10{]} J. Palmer, Kev.
  \href{https://github.com/jaredpalmer/kev}{GitHub}. Full (README, read
  live 14 October 2026).
\item
  {[}B11{]} Liquid AI, ``Open d1: Edge decision models for text, vision,
  and audio'', 7 October 2026.
  \href{https://www.liquid.ai/blog/d1-open}{liquid.ai}. Full.
\item
  {[}B12{]} Sanity, ``Decision model'' (glossary).
  \href{https://www.sanity.io/glossary/decision-model}{sanity.io}. Full.
\item
  {[}B13{]} L. Maes, Q. Le Lidec, D. Scieur, Y. LeCun and R.
  Balestriero, ``LeWorldModel: Stable End-to-End Joint-Embedding
  Predictive Architecture from Pixels'', 2026.
  \href{https://arxiv.org/abs/2603.19312}{arXiv:2603.19312}. Full.
\item
  {[}B14{]} J. Boylan and C. Hokamp, ``No Gaussian Required: Contrastive
  Inverse Dynamics for JEPA World Models'', 2026.
  \href{https://arxiv.org/abs/2608.17542}{arXiv:2608.17542}. Full.
\item
  {[}B15{]} P. Zhu et al., ``Spectral-Target Physical Latent Structuring
  for JEPA-Style World Models'', 2026.
  \href{https://arxiv.org/abs/2609.04264}{arXiv:2609.04264}. Full.
\item
  {[}B16{]} B. Terver et al., ``A Lightweight Library for Energy-Based
  Joint-Embedding Predictive Architectures'', 2026.
  \href{https://arxiv.org/abs/2602.03604}{arXiv:2602.03604}. Full.
\item
  {[}B17{]} A. Zholus et al., ``RoboJEPA: Scaling Robotic Latent World
  Models'', 2026.
  \href{https://arxiv.org/abs/2610.10515}{arXiv:2610.10515}. Full.
\item
  {[}B18{]} J. Bai and J. Xiong, ``Temporal-Distance JEPA: Plan-Aware
  Representation Learning for Latent World Model Predictive Control'',
  2026. \href{https://arxiv.org/abs/2607.25337}{arXiv:2607.25337}. Full.
\item
  {[}B19{]} A. Chahe, S. Cai and L. Zhou, ``Latent World Models with
  Monotone Planning Costs for Image-Goal Navigation'', 2026.
  \href{https://arxiv.org/abs/2608.09073}{arXiv:2608.09073}. Full.
\item
  {[}B20{]} Y. Zhu et al., ``MoRAE: Flow-Friendly Self-Supervised
  Latents for Text-to-Motion Generation'', 2026.
  \href{https://arxiv.org/abs/2607.29180}{arXiv:2607.29180}. Full.
\item
  {[}B21{]} J. Hong, D. Dohan, R. Singh, C. Sutton and M. Zaheer,
  ``Latent Programmer: Discrete Latent Codes for Program Synthesis'',
  ICML 2021.
  \href{https://proceedings.mlr.press/v139/hong21a/hong21a.pdf}{PMLR
  PDF}. Full.
\item
  {[}B22{]} X. Chen, D. Song and Y. Tian, ``Latent Execution for Neural
  Program Synthesis'', NeurIPS 2021.
  \href{https://arxiv.org/abs/2107.00101}{arXiv:2107.00101}. Full.
\item
  {[}B23{]} Meta FAIR CodeGen Team, ``CWM: An Open-Weights LLM for
  Research on Code Generation with World Models'', 2025.
  \href{https://arxiv.org/abs/2510.02387}{arXiv:2510.02387}. Full.
\item
  {[}B24{]} B. Rahmani, ``Debugging Code World Models'', 2026.
  \href{https://arxiv.org/abs/2602.07672}{arXiv:2602.07672}. Full.
\item
  {[}B25{]} M. Allamanis, P. Chanthirasegaran, P. Kohli and C. Sutton,
  ``Learning Continuous Semantic Representations of Symbolic
  Expressions'', ICML 2017.
  \href{https://arxiv.org/abs/1611.01423}{arXiv:1611.01423}. Full.
\item
  {[}B26{]} Y. Sun, J. Liu, D. Kroening and J. Xue, ``Agentic Model
  Checking'', 2026.
  \href{https://arxiv.org/abs/2605.21434}{arXiv:2605.21434}. Full.
\item
  {[}B27{]} A. Cunha, ``Verified LLM-Driven Synthesis for Concept
  Design'', 2026.
  \href{https://arxiv.org/abs/2607.15718}{arXiv:2607.15718}. Full.
\item
  {[}B28{]} B. K. Liu, Z. P. Chen and D. P. Woodruff, ``JEPA-Reasoner:
  Decoupling Latent Reasoning from Token Generation'', 2025.
  \href{https://arxiv.org/abs/2512.19171}{arXiv:2512.19171}. Full.
\item
  {[}B29{]} H. Pan et al., ``Toward Compiler World Models: Learning
  Latent Dynamics for Efficient Tensor Program Search'', 2026.
  \href{https://arxiv.org/abs/2606.09312}{arXiv:2606.09312}. Full.
\item
  {[}B30{]} Open Interface Engineering, ``Exact-proposer ladder:
  registered run'', openIE-dev/openie-research-experiments,
  \texttt{rce/\allowbreak{}bench/\allowbreak{}ladder/\allowbreak{}exact-\allowbreak{}proposers/\allowbreak{}}: registration committed
  before the run at \texttt{3030469d}, results at \texttt{70a0a776}
  (pull request \#21); ladder openIE-dev/rce \texttt{6d04a096}.
  \href{https://github.com/openIE-dev/openie-research-experiments/tree/main/rce/bench/ladder/exact-proposers}{GitHub}.
  Full.
\item
  {[}B31{]} Open Interface Engineering, ``Exact-proposer ladder on real
  defects: registered run'', openIE-dev/openie-research-experiments,
  \texttt{rce/\allowbreak{}bench/\allowbreak{}ladder/\allowbreak{}real-\allowbreak{}defects/\allowbreak{}}: registration committed before
  the run at \texttt{64e70389}, results at \texttt{00e3b005} (pull
  request \#22); ladder openIE-dev/rce \texttt{7dc4eded}.
  \href{https://github.com/openIE-dev/openie-research-experiments/tree/main/rce/bench/ladder/real-defects}{GitHub}.
  Full.
\item
  {[}B32{]} Open Interface Engineering, ``Exact-proposer ladder with
  renames: registered run'', openIE-dev/openie-research-experiments,
  \texttt{rce/\allowbreak{}bench/\allowbreak{}ladder/\allowbreak{}rung-\allowbreak{}e/\allowbreak{}}: registration committed before the
  run at \texttt{be5f370d}, results at \texttt{3a5227c2}, merged at
  \texttt{f76ba9d8} (pull request \#23); ladder openIE-dev/rce
  \texttt{8a5d90ea}.
  \href{https://github.com/openIE-dev/openie-research-experiments/tree/f76ba9d8d02d57b773afbf105333417d55f72e2c/rce/bench/ladder/rung-e}{GitHub}.
  Full.
\item
  {[}B33{]} Open Interface Engineering, ``RCE-JEPA H2, single-shot H15,
  and Shape D (DNULLMM1) Gear-0 on the 40 open residual tasks'',
  registered 22 November 2026
  (\texttt{rce-\allowbreak{}jepa-\allowbreak{}h15-\allowbreak{}and-\allowbreak{}shape-\allowbreak{}d-\allowbreak{}gear0-\allowbreak{}2026-\allowbreak{}11-\allowbreak{}22},
  \texttt{sha256:\ b64c4c67\allowbreak{}1ae85755.\allowbreak{}.\allowbreak{}.\allowbreak{}}) on \texttt{openIE-\allowbreak{}dev/\allowbreak{}rce} at
  commit \texttt{1762fb4b} (20 November 2026, binary
  \texttt{sha256:\ 1d3b0d07\allowbreak{}ff5b2e8f.\allowbreak{}.\allowbreak{}.\allowbreak{}}), Apple M5 Max dual-channel
  \texttt{reported\_\allowbreak{}j} receipts (\texttt{measured\_\allowbreak{}j\ =\ None},
  \texttt{board\_\allowbreak{}metered\ =\ false}).
  \href{https://research.openie.dev/data/\#rce-jepa}{research.openie.dev/data/}.
  Full.
\item
  {[}B34{]} Open Interface Engineering, ``FormJEPA-v1 latent Form repair
  proposer, 3-Gear Cascade, and rce repair rung f criterion evaluation
  on the 151-task ladder and 40 residual tasks'', registered 22 November
  2026 (\texttt{rce-\allowbreak{}jepa-\allowbreak{}form-\allowbreak{}jepa-\allowbreak{}v1-\allowbreak{}2026-\allowbreak{}11-\allowbreak{}22},
  \texttt{sha256:\ 22682e6e\allowbreak{}936ec120.\allowbreak{}.\allowbreak{}.\allowbreak{}};
  \texttt{rce-\allowbreak{}repair-\allowbreak{}rung-\allowbreak{}f-\allowbreak{}wave3-\allowbreak{}eval-\allowbreak{}2026-\allowbreak{}11-\allowbreak{}22},
  \texttt{sha256:\ 04aff4b2\allowbreak{}b44cdbea.\allowbreak{}.\allowbreak{}.\allowbreak{}}) on \texttt{openIE-\allowbreak{}dev/\allowbreak{}rce} at
  commit \texttt{1762fb4b} (Linux x86\_64 binary
  \texttt{sha256:\ 5ef9931f\allowbreak{}f23a310f.\allowbreak{}.\allowbreak{}.\allowbreak{}}), AMD Ryzen AI Max+ 395 Linux
  powercap RAPL \texttt{reported\_\allowbreak{}j} receipts
  (\texttt{measured\_\allowbreak{}j\ =\ None}, \texttt{board\_\allowbreak{}metered\ =\ false}).
  \href{https://research.openie.dev/data/\#rce-jepa}{research.openie.dev/data/}.
  Full.
\item
  {[}B35{]} Open Interface Engineering, ``FormEBM-v2 contrastive
  multi-step margin latent energy scorer (H5, H11) and H15-Iterative
  bounded checker-feedback residual closure (40/40) with DNULLMM1 Gear-0
  v2 crystallization'', registered 23 November 2026
  (\texttt{rce-\allowbreak{}jepa-\allowbreak{}form-\allowbreak{}ebm-\allowbreak{}v2-\allowbreak{}2026-\allowbreak{}11-\allowbreak{}23},
  \texttt{sha256:\ 128f17e6\allowbreak{}9361ea89.\allowbreak{}.\allowbreak{}.\allowbreak{}};
  \texttt{rce-\allowbreak{}jepa-\allowbreak{}h15-\allowbreak{}iterative-\allowbreak{}2026-\allowbreak{}11-\allowbreak{}23},
  \texttt{sha256:\ a7774d3b\allowbreak{}d03ae450.\allowbreak{}.\allowbreak{}.\allowbreak{}}) on \texttt{openIE-\allowbreak{}dev/\allowbreak{}rce} at
  commit \texttt{1762fb4b} (Linux x86\_64 binary
  \texttt{sha256:\ 33146449\allowbreak{}6009d98a.\allowbreak{}.\allowbreak{}.\allowbreak{}}), AMD Ryzen AI Max+ 395 / Radeon
  8060S Linux powercap RAPL \texttt{reported\_\allowbreak{}j} receipts
  (\texttt{measured\_\allowbreak{}j\ =\ None}, \texttt{board\_\allowbreak{}metered\ =\ false}).
  \href{https://research.openie.dev/living/rce-jepa/}{research.openie.dev/living/rce-jepa/}
  ·
  \href{https://research.openie.dev/data/\#rce-jepa}{research.openie.dev/data/}.
  Full.
\item
  {[}B36{]} Open Interface Engineering, ``FormEBM-v3 span-normalized
  two-head decomposed contrastive energy scorer (H5, H11, and zero-shot
  OOD H15 cross-span transfer)'', registered 23 November 2026
  (\texttt{rce-\allowbreak{}jepa-\allowbreak{}form-\allowbreak{}ebm-\allowbreak{}v3-\allowbreak{}2026-\allowbreak{}11-\allowbreak{}23},
  \texttt{sha256:\ cb7ee1ac\allowbreak{}8dc5c594.\allowbreak{}.\allowbreak{}.\allowbreak{}}) on \texttt{openIE-\allowbreak{}dev/\allowbreak{}rce} at
  commit \texttt{1762fb4b} (Linux x86\_64 binary
  \texttt{sha256:\ 386be0c2\allowbreak{}6389550e.\allowbreak{}.\allowbreak{}.\allowbreak{}}; checkpoint
  \texttt{sha256:\ 71e642f1\allowbreak{}bb0b9c06.\allowbreak{}.\allowbreak{}.\allowbreak{}}), AMD Ryzen AI Max+ 395 / Radeon
  8060S Linux powercap RAPL \texttt{reported\_\allowbreak{}j} receipts
  (\texttt{measured\_\allowbreak{}j\ =\ None}, \texttt{board\_\allowbreak{}metered\ =\ false}).
  \href{https://research.openie.dev/living/rce-jepa/}{research.openie.dev/living/rce-jepa/}
  ·
  \href{https://research.openie.dev/data/\#rce-jepa}{research.openie.dev/data/}.
  Full.
\end{enumerate}

\section{The thesis: software construction is pre-``Laws of
Coding''}\label{the-thesis-software-construction-is-pre-laws-of-coding}

Picture Europe in the 1660s. Kepler had published two regularities of
planetary motion in \emph{Astronomia nova} (1609) and a third in
\emph{Harmonice mundi} (1619) {[}T4{]}. Galileo and Huygens had derived
testable conclusions from hypotheses about motion {[}T5, Sec. 1{]}.
Gunners, millwrights and shipwrights built by rule of thumb and long
apprenticeship. The sky had been watched with enormous care for two
thousand years, and yet prediction was a patchwork. Newton, by the
account of George Smith, ``was familiar with seven different approaches
to calculating planetary orbits, all at roughly the same accuracy'', and
calculated positions ``were sometimes off by as much as a fourth of the
width of the Moon'' {[}T5, Sec. 2{]}. There were many good tables and no
single account of why the tables worked.

Then, in 1687, came the \emph{Principia}. A handful of axioms of motion
and one inverse-square force, and from them trajectories that could be
computed in advance, costs that could be reckoned before the cannon was
cast, and errors that pointed to new bodies rather than to bad luck.
Smith summarizes the change: after the \emph{Principia}, ``physical
theory gained primacy over observation for purposes of answering
specific questions about the world'' {[}T5, Sec. 1{]}. Galileo's results
for uniform gravity did not disappear; they became a limit case of
inverse-square gravity, and the theory went on to explain the planets'
deviations from Kepler's rules {[}T5, Secs. 1 and 2{]}.

The premise, in plain words: \textbf{the practice of software
construction, including today's AI-assisted coding by large language
models, is pre-``Laws of Coding''.} Here ``the practice of software
construction'' means the practitioners and systems, human developers and
code-generating models alike, that produce code by empirical trial,
testing and review rather than by derivation from a formal rule set.
That practice today lives in the 1660s. It has careful observation:
\href{mailto:pass@k}{\nolinkurl{pass@k}} (the chance that at least one
of k sampled programs passes the tests), used in Section 4 below, and
SWE-bench Verified (a reviewed set of repository issues), used in
Section 3 below. It has craft: prompt wording, retries, and trial and
error. Outputs of code-generating models are samples from a probability
distribution, so the same request can give a different program tomorrow;
when ChatGPT (gpt-3.5-turbo at its default temperature) was asked the
same coding question five times, for 47.56\% to 75.76\% of problems
across three datasets there was no test on which all five answers gave
the same output, and ``setting the temperature to 0 does not guarantee
determinism in code generation'' (Ouyang et al., 2023, abstract). The
regularities that the human side of the practice does show are, in
Lehman's words, ones that ``arise from the habits and practices of
people and organizations'' {[}T6, Sec. V-A, p.~1069{]}. Hoare and
colleagues named the same gap in 1987: ``What then are the laws of
programming, which provide the formal basis for the profession of
software engineering? Many programmers may be unable to quote even a
single law'' {[}T7, p.~673{]}. Dijkstra had already marked the limit of
the empirical method that the practice still relies on: ``program
testing can be a very effective way to show the presence of bugs, but is
hopelessly inadequate for showing their absence'' {[}T8, Argument
three{]}. There is no agreed conserved quantity for a program, no
equation of motion that a code generator obeys, and cost is measured
after the fact, in tokens and in electricity, rather than computed
before the work begins.

``Laws of Coding'' names the craft layer, defined with the other three
layers below. The RCE-JEPA proposal, stated here as conjecture, is that
this craft can gain a small set of rules from which the behaviour and
the cost of the code it produces can be computed before it runs, by
borrowing from the layers beneath it. The candidates are (L1) a
deterministic transition rule, (L2) an account of information and energy
per transition, (L3) a canonical Form for every program, which is
compression, and (L4) a composition rule under which L1 to L3 survive
when programs are joined. The
\href{https://research.openie.dev/about/rce/}{Reality Code Engine (RCE)}
proposes candidates for each; hypotheses H1 to H15 are the tests, stated
in numerical order in one table in Section 8 below. Nothing in this
section is a measured result unless it carries a number and a source.

The analogy has limits, stated here. Newton described a world nobody
designed; programs are made by people, so their rules are partly chosen
rather than discovered. Mechanics has continuous quantities and exact
conservation of energy; a program's state is discrete and finite, and
the energy a chip spends depends on the hardware, not only on the
program. Rice's theorem {[}T3{]} says that no general procedure decides
any non-trivial property of what an arbitrary program computes, and the
halting problem {[}T9{]} is the best-known case; RCE avoids these limits
only by working in a restricted language whose programs carry
termination bounds, and that restriction has a price in what can be
expressed. Finally, even the \emph{Principia} was not the ``clockwork
universe'': Smith notes that aspect ``was added by Laplace'' {[}T5, Sec.
1{]}; the axioms of motion alone did not settle every problem of
mechanics either, since for bodies that are not point-masses they ``must
be supplemented by further principles'' {[}T5, Sec. 5{]}; and
``Newtonian theory is now seen to hold only to high approximation in
limited circumstances'' {[}T5, Sec. 1{]}. The claim is an analogy, not
an equivalence.

\subsection{Four layers of law, from theory to
art}\label{four-layers-of-law-from-theory-to-art}

Four phrases are easy to blur, and this page keeps them apart. Each
names a different layer, running from theory through science to art, and
each is formalized to a different degree today. They are never used
interchangeably here.


\begin{center}
\openiefig{fig/rce-jepa-fig03-layers.png}{0.94}{https://research.openie.dev/figures/rce-jepa/layers.svg}
\end{center}


\emph{Figure 3. The spectrum. Left to right: ``Laws of Computation''
(theory), ``Laws of Programming Languages'' and ``Laws of Programming''
(science), ``Laws of Coding'' (art). The bar under each layer shows how
formalized it is today, a qualitative judgement explained in the table
below. RCE-JEPA is drawn as a bridge that carries evidence between
layers; it does not merge them. A qualitative diagram, not data. .}

{\footnotesize\setlength{\tabcolsep}{3.8pt}\renewcommand{\arraystretch}{1.15}
\begin{longtable}[]{@{}
  >{\raggedright\arraybackslash}p{(\linewidth - 6\tabcolsep) * \real{0.1285}}
  >{\raggedright\arraybackslash}p{(\linewidth - 6\tabcolsep) * \real{0.2905}}
  >{\raggedright\arraybackslash}p{(\linewidth - 6\tabcolsep) * \real{0.2905}}
  >{\raggedright\arraybackslash}p{(\linewidth - 6\tabcolsep) * \real{0.2905}}@{}}
\toprule\noalign{}
\begin{minipage}[b]{\linewidth}\raggedright
\textbf{Layer}
\end{minipage} & \begin{minipage}[b]{\linewidth}\raggedright
\textbf{Definition}
\end{minipage} & \begin{minipage}[b]{\linewidth}\raggedright
\textbf{Who works on it}
\end{minipage} & \begin{minipage}[b]{\linewidth}\raggedright
\textbf{How formalized today}
\end{minipage} \\
\midrule\noalign{}
\endhead
\bottomrule\noalign{}
\endlastfoot
``Laws of Computation'' & Physical and information limits: what can be
computed at all, and at what cost in information, time and energy. &
Shannon {[}T10{]}; Landauer {[}T11{]}; Bennett and Landauer {[}T12,
T13{]}; Turing {[}T9{]} and Rice {[}T3{]}; Lloyd {[}T14{]}; Wolpert
{[}T15{]}; Carette, James and Sabry {[}T16{]}. & Exact theorems for
limits: information in bits, the erasure floor \(k_BT\ln 2\),
undecidability. No unified account yet predicts the cost of a given
program on a given machine; Bennett and Landauer note their limits are
``terribly far away from the real limits of modern technology'' {[}T13,
p.~48{]}. \\
``Laws of Programming Languages'' & Semantics and type theory: the rules
that fix what a program in a language means. & Hoare's axiomatic
semantics {[}T17{]}; Milner {[}T18{]}; Harel {[}T19{]}; Lamport
{[}T20{]}; Berry and Gonthier {[}T21{]}; Halbwachs et al.~{[}T22{]};
Leroy's verified compiler and Letouzey's extraction (Section 1 below). &
Fully formal for some languages and compilers, machine-checked in
CompCert (Section 1 below). Formal meaning is a property a language is
designed to have, not one every language has. \\
``Laws of Programming'' & Algebraic laws of correct program
construction: equations by which programs are transformed and derived
from specifications. & Hoare, Hayes, He, Morgan, Roscoe, Sanders,
Sorensen, Spivey and Sufrin {[}T7{]}; Dijkstra's derivation of programs
and Bird's calculus of lists (Section 1 below). & Nearly one hundred
laws for a small language {[}T7{]}; by its authors' own account ``still
a long way from knowing how to apply them directly to the design of
correct and efficient programs on the scale required by modern
technology'' {[}T7, p.~686{]}. \\
``Laws of Coding'' & The empirical craft of producing working code by
people or models. This is where AI coding sits today. & Brooks
{[}T23{]}; Lehman {[}T6{]}; Hindle et al.~{[}T24{]}; the reports in
Section 3 below. & Empirical regularities and benchmarks, no laws.
Outputs of AI coders are samples, so the same request can give a
different program (Ouyang et al., 2023, quoted above). \\
\end{longtable}
}

\textbf{How RCE-JEPA relates to the layers.} RCE-JEPA is a bridge, not a
merger. It works in the ``Laws of Coding'' layer, where a learned model
proposes code, and it ties each proposal to the layers beneath: the Form
has a fixed meaning (``Laws of Programming Languages''), accepted Forms
can be transformed by stated rules (``Laws of Programming''), and their
cost is to be accounted against physical limits (``Laws of
Computation''); RCE counts that cost today in integer units, which are
not joules. Each hypothesis below names the layer it tests. If the
bridge holds, the craft gains rules borrowed from the science and the
theory; the layers stay distinct.

Mechanics has its \emph{Principia}. Computation does not, yet. The
pieces of ``Laws of Computation'' listed above are exact, but they have
not been joined into one formal account that predicts behaviour and cost
together, the way the \emph{Principia} joined Kepler's and Galileo's
regularities. This is stated as an open problem, not a solved one.

Nothing in that gap is a failure of the people who came before. Each era
did what its machines, its theory of information and its craft of
building and optimizing programs allowed. The timeline below says what
each era had to work with, with figures from the sources named, and what
it achieved.


\begin{center}
\openiefig{fig/rce-jepa-fig04-eras.png}{0.94}{https://research.openie.dev/figures/rce-jepa/eras.svg}
\end{center}


\emph{Figure 4. Five eras on the road toward ``Laws of Computation''.
Top lane: hardware, with figures from the sources in the table below.
Middle lane: information theory. Bottom lane: applied computer science
for developing and optimizing algorithms (compilers, complexity,
synthesis, verification). The right-hand box is the open question: a
unified ``Laws of Computation''. A historical diagram, not data. .}

{\footnotesize\setlength{\tabcolsep}{3.8pt}\renewcommand{\arraystretch}{1.15}
\begin{longtable}[]{@{}
  >{\raggedright\arraybackslash}p{(\linewidth - 8\tabcolsep) * \real{0.0673}}
  >{\raggedright\arraybackslash}p{(\linewidth - 8\tabcolsep) * \real{0.2332}}
  >{\raggedright\arraybackslash}p{(\linewidth - 8\tabcolsep) * \real{0.2332}}
  >{\raggedright\arraybackslash}p{(\linewidth - 8\tabcolsep) * \real{0.2332}}
  >{\raggedright\arraybackslash}p{(\linewidth - 8\tabcolsep) * \real{0.2332}}@{}}
\toprule\noalign{}
\begin{minipage}[b]{\linewidth}\raggedright
\textbf{Era}
\end{minipage} & \begin{minipage}[b]{\linewidth}\raggedright
\textbf{Hardware (memory, operations, energy)}
\end{minipage} & \begin{minipage}[b]{\linewidth}\raggedright
\textbf{Information theory}
\end{minipage} & \begin{minipage}[b]{\linewidth}\raggedright
\textbf{Applied computer science for algorithms}
\end{minipage} & \begin{minipage}[b]{\linewidth}\raggedright
\textbf{What it made possible, and where it stopped}
\end{minipage} \\
\midrule\noalign{}
\endhead
\bottomrule\noalign{}
\endlastfoot
1940s to 1950s & Single-copy machines whose memory ``was usually too
small'' and which were ``usually too slow'' {[}T8{]}. On the IBM 704, a
47-statement FORTRAN (FORmula TRANslation) program compiled ``in six
minutes, producing about 1000 instructions'' {[}T25, PDF p.~2{]}. &
Shannon (1948) makes information a measurable quantity: ``The
fundamental problem of communication is that of reproducing at one point
either exactly or approximately a message selected at another point''
{[}T10, Introduction, p.~1{]}. & The first optimizing compilers
{[}T25{]}. & Translation from formula to machine code became automatic.
Every instruction and every word of memory was dear, so effort went to
squeezing programs into machines, as Dijkstra recalls {[}T8{]}. \\
1960s & The four models of IBM System/360 (1964) differed by a factor of
16 in memory capacity and 50 in performance while sharing one
instruction set {[}T26{]}. & Landauer (1961) ties erasure of information
to heat {[}T11{]}; Solomonoff, Kolmogorov and Chaitin define complexity
as shortest description {[}T27, T28, T29{]}. & Hoare's axiomatic basis
for proving programs correct (1969) {[}T17{]}. & Correctness and
information got exact definitions on paper. Dijkstra (1972): as ``the
power of available machines grew by a factor of more than a thousand,
society's ambition to apply these machines grew in proportion''
{[}T8{]}, and the software crisis followed. \\
1970s to 1980s & Alto (1973): a writable control store of 4,096 words of
32 bits; VAX-11/780 (1977): a control store of 5,120 words of 96 bits;
Moore's 1975 projection of a doubling in transistor density every two
years {[}T26{]}. & Bennett (1973) shows that general-purpose computing
machines ``may be made logically reversible at every step, while
retaining their simplicity and their ability to do general
computations'' {[}T30, abstract, p.~525{]}; Bennett and Landauer (1985)
set out the physical limits {[}T13{]}. & Brooks (1986) separates
essential from accidental complexity {[}T23{]}; Hoare and eight
co-authors give nearly one hundred algebraic laws for a programming
language {[}T7{]}; Lehman states empirical laws of software evolution
{[}T6{]}. & Programs became algebraic objects. Hoare et al.~are candid
about the distance left: ``we are still a long way from knowing how to
apply them directly to the design of correct and efficient programs on
the scale required by modern technology'' {[}T7, p.~686{]}. \\
1990s to 2000s & Instruction-level parallelism and faster transistors
gave about 50\% annual performance growth from 1986 to about 2002
{[}T26{]}. & Lloyd (2000) bounds operations per second by energy
{[}T14{]}. & McKee (1995) proposes fundamental laws of computation
{[}T31{]}; machine-checked compilers and kernels, CompCert and seL4
(Section 1 below). & Proof reached real systems, at a price: seL4's
proof took about 20 person-years (Section 1 below). \\
2010s to today & Dennard scaling ``faded to almost nothing by 2012''; a
32-kilobyte cache access costs about 200 times the energy of a 32-bit
add, and off-chip DRAM (dynamic random-access memory) about 10 times a
last-level cache access; Google's TPU (Tensor Processing Unit) v1
performs 256 x 256 multiply-accumulates per cycle with 24 megabytes of
local memory, at more than 80 times a CPU's energy efficiency on its
workload {[}T26{]}. A 65-billion-parameter language model used about 3
to 4 joules per output token (Samsi et al., Section 4 below). &
Wolpert's stochastic thermodynamics of computation {[}T15{]}; prediction
shown equivalent to compression for language models {[}T32{]}; Carette,
James and Sabry rebuild models of computation on conservation of
information {[}T16{]}. & Learned code generators; model plus verifier
loops; deterministic simulation testing (see who holds these views,
below). & Checking is cheap and exact where proposers are cheap and
plentiful. A Python matrix multiply rewritten for the hardware ran
``more than 62,000x faster'' {[}T26{]}: the distance between a program
and its machine is now large, measurable and worth closing. \\
\end{longtable}
}

\textbf{Why progress is possible now (opinion, to be tested).} Three
things that were never available together are available today.
Machine-checked verifiers exist for real systems, so a rule can be
enforced rather than hoped for. Learned models can propose candidate
programs by the thousand, so the search for a program that satisfies a
rule is affordable. And energy can be read from counters on ordinary
hardware, so cost can be measured against prediction rather than
guessed. Each earlier era lacked at least one of the three. None of this
makes ``Laws of Computation'' inevitable; it makes them testable, which
is the hypotheses' purpose below.

\subsection{Symbols used in this
section}\label{symbols-used-in-this-section}

\begin{itemize}
\tightlist
\item
  \(P=(S,E,\delta,D,\lambda)\): a program as a functional state machine.
  \(S\) is a finite set of states; \(E\) a finite set of events
  (inputs); \(\delta: S\times E\to S\) the transition function; \(D\)
  the set of derivations, values computed from a state by pure
  functions; \(\lambda: S\times E\to O\) the output function into a
  finite output set \(O\).
\item
  \(\Omega_n\): the set of canonical Forms of size at most \(n\)
  symbols. A Form is the checker-accepted canonical representation of a
  program in RCE (Section 8).
\item
  \(X\): the set of surface programs (the text people or models write).
  \(\kappa: X\to\Omega\): the canonicalization map from surface program
  to Form.
\item
  \(H(\cdot)\): Shannon entropy in bits {[}T10{]}. \(K(x)\): Kolmogorov
  complexity, the length in bits of the shortest program that outputs
  \(x\) on a fixed universal machine {[}T27, T28{]}.
\item
  \(k_B=1.380649\times10^{-23}\ \mathrm{J\,K^{-1}}\): Boltzmann's
  constant. \(T\): temperature in kelvin (K). J: joule.
\item
  JEPA: Joint Embedding Predictive Architecture {[}LeCun 2022, cited
  above{]}. MDL: minimum description length {[}T33{]}. EBM: energy-based
  model.
\end{itemize}

\subsection{Part 1. Computation as a finite information
system}\label{part-1.-computation-as-a-finite-information-system}

Strip a program of its syntax, its comments, its indentation and its
author's habits, and what remains is a machine that moves from one state
to the next, one event at a time. That machine is finite, and a finite
thing can be counted.

\textbf{Definition.} A program is a deterministic functional state
machine \(P=(S,E,\delta,D,\lambda)\) with \(S\) and \(E\) finite and
\(\delta\) a total function. Determinism means \(\delta\) is a function,
not a relation: one state and one event give exactly one next state.

\textbf{Bound on information content.} If a Form is drawn from
\(\Omega_n\) by any distribution, its entropy is at most that of the
uniform distribution: for \(n\) outcomes, ``H is a maximum and equal to
log n when all the p\textsubscript{i} are equal'' {[}T10, Sec. 6,
property 2, p.~11{]}:

\[
H(P)\;\le\;\log_2\lvert\Omega_n\rvert\ \text{bits}.
\]

For an alphabet of \(a\) symbols,
\(\lvert\Omega_n\rvert\le\sum_{k\le n}a^k\), so \(H(P)\le n\log_2 a+1\)
bits. Every Form, however long its behaviour, is a finite message.

\textbf{Compression and understanding.} Kolmogorov complexity \(K(x)\)
{[}T27, T28{]} measures the shortest description of an object.
Kolmogorov takes as the ``relative complexity'' of an object y given x
``the minimal length l(p) of the `program' p for obtaining y from x''
{[}T28, Sec. 3, p.~165{]}; Solomonoff's induction {[}T27{]} and
Rissanen's MDL {[}T33{]} turn ``prefer the shortest description'' into a
rule for prediction and model choice. The bridge to learning is
Shannon's source coding theorem {[}T10, Sec. 9, Theorem 9, p.~16{]}: a
good code spends about \(-\log_2 p(x)\) bits on \(x\) (Shannon's second
proof gives a message of probability \(p_s\) a code of \(m_s\) binary
digits with \(\log_2(1/p_s)\le m_s<1+\log_2(1/p_s)\) {[}T10, Sec. 9,
p.~17{]}), so a better predictor is a better compressor. Delétang et
al.~state it directly: ``maximizing the log2-likelihood (of the data) is
equivalent to minimizing the number of bits required per message''
{[}T32, Sec. 1{]}. \(K\) itself is not computable: Kolmogorov already
noted that the conditional complexity ``cannot be effectively calculated
(generally recursive) even if it is known to be finite for all x and y''
{[}T28, Sec. 3, p.~165{]}, and Schmidhuber notes ``there is no general
way of finding the shortest program computing any given data'' {[}T34,
Sec. 1{]}. In a restricted language with a canonical Form, the Form
length is a computable upper bound on \(K\) up to an additive constant.

\subsection{Part 2. Compression into Form and
template}\label{part-2.-compression-into-form-and-template}

Ask a hundred programmers to write a shopping cart and you will get a
hundred programs. Underneath the hundred there are perhaps three or four
ideas. The surface is a crowd; the Form is the face it shares.

\textbf{Equivalence classes.} Define
\(x\sim y \iff \kappa(x)=\kappa(y)\). Each Form \(f\in\Omega\) is the
representative of the class \([f]=\kappa^{-1}(f)\) of
\(m_f=\lvert[f]\rvert\) surface programs. Because \(\kappa\) is a
function, the chain rule of entropy, \(H(x,y)=H(x)+H_x(y)\) {[}T10, Sec.
6, property 5, p.~12{]}, gives

\[
H(X)\;=\;H(\kappa(X))\;+\;H\big(X\mid\kappa(X)\big).
\]

A predictor that outputs the Form rather than the surface text is spared
the second term. If surface variants within a class are equally likely,
the saving per Form is

\[
\Delta H_f=\log_2 m_f\ \text{bits},\qquad \rho_f=\frac{\lvert x\rvert}{\lvert f\rvert}\quad(\text{bytes of surface code per byte of Form}).
\]

This is why the predictor works on the quotient space
\(X/\!\sim\;\cong\;\Omega\): it learns a smaller, cleaner target, and
every point in it is something the checker can judge.

\textbf{Templates as low-dimensional manifolds.} A template or a no-code
schema is a Form with holes. If a template has \(h\) holes each taking
one of \(v\) values, it covers \(v^h\) Forms and is addressed by
\(h\log_2 v\) bits, typically far fewer than
\(\log_2\lvert\Omega_n\rvert\). Program space is vast; the programs
people actually want lie on thin sheets within it. This is a conjecture,
consistent with the observation that code is more predictable than
natural language {[}T24{]} and with learned code embeddings in which
``semantically similar objects are mapped to close vectors'' {[}T35,
Sec. 1{]}.

\textbf{Corpus measurements.} H1 has a published floor on the
hand-written population. The median of \(\rho\) = b(code)/b(P) over 185
programs is 9.34 against the emitted Python, above the threshold of 5.
The run was registered before it was measured. The data section of the
proposer above gives the interquartile range, the version-1 figure of
9.31 over 173 programs, and the compressed-description check. The saving
\(\Delta H\) on a lift corpus is still planned. That corpus is empty
until the lift pipeline is public. The 9.34 figure is a measured floor
on hand-written programs. It is not a measurement of lifted or
model-written programs.

\subsection{Part 3. Kinematics: a Newtonian form of
computing}\label{part-3.-kinematics-a-newtonian-form-of-computing}

Throw a stone and, given where it starts, how fast, and what pushes on
it, Newton tells you where it lands. A deterministic program has the
same character, and more strictly: no air, no wind, no rounding of the
world. Given where it starts and what events arrive, its entire future
is fixed.

\textbf{Trajectory.} For initial state \(s_0\) and event stream
\(e_0,e_1,\dots,e_{T-1}\),

\[
s_{t+1}=\delta(s_t,e_t),\qquad t=0,\dots,T-1.
\]

The correspondence proposed here, as analogy:

\begin{itemize}
\tightlist
\item
  state \(s_t\) plays the part of position;
\item
  an event \(e_t\) plays the part of an applied impulse;
\item
  \(\delta\) plays the part of the equation of motion;
\item
  the event log is a record from which the trajectory can be replayed
  exactly, which is reversibility in the record even when \(\delta\)
  discards information; Bennett shows that ``it is possible to reprogram
  any deterministic computation as a sequence of logically reversible
  steps, provided the computation is allowed to save a copy of its
  input'' {[}T12{]};
\item
  step count or fuel \(t\) plays the part of time;
\item
  energy per transition, in joules, is the cost.
\end{itemize}

\textbf{Energy floor.} Landauer argued that logical irreversibility ``is
associated with physical irreversibility and requires a minimal heat
generation, per machine cycle, typically of the order of kT for each
irreversible function'' {[}T11, abstract, p.~183{]}; resetting a bit
cuts the entropy by k log\textsubscript{e} 2 and so supplies ``0.6931 kT
per restored bit to the surroundings'' {[}T11, Sec. 4, p.~187{]}.
Landauer's principle, as Bennett later restates it, holds that a
logically irreversible operation such as erasure ``must be accompanied
by a corresponding entropy increase'' in the environment {[}T12,
abstract{]}. The floor per erased bit is

\[
E_{\min}=k_B T\ln 2\approx 2.87\times10^{-21}\ \mathrm{J}\quad\text{at } T=300\ \mathrm{K}.
\]

If transition \(t\) erases \(b_t\) bits, then
\(E_{\text{total}}\ge k_B T\ln 2\sum_t b_t\). Lloyd adds an upper limit
on speed: ``the speed with which a physical device can process
information is limited by its energy'' {[}T14, abstract{]}, at most
\(2E/\pi\hbar\) operations per second for energy \(E\), where
\(\hbar=1.0546\times10^{-34}\ \mathrm{J\,s}\) is the reduced Planck
constant. Real chips spend many orders of magnitude above the Landauer
floor per operation; the RCE project has no measured
joules-per-operation figure at the time of writing, and H9 (Section 8
below) is the planned measurement. Kolmogorov flagged the missing cost
dimension in 1965: his complexity measures ``do not allow for the
`difficulty' of preparing a program p for passing from an object x to an
object y'' {[}T28, Sec. 4, p.~167{]}. Wolpert, quoting Boyd, notes as a
caution that ``different implementations of the same computation (input
to output) can have radically different thermodynamic consequences''
{[}T15, Sec. I.B{]}: energy belongs to the program and the machine
together.

\textbf{Definition (kinematic program).} \(P\) is kinematic if, for
every \(s_0\) and every finite event stream,

\[
\text{(i)}\ \operatorname{Var}[s_T]=0,\qquad \text{(ii)}\ E_{\text{total}}=\sum_{t=0}^{T-1}c(s_t,e_t)\ \le\ c_{\max}\,F
\]

where \(c: S\times E\to\mathbb{R}_{\ge0}\) is a cost table in joules for
a stated machine, \(c_{\max}\) its largest entry, and \(F\) the fuel
bound (maximum steps) from the termination proof. Condition (i) says the
trajectory has zero variance; condition (ii) says its cost can be
computed, or at least bounded, before it runs. RCE's own cost table, in
Appendix C.6.4 of its specification, is in integer units, not joules; a
table in joules per unit for a stated machine is a proposal in the
specification repository, and H9 is the test.

\textbf{Contrast.} A stochastic generator returns a sample
\(x\sim p_\theta(x\mid\text{prompt})\) from a model with parameters
\(\theta\). Its output is a distribution, its correctness is a
probability, and its cost is known only after sampling. RCE-JEPA keeps
the generator stochastic but makes the thing it generates kinematic, and
puts a deterministic checker between them.

\subsection{Part 4. A world model whose physics is known
exactly}\label{part-4.-a-world-model-whose-physics-is-known-exactly}

Most world models are built in fog. A robot learns physics from blurry
cameras; a game agent infers rules it cannot read. Code is the rare
world where the rules are written down.

\textbf{World model.} In the JEPA sense {[}LeCun 2022, cited above{]}, a
world model is a predictor \(\hat z_{t+1}=g_\phi(z_t,a_t)\) of the next
latent representation \(z_{t+1}=\mathrm{enc}(s_{t+1})\) from the current
one and an action \(a_t\), with learned parameters \(\phi\). For code
the world is the Form's state space and the action is an event, so the
target is computed exactly by \(\delta\):

\[
\mathcal{L}(\phi)=\mathbb{E}\,\big\lVert g_\phi\big(\mathrm{enc}(s_t),e_t\big)-\mathrm{enc}\big(\delta(s_t,e_t)\big)\big\rVert^2 .
\]

There is no aleatoric (irreducible) noise in the target. MuZero starts
from the observation that lookahead planning algorithms ``all rely on
knowledge of the environment's dynamics, such as the rules of the game
or an accurate simulator'', and learns a model instead {[}T1, Sec. 1{]};
for code, those rules are \(\delta\), available in full.

\textbf{Contrast with models that learn the physics.} CWM, the Code
World Model {[}B23{]}, is a decoder mid-trained on Python execution
traces and on agent trajectories. Trace prediction in that report is a
learned interpreter, and the authors call it a ``first step'', with a
neural debugger left as future work. V-JEPA 2 {[}B1{]} learns its
physics from video, including 62 hours of robot video in the reported
summary, and the action-conditioned predictor uses fewer than 62 hours
of unlabeled robot video. Its model predictive control took about 16
seconds per action. In the Form, \(\delta\) is the physics, and it is
exact, so there is no aleatoric noise to fit. What a proposer still has
to learn is set out in Section 7 below: intent to Form, search guidance,
and the residual the ladder does not close.

\textbf{Intent decoding as an inverse problem.} Given an intent \(I\)
stated as constraints, find a Form whose trajectories satisfy it. This
is written as energy minimization in the sense of energy-based models
{[}T36, T2{]}:

\[
\mathcal{E}(I,f)=\alpha\,V(f)+\beta\,d\big(\tau(f),\sigma(I)\big)+\gamma\,C(f),\qquad z^\star=\arg\min_{z}\ \mathcal{E}\big(I,\mathrm{dec}(z)\big).
\]

Here \(V(f)\) is the number of checker violations (dimensionless);
\(\tau(f)\) the set of trajectories of \(f\) on the intent's example
events; \(\sigma(I)\) the trajectories the intent requires; \(d\) a
distance between them (dimensionless count of mismatched states);
\(C(f)\) the predicted energy per run in joules; \(\alpha,\beta\ge0\)
dimensionless weights and \(\gamma\ge0\) in \(\mathrm{J}^{-1}\);
\(\mathrm{dec}\) the decoder from latent point to Form. The minimizer is
decoded, then checked; the checker, not \(\mathcal{E}\), decides
acceptance. Du and Mordatch observe that ``summation of two or more
energies corresponds to satisfying all their goals or constraints''
{[}T36, Sec. 1{]}, which is why the three terms can be added. LeCun et
al.~add a caution: ``combining two, separately trained energy-based
models is not straightforward: there is no a priori guarantee that their
energy scales are commensurate'' {[}T2, Sec. 1.3, p.~6{]}. The weights
\(\alpha,\beta,\gamma\) must therefore be calibrated, not assumed, and
their values are reported with each experiment.

\subsection{Part 5. Consequences as testable
predictions}\label{part-5.-consequences-as-testable-predictions}

A thesis that cannot lose is not a thesis. Each part above maps to a
test that could prove it wrong.

\begin{itemize}
\tightlist
\item
  Part 1 and T2 (compression into Form) are tested by H1, H3 and H10.
\item
  Part 3 (kinematics, energy accounting) is tested by H2, H4, H8 and H9.
\item
  Part 4 (world model, inverse problem) is tested by H5 to H7.
\item
  The proposer is tested by H11 to H13, on the residual of the
  deterministic ladder that H14 measures. H15 compares proposer classes
  on that residual and is unfinished. All fifteen hypotheses, with
  metrics and thresholds, are in one table in numerical order in Section
  8 below.
\end{itemize}


\begin{center}
\openiefig{fig/rce-jepa-fig05-thesis.png}{0.94}{https://research.openie.dev/figures/rce-jepa/thesis.svg}
\end{center}


\emph{Figure 5. Information and kinematics. Left: program space, where
many surface programs X express the same behaviour. Middle:
canonicalization κ collapses each class onto one Form on a
low-dimensional manifold Ω, saving ΔH = log2 m bits for a class of m
equally likely variants. Right: a Form moves deterministically, s(t+1) =
δ(s(t), e(t)), each transition carrying a cost c in joules, with the
total bounded by c\_max · F before execution. A conceptual diagram, not
data. .}

\subsection{The provable chain: hypothesis, layer,
evidence}\label{the-provable-chain-hypothesis-layer-evidence}

A claim about ``Laws of Coding'' is only as strong as its weakest link.
The table names, for each hypothesis, which of the four layers of law it
tests, then follows it through three layers of implementation: the
algorithm (the mathematics of the Form and the predictor), the software
environment (checker, corpus, harness) and the hardware (the machine and
its meter). Each link carries one of four labels. \textbf{{[}C{]} cited
literature:} a published result, read as stated in the ledger.
\textbf{{[}M{]} measured data:} a number measured by the OpenIE authors
in a run registered before it ran, with its source in a public
repository. \textbf{{[}X{]} measured, not supported:} a registered
measurement whose result does not support the hypothesis as intended.
\textbf{{[}P{]} planned experiment:} a test not yet run. A link that is
partly measured carries both {[}M{]} and {[}P{]}. Four hypotheses have
measured links, as the Status section at the end of the page reports.
H15 has no measured link. The four are: H1 on the hand-written programs
of the Form corpus (the data section of the proposer above), H7 on its
fuzzed arm only, H8 on six targets, and H14, whose three registered runs
{[}B30, B31, B32{]} support it by its registered rule (Section 8 below).
A measured link records that a registered run exists and what it found;
it is not by itself a verdict of support. The chain is a plan for
turning every remaining {[}P{]} into {[}M{]} or {[}X{]}, and any {[}P{]}
can still break it.


\begin{center}
\openiefig{fig/rce-jepa-fig06-chain.png}{0.94}{https://research.openie.dev/figures/rce-jepa/chain.svg}
\end{center}


\emph{Figure 6. The provable chain for ``Laws of Coding''. Rows:
hypotheses H1 to H14, each tagged with the layer of law it tests. H15 is
not drawn on the figure. The table adds it as planned work, with no
measured link. Columns: implementation layers, algorithm, software
environment, hardware. Blue: cited literature. Green: measured data.
Vermilion: measured, not supported. Orange: planned experiment. Grey:
not applicable. The table below gives the content and source of every
cell. .}

{\scriptsize\setlength{\tabcolsep}{3.0pt}\renewcommand{\arraystretch}{1.14}
\begin{longtable}[]{@{}
  >{\raggedright\arraybackslash}p{(\linewidth - 10\tabcolsep) * \real{0.0286}}
  >{\raggedright\arraybackslash}p{(\linewidth - 10\tabcolsep) * \real{0.1238}}
  >{\raggedright\arraybackslash}p{(\linewidth - 10\tabcolsep) * \real{0.1095}}
  >{\raggedright\arraybackslash}p{(\linewidth - 10\tabcolsep) * \real{0.2429}}
  >{\raggedright\arraybackslash}p{(\linewidth - 10\tabcolsep) * \real{0.2476}}
  >{\raggedright\arraybackslash}p{(\linewidth - 10\tabcolsep) * \real{0.2476}}@{}}
\toprule\noalign{}
\begin{minipage}[b]{\linewidth}\raggedright
\textbf{ID}
\end{minipage} & \begin{minipage}[b]{\linewidth}\raggedright
\textbf{Hypothesis}
\end{minipage} & \begin{minipage}[b]{\linewidth}\raggedright
\textbf{Layer of law tested}
\end{minipage} & \begin{minipage}[b]{\linewidth}\raggedright
\textbf{Algorithm}
\end{minipage} & \begin{minipage}[b]{\linewidth}\raggedright
\textbf{Software environment}
\end{minipage} & \begin{minipage}[b]{\linewidth}\raggedright
\textbf{Hardware}
\end{minipage} \\
\midrule\noalign{}
\endhead
\bottomrule\noalign{}
\endlastfoot
H1 & Form much smaller than source & ``Laws of Programming Languages''
(canonical meaning) and ``Laws of Coding'' & {[}C{]} Chain rule:
predicting the Form saves \(H(X\mid\kappa(X))\) bits {[}T10{]} & {[}M{]}
Median \(\rho\) of 9.34 over 185 hand-written programs of the Form
corpus against the emitted Python, above the threshold of 5, registered
before the run (data section of the proposer above). {[}P{]} \(\rho\) on
the canonical serialization over the lift corpus, empty until the lift
pipeline is public & Not applicable (a count of bits) \\
H2 & Checker cheap relative to generation & ``Laws of Programming
Languages'' & {[}C{]} Validating a result can be easier than proving the
generator: CompCert's register allocator (Leroy, Section 4 below) &
{[}P{]} Checker wall time and cycles per candidate & {[}P{]}
\(C_{\text{check}}/E_{\text{model}}\) in joules on a named meter \\
H3 & Latent prediction reaches useful acceptance & ``Laws of Coding'' &
{[}C{]} Prediction in representation space (I-JEPA and V-JEPA 2, Section
4 below; LLM-JEPA, Huang, LeCun and Balestriero, Section 5 below) &
{[}P{]} \href{mailto:pass@1}{\nolinkurl{pass@1}} on held-out
intent-to-Form tasks & Not applicable \\
H4 & Energy per accepted program falls & ``Laws of Computation'' &
{[}C{]} Landauer floor \(k_BT\ln 2\) per erased bit {[}T11, T12{]} &
{[}P{]} \(J_{\text{acc}}\) for both arms in one harness & {[}C{]} About
3 to 4 J per token for a 65B model (Samsi et al.); memory access
dominates arithmetic {[}T26{]}. {[}P{]} Device meter, idle subtracted \\
H5 & Energy ranking favours valid Forms & ``Laws of Coding'' & {[}C{]}
Energies add to combine constraints {[}T36{]} & {[}P{]} AUROC of energy
for accepted against requires change & Not applicable \\
H6 & Reference apps are encodable & ``Laws of Programming Languages'' &
{[}C{]} No general procedure decides non-trivial semantic properties
{[}T9, T3{]}, hence the restricted language & {[}P{]} 7 of 7 reference
apps as checked Forms & {[}P{]} Replay of the same Forms across machines
(with H8) \\
H7 & Security independent of the model & ``Laws of Programming
Languages'' & {[}C{]} After about six CPU-years, Csmith found no
wrong-code errors in CompCert (Yang et al., Section 4 below) & {[}M{]}
Fuzzed arm: 0 violations of five stated invariants over 1,275
checker-accepted mutants in two runs, each registered before it ran
(\texttt{rce/\allowbreak{}bench/\allowbreak{}adversarial/\allowbreak{}}). {[}P{]} Model-generated adversarial
Forms, an arm that does not exist yet & {[}C{]} Speculative-execution
side channels lie outside an instruction set's semantics {[}T26{]}; out
of the Form's scope, stated \\
H8 & Determinism & ``Laws of Programming Languages'' and ``Laws of
Computation'' & {[}C{]} Any deterministic computation can be made
logically reversible given a copy of its input {[}T12{]} & {[}M{]}
Bit-identical replay of 10,500 seeded event streams with 0 mismatching
lines against the reference, registered before the run
(\texttt{rce/\allowbreak{}bench/\allowbreak{}determinism/\allowbreak{}}) & {[}C{]} Interrupts arrive ``at
unpredictable and irreproducible moments'' {[}T8{]}. {[}M{]} The same
run on six targets: four instruction sets (x86\_64, aarch64, s390x
big-endian, wasm32) and three operating systems (Linux, macOS,
Windows) \\
H9 & Cost predicted before execution & ``Laws of Computation'' & {[}C{]}
Energy bounds speed {[}T14{]}; cost bound \(c_{\max}F\) (Part 3 above) &
{[}P{]} Cost-table prediction per Form & {[}C{]} Implementations of one
computation differ in thermodynamic cost {[}T15{]}. {[}P{]} Median error
within 20\% of measured joules \\
H10 & Learning on the quotient saves data & ``Laws of Coding'' & {[}C{]}
Better prediction is better compression {[}T32{]}; minimum description
length {[}T33{]} & {[}P{]} Equal acceptance with at most half the
training examples & {[}P{]} Training energy measured for both
predictors \\
H11 & Energy monotone in edit distance & ``Laws of Coding'' & {[}C{]} A
monotone cost ranking makes plans that deviate further from the goal
cost more (Chahe et al., {[}B19{]} in the proposer section) & {[}P{]}
Share of families where energy rises with edit distance & Not
applicable \\
H12 & Decision head calibrated on gate outcomes & ``Laws of Coding'' and
``Laws of Computation'' & {[}C{]} Calibrated typed decisions are
measurable: pooled Choice ECE 0.028 for one hosted decision model
({[}B7{]}) & {[}P{]} ECE per decision type after 10\textsuperscript{4}
receipts, against a frequency table & {[}P{]} Latency and joules per
decision on a named device \\
H13 & Discrete against continuous latent & ``Laws of Coding'' & {[}C{]}
Discrete tokens kept long-horizon prediction stable where continuous
latents accumulated error ({[}B3{]}) & {[}P{]} Planning success at edit
lengths 1, 5 and 20 & Not applicable \\
H14 & Deterministic ladder share & ``Laws of Coding'' & {[}C{]} Exact
gears (lookup, rules, synthesis) run before any model in the RCE paper's
cascade, with falsifier RCE-F1 on the share they close with no model
call & {[}M{]} Share of a registered task set closed by the public
ladder with zero model calls: three registered runs {[}B30, B31, B32{]};
run 3 closes 111 of 151 (73.5\%), at or above the registered 10\%, so
H14 is supported by its registered rule, and 1 of 41 changes from real
history (2.4\%) (Section 8 below) & Not applicable \\
H15 & Proposer classes on the residual. Unfinished & ``Laws of Coding''
& {[}P{]} Four classes on the 40 residual tasks, equal energy, gate
acceptance. No run & {[}P{]} The comparison has not been run & {[}P{]}
Equal energy when a named meter exists. No joule figure is stated \\
\end{longtable}
}

\subsection{Who holds these views and
why}\label{who-holds-these-views-and-why}

Ideas rarely arrive alone. Each part of the thesis echoes arguments
made, for their own reasons, by people who have never heard of RCE. They
are listed as convergence observations: agreement in direction, never
confirmation of the RCE project's claims. Where the OpenIE authors read
a source in full it is quoted with its section; where they did not, the
entry says so and quotes nothing.

\textbf{Those who already think this way.} The idea that programming and
computation should have laws of their own did not originate with the RCE
project, and it is older than most of the field. This page joins a
conversation that has run for at least forty years and lists it as
convergence, not priority.

\begin{itemize}
\tightlist
\item
  \textbf{Hoare and colleagues, ``Laws of Programming'' (1987).} Nine
  authors argued that ``conventional procedural programs are
  mathematical expressions, and that they are subject to a set of laws
  as rich and elegant as those of any other branch of mathematics,
  engineering, or natural science'' {[}T7, p.~672{]}, and gave nearly
  one hundred of them. They also named the gap that remains (quoted in
  the timeline above) {[}T7, p.~686{]}.
\item
  \textbf{Bennett and Landauer, ``The Fundamental Physical Limits of
  Computation'' (1985).} ``We are looking for general laws that must
  govern all information processing, no matter how it is accomplished''
  {[}T13, p.~48{]}. They add a caution this page shares: ``To the extent
  that we have found limits, they are terribly far away from the real
  limits of modern technology'' {[}T13, p.~48{]}. Their result for
  reversible machines: ``There is thus no minimum amount of energy that
  must be expended in order to perform any given computation'' {[}T13,
  p.~51{]}, while erasure and error correction still dissipate, since,
  in their words, ``Any correctional device must therefore dissipate
  some energy'' {[}T13, p.~51{]}.
\item
  \textbf{McKee, ``A few (proposed) fundamental laws of computation''
  (1995).} A short essay in IEEE Computer proposing laws for the field
  {[}T31{]}. The OpenIE authors have not obtained the text, so it is not
  quoted or summarized beyond its title.
\item
  \textbf{Carette, James and Sabry, ``Embracing the laws of physics''
  (2022).} They note that quantum theory and thermodynamics ``imply that
  information is a conserved quantity of physical processes and hence of
  primitive computational operations'', and set out ``to redevelop
  foundational computational models in a way that embraces the principle
  of conservation of information'' {[}T16, abstract, PDF p.~2{]}. Their
  closing summary: ``The entire edifice of computer science including
  its mainstream models of computations, programming languages, and
  logics is founded on classical physics'' {[}T16, Sec. 5.5, p.~45{]}.
\item
  \textbf{Lehman (1980)} states five empirical regularities of software
  evolution for programs embedded in the world {[}T6{]}. The first,
  ``Continuing Change'': a program ``that is used and that as an
  implementation of its specification reflects some other reality,
  undergoes continual change or becomes progressively less useful''
  {[}T6, Table I, p.~1068{]}. He is explicit that these are not physics:
  ``They are, however, clearly not immutable'' in the way physical and
  chemical regularities are, since they ``arise from the habits and
  practices of people and organizations'' {[}T6, Sec. V-A, p.~1069{]},
  the same distance this page keeps between ``Laws of Coding'' and
  ``Laws of Computation''. \textbf{Brooks} distinguished ``essential
  tasks, the fashioning of the complex conceptual structures that
  compose the abstract software entity'' from accidental ones {[}T23,
  PDF p.~1{]}. \textbf{Dijkstra} observed that a computer spans a ratio
  of time grains of ``10\textsuperscript{10} or more'', which makes it
  ``the first to provide us with an environment where highly
  hierarchical artefacts are both possible and necessary'' {[}T8{]}.
\end{itemize}

The groups below each hold one strand of the same idea.

\textbf{Compression as understanding (Part 1, T2).} Solomonoff,
Kolmogorov and Chaitin defined complexity as shortest description
{[}T27, T28, T29{]}; Rissanen made it a model selection rule {[}T33{]};
Hutter's prize and AIXI and Sutskever's talks extend it to intelligence
(not read here; not quoted). Delétang and colleagues at DeepMind argue
that ``a model that compresses well generalizes well'' {[}T32, Sec. 1{]}
and show large language models compress images and audio. Schmidhuber
argues that if the history of the universe were computable, its simplest
explanation ``would be the shortest program that computes it'' {[}T34,
Sec. 1{]}. Disagreement: Delétang et al.~warn that ``scaling beyond a
certain point will deteriorate the compression performance since the
model parameters need to be accounted for'' {[}T32, Sec. 1{]}, and
Chollet argues that ``solely measuring skill at any given task falls
short of measuring intelligence'' {[}T37, abstract{]}; compression of a
fixed corpus is not the same as efficient learning of new skills.

\textbf{Computation as physics (Part 3).} Landauer {[}T11{]}, Bennett
{[}T12, T30{]} and Lloyd {[}T14{]} hold that information processing is a
physical process with energy costs and limits. Lloyd: ``Energy limits
speed. Entropy limits memory.'' {[}T14, Sec. 1.5{]}. Wolpert's
stochastic thermodynamics treats computation, including finite automata,
as non-equilibrium physics: ``One of the most prominent aspects of
computers is that they are inherently nonequilibrium systems'' {[}T15,
abstract{]}; finite automata are treated in Sec. XII {[}T15{]}. Fredkin
and Toffoli's conservative logic {[}T38{]}, Feynman's lectures on
computation, Wheeler's ``it from bit'' and Wolfram's computational
universe point the same way (not read here). Disagreement: Bennett
records objections that Landauer's principle is false, or, being a
consequence of the second law of thermodynamics, either unnecessary or
insufficient to exorcise Maxwell's demon, and answers them {[}T12{]};
the debate is about foundations, not about the size of the floor. A
second disagreement is historical: Landauer wrote ``We believe that
devices exhibiting logical irreversibility are essential to computing''
{[}T11, Sec. 3, p.~186{]}, and Bennett's 1973 construction answered that
they are not, at the price of about twice as many steps and a large
temporary history {[}T30, pp.~525, 529{]}; Bennett then estimated ``kT
ln 2 is, as Landauer has argued {[}1{]}, an approximate lower bound''
only for machines that throw away about one bit per step {[}T30,
p.~532{]}.

\textbf{State machines and deterministic semantics (Part 1, T3).}
Turing's machines {[}T9{]}, Milner's calculus {[}T18{]}, Harel's
statecharts {[}T19{]}, Lamport's TLA, which treats a program as a set of
behaviours over states {[}T20{]}, the synchronous languages Esterel and
Lustre (Berry, Halbwachs) {[}T21, T22{]}, Elliott's functional reactive
programming {[}T39{]}, the Elm architecture, deterministic simulation
testing at FoundationDB and TigerBeetle, and Hoare's and Dijkstra's view
of programs as predicates {[}T17{]} all treat a program as a state
machine with exact semantics. They are cited here as history. The
FoundationDB and TigerBeetle pages and the Elm guide were read in full
for the convergence section below (its FoundationDB, TigerBeetle and Elm
Architecture rows); the others were not read in full for this section,
so none of them are quoted. Limits: Turing {[}T9{]} and Rice {[}T3{]}
bound what can be decided about arbitrary programs; deterministic
semantics is easy to state for a closed language and hard to keep across
real hardware, threads and networks.

\textbf{Programs as data on a low-dimensional manifold (Part 2).} Hindle
et al.~found software ``natural'', more repetitive than English
{[}T24{]} (not read in full here). Alon et al.'s code2vec maps code to
vectors where ``semantically similar objects are mapped to close
vectors'' {[}T35, Sec. 1{]}. Sketch-based synthesis (Solar-Lezama),
DreamCoder's library learning (Ellis et al.~2021, cited in Section 2
below; its supplement describes library learning as ``automatic
refactoring'' {[}Ellis et al.~2021, supplement S3.2, p.~39{]}) and
Allamanis's survey of machine learning on code point the same way
(Solar-Lezama and Allamanis not read here). Disagreement: code2vec's
authors report that ``When given uninformative, obfuscated or
adversarial variable names, the prediction of the label is usually less
accurate'' {[}T35, Sec. 7{]}: learned embeddings can track names rather
than meaning, which is the gap a canonical Form is meant to close.

\textbf{World models and latent prediction (Part 4).} LeCun's JEPA
programme at Meta FAIR (cited above), Ha and Schmidhuber, who ``learn a
compressed spatial and temporal representation of the environment''
{[}T40, abstract{]}, MuZero {[}T1{]}, physics-informed neural networks,
where prior physical knowledge ``can act as a regularization agent that
constrains the space of admissible solutions to a manageable size''
{[}T41, Sec. 1{]}, and Hamiltonian neural networks, motivated by models
whose simulation ``drifts over time to higher or lower energy states''
{[}T42, Sec. 1{]}, all build structure into learned dynamics. Lagrangian
networks (Cranmer), Dreamer (Hafner) and Koopman operator learning share
the aim (not read here). Disagreement: Ha and Schmidhuber found that a
learned world model ``will be exploitable by the controller, even if in
the actual environment such exploits do not exist'' {[}T40, Sec. 4.5{]};
Greydanus et al.~note their method assumes a conserved quantity and so
cannot model friction {[}T42, Sec. 3.2{]}. For code, the checker is the
guard against exploiting the model's errors.

\textbf{Energy-based models (Part 4).} Hopfield {[}T43{]} and Hinton
{[}T44{]} framed learning as shaping an energy surface (not read here);
Du and Mordatch {[}T36{]} show implicit generation by descending a
learned energy, and LeCun (cited above) places energy minimization at
the centre of JEPA inference. Disagreement within the field concerns
training stability and sampling cost, which Du and Mordatch discuss
{[}T36{]}.

\subsection{Lane 3, seen from outside the
project}\label{lane-3-seen-from-outside-the-project}

Section 3 below records programmes that reconceive what computation,
programming or coding is. They are convergence observations. The
commitment section below maps some of the same names onto the seven
commitments, one by one.

\begin{itemize}
\tightlist
\item
  \textbf{Unison.} A definition is the hash of its syntax tree. Names
  are a view. Two texts with the same tree are the same definition. That
  is the identity half of a canonical Form.
\item
  \textbf{egg.} An e-graph holds a set of equivalent programs as one
  structure (Willsey et al., Section 3 below). The Laws of Coding use
  the same quotient: many surface programs, one Form.
\item
  \textbf{Interaction nets.} HVM2 (the Higher Order Company's
  interaction-combinator model) treats computation as local rewriting,
  with parallelism inherent in the model. Section 3 cites the model and
  does not carry that paper's figures.
\item
  \textbf{Proof assistants.} Lean, Dafny and Verus make correctness a
  checkable object. The vericoding benchmark is the measured floor cited
  in Section 3. AlphaProof, DeepSeek-Prover-V2, Kimina-Prover Preview
  and Goedel-Prover use the Lean kernel as the gate. In the reports read
  here, the domain is mathematics.
\item
  \textbf{Elliott, denotational design.} ``The instance's meaning is the
  meaning's instance.'' A semantic model defines correctness and carries
  it to the implementation (Elliott, 2009, revised 2016). It is a manual
  discipline. It is the same demand that a Form have one meaning.
\item
  \textbf{Propagators.} Radul and Sussman (2009) describe computation as
  autonomous machines adding information to shared cells, unifying
  expressions and constraints. The report is theory with a Scheme
  implementation.
\item
  \textbf{Lean, as an engineering programme.} The Lean FRO (Focused
  Research Organization) Year 1 roadmap, August 2023 to July 2024, lists
  a standard library, a module system, Lake and Reservoir, proof
  automation, machine-learning integration and a new code generator.
  Later years are not used here.
\item
  \textbf{Chollet.} Intelligence as skill-acquisition efficiency, with
  the Abstraction and Reasoning Corpus (ARC) as the measure {[}T37{]}.
  The ARC Prize results in Section 3 are the later measurement of that
  programme. Talks after 2019 are not cited on their own.
\end{itemize}


\begin{center}
\openiefig{fig/rce-jepa-fig07-groups.png}{0.94}{https://research.openie.dev/figures/rce-jepa/groups.svg}
\end{center}


\emph{Figure 7. Which communities bear on which part of the thesis. A
filled dot marks a convergence observation, not confirmation. Part 1:
finite information system; T2: compression into Form; T3: kinematics and
energy; T4: world model and inverse problem. .}

\subsection{Thesis references}\label{thesis-references}

Ledger: Full = read in full; Partial = read in part; Unavailable = not
read (paywalled or not retrieved), listed with identifier only and never
quoted.

\begin{enumerate}
\def\labelenumi{\arabic{enumi}.}
\tightlist
\item
  {[}T1{]} J. Schrittwieser et al., ``Mastering Atari, Go, Chess and
  Shogi by Planning with a Learned Model'', Nature 588 (2020).
  \href{https://arxiv.org/abs/1911.08265}{arXiv:1911.08265}. Full.
\item
  {[}T2{]} Y. LeCun et al., ``A tutorial on energy-based learning'',
  2006. \href{http://yann.lecun.com/exdb/publis/pdf/lecun-06.pdf}{PDF}.
  Full.
\item
  {[}T3{]} H. G. Rice, ``Classes of recursively enumerable sets and
  their decision problems'', Trans. AMS 74 (1953).
  \href{https://doi.org/10.1090/S0002-9947-1953-0053041-6}{doi:10.1090/S0002-9947-1953-0053041-6}.
  Unavailable.
\item
  {[}T4{]} M. A. Di Liscia, ``Johannes Kepler'', Stanford Encyclopedia
  of Philosophy.
  \href{https://plato.stanford.edu/entries/kepler/}{plato.stanford.edu}.
  Full.
\item
  {[}T5{]} G. E. Smith, ``Newton's Philosophiae Naturalis Principia
  Mathematica'', Stanford Encyclopedia of Philosophy.
  \href{https://plato.stanford.edu/entries/newton-principia/}{plato.stanford.edu}.
  Full.
\item
  {[}T6{]} M. M. Lehman, ``Programs, life cycles, and laws of software
  evolution'', Proceedings of the IEEE 68(9) (1980).
  \href{https://doi.org/10.1109/PROC.1980.11805}{doi:10.1109/PROC.1980.11805}.
  Full.
\item
  {[}T7{]} C. A. R. Hoare, I. J. Hayes, He Jifeng, C. C. Morgan, A. W.
  Roscoe, J. W. Sanders, I. H. Sorensen, J. M. Spivey and B. A. Sufrin,
  ``Laws of Programming'', Communications of the ACM 30(8), 672-686
  (1987).
  \href{https://doi.org/10.1145/27651.27653}{doi:10.1145/27651.27653};
  \href{https://www.cs.ox.ac.uk/bill.roscoe/publications/20.pdf}{PDF
  copy}. Full.
\item
  {[}T8{]} E. W. Dijkstra, ``The Humble Programmer'', ACM Turing Lecture
  1972, EWD 340.
  \href{https://www.cs.utexas.edu/~EWD/transcriptions/EWD03xx/EWD340.html}{Transcription};
  \href{https://doi.org/10.1145/355604.361591}{doi:10.1145/355604.361591}.
  Full.
\item
  {[}T9{]} A. M. Turing, ``On computable numbers'', Proc. London Math.
  Soc. 42 (1936).
  \href{https://doi.org/10.1112/plms/s2-42.1.230}{doi:10.1112/plms/s2-42.1.230}.
  Unavailable.
\item
  {[}T10{]} C. E. Shannon, ``A mathematical theory of communication'',
  Bell System Technical Journal 27 (1948).
  \href{https://doi.org/10.1002/j.1538-7305.1948.tb01338.x}{doi:10.1002/j.1538-7305.1948.tb01338.x};
  \href{https://people.math.harvard.edu/~ctm/home/text/others/shannon/entropy/entropy.pdf}{reprint}
  (page numbers refer to this reprint). Full (all 55 pp.~of the reprint,
  Parts I to V and Appendices 1 to 7).
\item
  {[}T11{]} R. Landauer, ``Irreversibility and heat generation in the
  computing process'', IBM J. Res. Dev. 5 (1961).
  \href{https://doi.org/10.1147/rd.53.0183}{doi:10.1147/rd.53.0183}.
  Full.
\item
  {[}T12{]} C. H. Bennett, ``Notes on Landauer's principle, reversible
  computation, and Maxwell's Demon'', 2003.
  \href{https://arxiv.org/abs/physics/0210005}{arXiv:physics/0210005}.
  Full.
\item
  {[}T13{]} C. H. Bennett and R. Landauer, ``The Fundamental Physical
  Limits of Computation'', Scientific American 253(1), 48-57 (1985).
  \href{https://doi.org/10.1038/scientificamerican0785-48}{doi:10.1038/scientificamerican0785-48};
  \href{http://physics.bu.edu/~pankajm/PY541/Bennett-Landauer.pdf}{PDF
  copy}. Full (all text pages of the linked copy, pp.~48-56).
\item
  {[}T14{]} S. Lloyd, ``Ultimate physical limits to computation'',
  Nature 406 (2000).
  \href{https://arxiv.org/abs/quant-ph/9908043}{arXiv:quant-ph/9908043}.
  Full.
\item
  {[}T15{]} D. H. Wolpert, ``The stochastic thermodynamics of
  computation'', J. Phys. A 52 (2019).
  \href{https://arxiv.org/abs/1905.05669}{arXiv:1905.05669}. Full.
\item
  {[}T16{]} J. Carette, R. P. James and A. Sabry, ``Embracing the laws
  of physics: Three reversible models of computation'', Advances in
  Computers (2022).
  \href{https://doi.org/10.1016/bs.adcom.2021.11.009}{doi:10.1016/bs.adcom.2021.11.009};
  \href{https://www.cas.mcmaster.ca/~carette/publications/EmbracingTheLaws.pdf}{author
  copy}. Full.
\item
  {[}T17{]} C. A. R. Hoare, ``An axiomatic basis for computer
  programming'', CACM 12 (1969).
  \href{https://doi.org/10.1145/363235.363259}{doi:10.1145/363235.363259}.
  Unavailable.
\item
  {[}T18{]} R. Milner, A Calculus of Communicating Systems, LNCS 92
  (1980).
  \href{https://doi.org/10.1007/3-540-10235-3}{doi:10.1007/3-540-10235-3}.
  Unavailable.
\item
  {[}T19{]} D. Harel, ``Statecharts: a visual formalism for complex
  systems'', Sci. Comput. Program. 8 (1987).
  \href{https://doi.org/10.1016/0167-6423(87)90035-9}{doi:10.1016/0167-6423(87)90035-9}.
  Unavailable.
\item
  {[}T20{]} L. Lamport, ``The temporal logic of actions'', ACM TOPLAS 16
  (1994).
  \href{https://doi.org/10.1145/177492.177726}{doi:10.1145/177492.177726}.
  Unavailable.
\item
  {[}T21{]} G. Berry and G. Gonthier, ``The Esterel synchronous
  programming language'', Sci. Comput. Program. 19 (1992).
  \href{https://doi.org/10.1016/0167-6423(92)90005-V}{doi:10.1016/0167-6423(92)90005-V}.
  Unavailable.
\item
  {[}T22{]} N. Halbwachs et al., ``The synchronous data flow programming
  language LUSTRE'', Proc. IEEE 79 (1991).
  \href{https://doi.org/10.1109/5.97300}{doi:10.1109/5.97300}.
  Unavailable.
\item
  {[}T23{]} F. P. Brooks, ``No Silver Bullet: Essence and Accidents of
  Software Engineering'', IEEE Computer 20(4) (1987; first presented
  1986).
  \href{https://doi.org/10.1109/MC.1987.1663532}{doi:10.1109/MC.1987.1663532};
  \href{https://worrydream.com/refs/Brooks_1986_-_No_Silver_Bullet.pdf}{PDF
  copy}. Full.
\item
  {[}T24{]} A. Hindle et al., ``On the naturalness of software'', ICSE
  (International Conference on Software Engineering) 2012.
  \href{https://doi.org/10.1109/ICSE.2012.6227135}{doi:10.1109/ICSE.2012.6227135}.
  Unavailable.
\item
  {[}T25{]} J. W. Backus et al., ``The FORTRAN Automatic Coding
  System'', Western Joint Computer Conference (1957).
  \href{https://archive.computerhistory.org/resources/text/Fortran/102663113.05.01.acc.pdf}{Computer
  History Museum copy}. Full.
\item
  {[}T26{]} J. L. Hennessy and D. A. Patterson, ``A New Golden Age for
  Computer Architecture'', Communications of the ACM 62(2) (2019).
  \href{https://doi.org/10.1145/3282307}{doi:10.1145/3282307};
  \href{https://cacm.acm.org/research/a-new-golden-age-for-computer-architecture/}{CACM
  web text}. Full.
\item
  {[}T27{]} R. J. Solomonoff, ``A formal theory of inductive inference,
  Part I'', Information and Control 7 (1964).
  \href{https://doi.org/10.1016/S0019-9958(64)90223-2}{doi:10.1016/S0019-9958(64)90223-2}.
  Unavailable.
\item
  {[}T28{]} A. N. Kolmogorov, ``Three approaches to the quantitative
  definition of information'', Problems of Information Transmission 1
  (1965); English reprint, International Journal of Computer Mathematics
  2 (1968),
  \href{https://doi.org/10.1080/00207166808803030}{doi:10.1080/00207166808803030}.
  Full (page numbers from the 1968 reprint).
\item
  {[}T29{]} G. J. Chaitin, ``On the length of programs for computing
  finite binary sequences'', J. ACM 13 (1966).
  \href{https://doi.org/10.1145/321356.321363}{doi:10.1145/321356.321363}.
  Unavailable.
\item
  {[}T30{]} C. H. Bennett, ``Logical Reversibility of Computation'', IBM
  J. Res. Dev. 17(6), 525-532 (1973).
  \href{https://doi.org/10.1147/rd.176.0525}{doi:10.1147/rd.176.0525}.
  Full.
\item
  {[}T31{]} G. McKee, ``A Few (Proposed) Fundamental Laws of
  Computation'', IEEE Computer 28(1), 120 (1995), Open Channel. No DOI
  located;
  \href{https://vldb.org/dblp/db/journals/computer/computer28.html}{DBLP
  table of contents}. Unavailable.
\item
  {[}T32{]} G. Delétang et al., ``Language Modeling Is Compression'',
  ICLR 2024. \href{https://arxiv.org/abs/2309.10668}{arXiv:2309.10668}.
  Full.
\item
  {[}T33{]} J. Rissanen, ``Modeling by shortest data description'',
  Automatica 14 (1978).
  \href{https://doi.org/10.1016/0005-1098(78)90005-5}{doi:10.1016/0005-1098(78)90005-5}.
  Unavailable.
\item
  {[}T34{]} J. Schmidhuber, ``Driven by Compression Progress'', 2009.
  \href{https://arxiv.org/abs/0812.4360}{arXiv:0812.4360}. Full.
\item
  {[}T35{]} U. Alon et al., ``code2vec: Learning Distributed
  Representations of Code'', POPL 2019.
  \href{https://arxiv.org/abs/1803.09473}{arXiv:1803.09473}. Full.
\item
  {[}T36{]} Y. Du and I. Mordatch, ``Implicit Generation and Modeling
  with Energy-Based Models'', NeurIPS 2019.
  \href{https://arxiv.org/abs/1903.08689}{arXiv:1903.08689}. Full.
\item
  {[}T37{]} F. Chollet, ``On the Measure of Intelligence'', 2019.
  \href{https://arxiv.org/abs/1911.01547}{arXiv:1911.01547}. Full.
\item
  {[}T38{]} E. Fredkin and T. Toffoli, ``Conservative logic'', Int. J.
  Theor. Phys. 21 (1982).
  \href{https://doi.org/10.1007/BF01857727}{doi:10.1007/BF01857727}.
  Unavailable.
\item
  {[}T39{]} C. Elliott and P. Hudak, ``Functional reactive animation'',
  ICFP 1997.
  \href{https://doi.org/10.1145/258948.258973}{doi:10.1145/258948.258973}.
  Unavailable.
\item
  {[}T40{]} D. Ha and J. Schmidhuber, ``World Models'', 2018.
  \href{https://arxiv.org/abs/1803.10122}{arXiv:1803.10122}. Full.
\item
  {[}T41{]} M. Raissi, P. Perdikaris and G. E. Karniadakis, ``Physics
  Informed Deep Learning (Part I)'', 2017.
  \href{https://arxiv.org/abs/1711.10561}{arXiv:1711.10561}. Full.
\item
  {[}T42{]} S. Greydanus, M. Dzamba and J. Yosinski, ``Hamiltonian
  Neural Networks'', NeurIPS 2019.
  \href{https://arxiv.org/abs/1906.01563}{arXiv:1906.01563}. Full.
\item
  {[}T43{]} J. J. Hopfield, ``Neural networks and physical systems with
  emergent collective computational abilities'', PNAS 79 (1982).
  \href{https://doi.org/10.1073/pnas.79.8.2554}{doi:10.1073/pnas.79.8.2554}.
  Unavailable.
\item
  {[}T44{]} G. E. Hinton, ``Training products of experts by minimizing
  contrastive divergence'', Neural Computation 14 (2002).
  \href{https://doi.org/10.1162/089976602760128018}{doi:10.1162/089976602760128018}.
  Unavailable.
\end{enumerate}

\section{Convergence by commitment}\label{convergence-by-commitment}

The RCE thesis rests on seven commitments. Each has been reached
independently, often more than once, by people working for their own
reasons. The tables below list projects whose primary sources (paper,
official documentation or repository) were opened and checked against
the sentence written here. Each row says what the project established
and where it stops relative to the commitment. These are convergence
observations: agreement in direction, not confirmation of RCE's claims
and not a ranking. Items from OpenIE, the organization behind RCE, are
marked \emph{same house} and are never used as evidence.

Acronyms used in this section: RCE, the Reality Code Engine; AST,
abstract syntax tree; IR, intermediate representation; LLM, large
language model; JSON, JavaScript Object Notation; WASM, WebAssembly;
SMT, satisfiability modulo theories (Z3 is an SMT solver); DST,
deterministic simulation testing; CEGIS, counterexample-guided inductive
synthesis.


\begin{center}
\openiefig{fig/rce-jepa-fig08-convergence.png}{0.94}{https://research.openie.dev/figures/rce-jepa/convergence.svg}
\end{center}


\emph{Figure 8. Which commitment each verified outside project reaches,
according to its own primary source. A dot is a convergence observation,
not a score; RCE is not drawn. Select the figure to enlarge it. .}

\subsection{1. One canonical form; surface text is a
view}\label{one-canonical-form-surface-text-is-a-view}

{\footnotesize\setlength{\tabcolsep}{3.8pt}\renewcommand{\arraystretch}{1.15}
\begin{longtable}[]{@{}
  >{\raggedright\arraybackslash}p{(\linewidth - 4\tabcolsep) * \real{0.2000}}
  >{\raggedright\arraybackslash}p{(\linewidth - 4\tabcolsep) * \real{0.4000}}
  >{\raggedright\arraybackslash}p{(\linewidth - 4\tabcolsep) * \real{0.4000}}@{}}
\toprule\noalign{}
\begin{minipage}[b]{\linewidth}\raggedright
\textbf{Project}
\end{minipage} & \begin{minipage}[b]{\linewidth}\raggedright
\textbf{What it established}
\end{minipage} & \begin{minipage}[b]{\linewidth}\raggedright
\textbf{Where it stops}
\end{minipage} \\
\midrule\noalign{}
\endhead
\bottomrule\noalign{}
\endlastfoot
Intentional Programming, Simonyi, 1995 {[}C1{]} & A program is a tree of
nodes, each an instance of an intention; users edit the tree and view it
in any notation; computational intent can be kept separate from
implementation & A research system at Microsoft; the report describes no
determinism or replay model (a transaction file feeds version control,
but no replay semantics are given) \\
Unison {[}C2{]} & Each definition is identified by a hash of its syntax
tree; names are separately stored metadata & The hash is of structure:
two definitions with the same behaviour and different structure have
different hashes \\
Equality saturation and egg {[}C3; C4{]} & An e-graph holds many
equivalent versions of a program at once; a cost model picks one (Tate
et al.: a pseudo-Boolean solver over a static cost model; egg: an
extractor with a user-provided cost function) & An optimization and
synthesis technique applied per function (intraprocedural in Tate et
al.) or over terms in a user-defined language (egg, which has no
built-in support for binding); saturation is bounded when it does not
terminate, and the e-graph is not kept as the program's stored form \\
Tacit, 2026 {[}C5{]} & The AST is the source of truth; every valid AST
has one byte-exact canonical text; definitions are addressed by BLAKE3
hash; malformed code becomes typed Hole nodes & One target (LLVM); no
state-machine semantics stated \\
X07, 2026 {[}C6{]} & Canonical source is a JSON AST, edited with JSON
Patch; a 646-code diagnostic catalogue with enforced quickfix coverage &
Canonical bytes, not canonical behaviour \\
Codifide, 2026 {[}C7{]} & The canonical form is a typed hypergraph with
a deterministic byte form and SHA-256 hash; surface syntax is a
projection; intent is part of a definition's identity & Described by its
authors as a contract-and-dispatch language, not general computation \\
NAIL {[}C8{]} & \texttt{nail\ canonicalize} normalizes any program;
\texttt{nail\ check\ -\allowbreak{}\/\allowbreak{}-\allowbreak{}strict} rejects non-canonical input & An RFC
8785-inspired JSON subset that, in its own words, does not claim full
RFC 8785 compliance; byte canonicalization only \\
EqNet, Allamanis et al., 2017 {[}C9{]} & Neural representations trained
on known equivalence classes so that equivalent expressions get
identical or highly similar vectors & Boolean and polynomial expressions
(supervised), not programs \\
\end{longtable}
}

\subsection{2. A program is a deterministic state machine over a
log}\label{a-program-is-a-deterministic-state-machine-over-a-log}

{\footnotesize\setlength{\tabcolsep}{3.8pt}\renewcommand{\arraystretch}{1.15}
\begin{longtable}[]{@{}
  >{\raggedright\arraybackslash}p{(\linewidth - 4\tabcolsep) * \real{0.2409}}
  >{\raggedright\arraybackslash}p{(\linewidth - 4\tabcolsep) * \real{0.3796}}
  >{\raggedright\arraybackslash}p{(\linewidth - 4\tabcolsep) * \real{0.3796}}@{}}
\toprule\noalign{}
\begin{minipage}[b]{\linewidth}\raggedright
\textbf{Project}
\end{minipage} & \begin{minipage}[b]{\linewidth}\raggedright
\textbf{What it established}
\end{minipage} & \begin{minipage}[b]{\linewidth}\raggedright
\textbf{Where it stops}
\end{minipage} \\
\midrule\noalign{}
\endhead
\bottomrule\noalign{}
\endlastfoot
Abstract State Machines, Gurevich, 2000 {[}C10{]} & The paper's title
states the result: sequential abstract state machines capture sequential
algorithms & A theorem, not a system: every sequential algorithm,
defined by three postulates (sequential time, abstract state, bounded
exploration; Def. 6.1), has an equivalent sequential abstract state
machine that matches it step for step (Main Theorem 6.13). That the
postulates cover all sequential algorithms is argued, not proved \\
Event sourcing, Fowler, 2005 {[}C11{]} & ``Capture all changes to an
application state as a sequence of events'' & A pattern, published as a
draft; replay stays consistent only if gateways record external query
responses and suppress external updates during replay \\
The Elm Architecture {[}C12{]} & Model, View, Update: state, a view of
the state, and an update from messages & One language and one target
(the browser) \\
Dedalus and the CALM theorem, Hellerstein, Alvaro et al.~{[}C13; C14{]}
& Dedalus: Datalog with explicit time, modelling mutable state and
asynchronous messaging. CALM: a program has a consistent,
coordination-free distributed implementation if and only if it is
monotonic (proved by Ameloot et al.; {[}C14{]} gives a proof sketch) &
Dedalus and Bloom are research languages and CALM is a theorem; neither
emits to other languages \\
Durable execution: Temporal, Restate {[}C15; C16{]} & Temporal: workflow
code ``must be deterministic to support replay''. Restate: every step
and its result are recorded in a journal & Replay serves resilience, not
program identity \\
TigerBeetle deterministic simulation testing, 2026 {[}C17{]} & Logical
determinism (like FoundationDB) plus physical determinism: replicas
converge to byte-by-byte identical state; real code runs in a simulator
& One database, built by hand to this discipline \\
FoundationDB simulation testing {[}C18{]} & ``A deterministic simulation
of an entire FoundationDB cluster within a single-threaded process'';
determinism ``allows perfect repeatability of a simulated run'' & A
testing discipline for one database, not a program form \\
musil {[}C19{]} & An explicit-state model checker as a Python library;
its simulate mode runs the real node code on a seeded, replayable
simulated network & Bounded state spaces; it checks the model, and
simulation finds bugs rather than proving their absence \\
AILANG; agent-scroll {[}C20; C21{]} & AILANG: every construct has
deterministic semantics, with replayable execution and structured
traces. agent-scroll: canonical, byte-deterministic, hash-chained
transcripts with a verifier that catches any single-byte change & AILANG
is functional, not a state-machine form; agent-scroll logs
conversations, not programs \\
\end{longtable}
}

\subsection{3. Every target is a projection: write once, emit
many}\label{every-target-is-a-projection-write-once-emit-many}

{\footnotesize\setlength{\tabcolsep}{3.8pt}\renewcommand{\arraystretch}{1.15}
\begin{longtable}[]{@{}
  >{\raggedright\arraybackslash}p{(\linewidth - 4\tabcolsep) * \real{0.1938}}
  >{\raggedright\arraybackslash}p{(\linewidth - 4\tabcolsep) * \real{0.4031}}
  >{\raggedright\arraybackslash}p{(\linewidth - 4\tabcolsep) * \real{0.4031}}@{}}
\toprule\noalign{}
\begin{minipage}[b]{\linewidth}\raggedright
\textbf{Project}
\end{minipage} & \begin{minipage}[b]{\linewidth}\raggedright
\textbf{What it established}
\end{minipage} & \begin{minipage}[b]{\linewidth}\raggedright
\textbf{Where it stops}
\end{minipage} \\
\midrule\noalign{}
\endhead
\bottomrule\noalign{}
\endlastfoot
MLIR {[}C22{]} & A reusable compiler infrastructure spanning levels of
abstraction, hardware targets and domains & Starts below the source: it
has no parser generation or AST construction (Sec. 7), and its paper
names translation validation across an extensible dialect ecosystem as
an open problem (Sec. 2) \\
Model Driven Architecture, OMG, 2003 {[}C23{]} &
Computation-independent, platform-independent and platform-specific
viewpoints & Petre reports, from Forward and Lethbridge, that keeping
the model up to date with the code was the biggest perceived problem
(68\% agreement) {[}C24{]} \\
Protocol Buffers; Lumos {[}C25; C26{]} & Define a schema once and
generate bindings in many languages (Lumos: Rust and TypeScript for
Solana) & Data shapes, not behaviour; the same Protocol Buffers message
can have many binary serializations (``You cannot compare two messages
for equality without fully parsing them'') \\
Forge, 2026 {[}C27{]} & ``Many frontends feed one verified intermediate
representation, and that representation feeds many backends'' & Board
and peripheral descriptions, not general programs \\
Faber Romanus, 2026 {[}C28{]} & One high-level IR projected to Rust, Go,
TypeScript, Swift and Faber, and through a lower IR to LLVM and WASM
text, with Metal and CUDA kernels; support is measured and stated per
target & Agreement across targets is measured by corpus coverage, not
proved \\
Intent, 2026 {[}C29{]} & One source file compiles to Rust, JavaScript
and WASM; contracts are enforced at runtime in every target and
optionally checked by Z3 & Three targets \\
NanoLang {[}C30{]} & Core semantics proved in Coq with no axioms; C is
the production target & PTX and RISC-V are experimental; LLVM and WASM
are planned \\
\end{longtable}
}

\subsection{4. Store intent; derive the
implementation}\label{store-intent-derive-the-implementation}

{\footnotesize\setlength{\tabcolsep}{3.8pt}\renewcommand{\arraystretch}{1.15}
\begin{longtable}[]{@{}
  >{\raggedright\arraybackslash}p{(\linewidth - 4\tabcolsep) * \real{0.2925}}
  >{\raggedright\arraybackslash}p{(\linewidth - 4\tabcolsep) * \real{0.3537}}
  >{\raggedright\arraybackslash}p{(\linewidth - 4\tabcolsep) * \real{0.3537}}@{}}
\toprule\noalign{}
\begin{minipage}[b]{\linewidth}\raggedright
\textbf{Project}
\end{minipage} & \begin{minipage}[b]{\linewidth}\raggedright
\textbf{What it established}
\end{minipage} & \begin{minipage}[b]{\linewidth}\raggedright
\textbf{Where it stops}
\end{minipage} \\
\midrule\noalign{}
\endhead
\bottomrule\noalign{}
\endlastfoot
FORTRAN, Backus et al., 1957; deductive synthesis, Manna and Waldinger,
1980 {[}C31; C32{]} & Automatic coding from formulas; programs extracted
from constructive proofs & FORTRAN became a compiler; deductive
synthesis had no implementation at the time \\
Spec-driven development, as classified by Böckeler {[}C33{]} & Three
levels: spec-first, spec-anchored and spec-as-source, in which ``only
the spec is edited by the human, the human never touches the code''; of
Kiro, spec-kit and Tessl, only Tessl explores spec-as-source & The spec
is natural language, and Böckeler saw code generated repeatedly from one
spec differ; nothing checks that the code is the spec \\
Vericoding benchmark, 2025 {[}C34{]} & 12,504 formal specifications in
Dafny, Verus and Lean; at least one of nine off-the-shelf LLMs, with
five attempts each, solved 82\% (Dafny), 44\% (Verus) and 27\% (Lean) of
the tasks, and the best single model reached 67.5\%, 30.9\% and 17.9\%
(Table 3); adding natural-language descriptions gave no statistically
significant gain on the one subset tested (Verina, 157 tasks; Sec. 4.1)
& Specification languages are the proof languages; no runtime form or
multi-target emission; among successes about 9\% of specs were too weak
to capture intent (Sec. 4) \\
Codifide {[}C7{]} & Every definition carries an intent that compilation
never discards & Contract-and-dispatch scope \\
\end{longtable}
}

\textbf{Recommendation.} The vericoding benchmark {[}C34{]} is used as
an outside benchmark for the step from specification to Form: a slice of
its state-machine-shaped Dafny, Verus and Lean specifications is lifted
into Forms and the share the RCE gate accepts is reported beside the
published vericoding results.

\subsection{5. The model proposes; a deterministic gate
decides}\label{the-model-proposes-a-deterministic-gate-decides}

{\footnotesize\setlength{\tabcolsep}{3.8pt}\renewcommand{\arraystretch}{1.15}
\begin{longtable}[]{@{}
  >{\raggedright\arraybackslash}p{(\linewidth - 4\tabcolsep) * \real{0.2973}}
  >{\raggedright\arraybackslash}p{(\linewidth - 4\tabcolsep) * \real{0.3514}}
  >{\raggedright\arraybackslash}p{(\linewidth - 4\tabcolsep) * \real{0.3514}}@{}}
\toprule\noalign{}
\begin{minipage}[b]{\linewidth}\raggedright
\textbf{Project}
\end{minipage} & \begin{minipage}[b]{\linewidth}\raggedright
\textbf{What it established}
\end{minipage} & \begin{minipage}[b]{\linewidth}\raggedright
\textbf{Where it stops}
\end{minipage} \\
\midrule\noalign{}
\endhead
\bottomrule\noalign{}
\endlastfoot
CompCert; seL4 {[}C35; C36{]} & Machine-checked compiler and kernel;
CompCert checks an unverified heuristic's output with a verified
checker; seL4's proof cost about 20 person-years & Proof effort per
artifact \\
Sketching, Solar-Lezama et al., 2006 (the counterexample-driven loop
later named CEGIS) {[}C37{]} & Against a human-written partial program
(sketch) and an executable specification, a synthesizer proposes, a
checker returns a counterexample, and the loop repeats & Proposers are
SAT solvers, not learned models; limited to finite programs (bounded
input, guaranteed termination), with guarantees only for the bounded
instance \\
AlphaProof; Clover; AutoVerus {[}C38; C39; C40{]} & A model proposes;
Lean, Dafny or Verus checks & Gate on proofs or code text; no canonical
program form \\
Agentic model checking, 2026; LLM-driven verified synthesis for concept
design, 2026 {[}C41; C42{]} & ``Agents propose, solvers verify'' (an LLM
still triages counterexamples for realism); in the second, invariants
alone under-constrained designs, which differed across runs, some of
them implausible & {[}C41{]} verifies existing code with LLM-inferred
specs; checks are bounded (default unwind 4) and spec-relative, and the
authors name over-weak specs that ``cause bugs to be missed silently''
as the primary failure mode. {[}C42{]} synthesizes coordination rules
for design models (not code), checked only within a bounded Alloy scope,
with an LLM translating the design into Alloy. Both warn that validity
is not intent \\
Vera, 2026; Vow; Thermite {[}C43; C44; C45{]} & Vera: mandatory
contracts, proved by the Z3 SMT solver where decidable and otherwise
checked at run time: ``The model doesn't need to be right, it needs to
be checkable''. Vow: contracts statically verified by ESBMC bounded
model checking. Thermite: contracts with counterexamples on failure and
a recorded assurance level (Verus/Z3, Lean, Kani bounded, runtime) &
Gate at the function; no replay or promotion \\
Decision models, a class of concept {[}C46{]} & A model whose entire
output is a typed answer (one option from a closed set, a bounded score,
or a yes/no probability) with a probability for every option, trained
for calibration & Needs a closed option set. The RCE project builds its
own decision head and uses no outside decision model \\
\end{longtable}
}

\subsection{6. Accepted work is promoted and
reused}\label{accepted-work-is-promoted-and-reused}

{\footnotesize\setlength{\tabcolsep}{3.8pt}\renewcommand{\arraystretch}{1.15}
\begin{longtable}[]{@{}
  >{\raggedright\arraybackslash}p{(\linewidth - 4\tabcolsep) * \real{0.2828}}
  >{\raggedright\arraybackslash}p{(\linewidth - 4\tabcolsep) * \real{0.3586}}
  >{\raggedright\arraybackslash}p{(\linewidth - 4\tabcolsep) * \real{0.3586}}@{}}
\toprule\noalign{}
\begin{minipage}[b]{\linewidth}\raggedright
\textbf{Project}
\end{minipage} & \begin{minipage}[b]{\linewidth}\raggedright
\textbf{What it established}
\end{minipage} & \begin{minipage}[b]{\linewidth}\raggedright
\textbf{Where it stops}
\end{minipage} \\
\midrule\noalign{}
\endhead
\bottomrule\noalign{}
\endlastfoot
Gödel machine, Schmidhuber; Löbian obstacle, Yudkowsky and Herreshoff
{[}C47; C48{]} & Self-rewrite only after a proof, under its own axioms,
that the rewrite has higher utility than continuing to search; in the
straightforward formalism, a proof-based agent cannot trust a successor
that uses its own proof system (Löb's theorem), so trust shrinks with
each rewrite, and the paper's technical workarounds are only partial &
Theory \\
Voyager; STOP; Darwin Gödel Machine; AlphaEvolve {[}C49; C50; C51;
C52{]} & A growing library of executable skills; scaffolds that improve
themselves (the model's weights are unchanged); archives of
self-modified agents (kept if they still compile and can edit code;
parents chosen by benchmark score); code evolved under automatic
evaluators & Selection by score or evaluator, not by proof at an
interface. Voyager's evaluator is a GPT-4 critic; AlphaEvolve's scores
may include LLM-graded feedback, and its correctness came from
construction, randomized testing or human and verification review
outside the loop \\
Getafix; Sorald; TBar {[}C53; C54; C55{]} & Fix patterns mined from past
fixes or written as templates, reused on new cases without a model; with
perfect fault localization TBar correctly fixes 74 of 395 Defects4J bugs
(plausible patches for 101), and 43 with standard fault localization &
Patterns are per bug class, learned or written for one language
(Getafix's method itself is language-agnostic); Getafix checks only its
top fix, by re-running the static analyzer \\
\end{longtable}
}

\subsection{7. Energy or cost is the
objective}\label{energy-or-cost-is-the-objective}

{\footnotesize\setlength{\tabcolsep}{3.8pt}\renewcommand{\arraystretch}{1.15}
\begin{longtable}[]{@{}
  >{\raggedright\arraybackslash}p{(\linewidth - 4\tabcolsep) * \real{0.2348}}
  >{\raggedright\arraybackslash}p{(\linewidth - 4\tabcolsep) * \real{0.4000}}
  >{\raggedright\arraybackslash}p{(\linewidth - 4\tabcolsep) * \real{0.3652}}@{}}
\toprule\noalign{}
\begin{minipage}[b]{\linewidth}\raggedright
\textbf{Project}
\end{minipage} & \begin{minipage}[b]{\linewidth}\raggedright
\textbf{What it established}
\end{minipage} & \begin{minipage}[b]{\linewidth}\raggedright
\textbf{Where it stops}
\end{minipage} \\
\midrule\noalign{}
\endhead
\bottomrule\noalign{}
\endlastfoot
Landauer, 1961; Bennett, 1973 {[}C56; C57{]} & Erasing a bit has a
thermodynamic floor; computation can in principle be made logically
reversible & Physics, not a software objective \\
FrugalGPT, 2023 {[}C58{]} & A cascade of models matched the best single
model with up to 98\% lower cost on its tasks & Dollar cost; no
deterministic tiers below the models \\
Cost-of-pass, 2025; AI Agents That Matter, 2024 {[}C59; C60{]} &
Expected monetary cost of a correct solution; report and jointly
optimize cost with accuracy & Dollars and tokens, not joules; the
Cost-of-Pass authors note energy could replace dollars but do not
measure it \\
Compiler world models, 2026 {[}C61{]} & Rolls a tensor program's state
forward through scheduling actions in latent space, then ranks candidate
schedules within a workload before measuring on hardware (a relative
score, not a latency prediction) & Tensor schedules only (TVM
AutoScheduler); ranks candidates, does not generate them \\
Codifide {[}C7{]} & Selects among satisfied candidates by declared cost
& Declared cost, not measured energy \\
JouleOS (OpenIE) {[}same house{]} & Energy budgets and the
\texttt{est\_j}, \texttt{reported\_\allowbreak{}j}, \texttt{measured\_\allowbreak{}j} labels &
Same house: listed for disclosure, not as evidence \\
\end{longtable}
}

\subsection{How four commitments meet the later
lanes}\label{how-four-commitments-meet-the-later-lanes}

Section 3 below records three lanes. Figure 8 is unchanged. The notes
here are convergence readings of sources already in the tables, plus the
lane reports, not a score.

\begin{itemize}
\tightlist
\item
  \textbf{One canonical form.} Unison {[}C2{]} identifies a definition
  by a hash of its syntax tree, and names are a view. egg {[}C4{]}
  stores equivalent programs as one e-graph. Both treat structure, not
  surface text, as the object.
\item
  \textbf{The model proposes and a deterministic gate decides.}
  AlphaEvolve {[}C52{]} and AlphaProof {[}C38{]} are hybrid instances,
  together with the vericoding benchmark {[}C34{]}. Section 3 adds the
  Lean provers DeepSeek-Prover-V2, Kimina-Prover Preview and
  Goedel-Prover, where the Lean kernel is the gate, and type-constrained
  decoding, where the type system acts while the model samples.
\item
  \textbf{Promotion and reuse.} DreamCoder (Ellis et al., Section 2
  below) learns a library by compression. LILO (Grand et al., Section 3
  below) pairs a large language model with symbolic refactoring. The
  library is the reuse step in those systems. The domains they report
  are small.
\item
  \textbf{Cost.} In the CWM report, a best of 16 samples at test time
  moves SWE-bench Verified \href{mailto:pass@1}{\nolinkurl{pass@1}} from
  53.9\% to 65.8\% {[}B23{]}. That is a compute cost on one benchmark,
  not a joule measurement. Mercury reports throughput of 1,109 and 737
  tokens per second on an H100 (Inception Labs, Section 3 below).
  Throughput is not joules. This page still counts cost as electrical
  energy, idle subtracted when a meter reading exists, and does not
  convert those figures into joules.
\end{itemize}

\subsection{What the RCE paper states (OpenIE's own
work)}\label{what-the-rce-paper-states-openies-own-work}

The following are from
\href{https://research.openie.dev/papers/rce/}{the RCE paper} (revision
df38742, 6 November 2026). They are listed here as the project's own
work, not as a comparison with the projects above.

\begin{itemize}
\tightlist
\item
  \textbf{Canonical form.} The canonical serialization is injective on
  Forms, so two Forms that differ have different shape hashes.
  Equivalence of behaviour is defined separately (fit equivalence), and
  in the engine's design equivalent variants are grouped by sound
  rewrites in an e-graph; the paper does not claim that equivalent
  behaviour always yields one canonical form.
\item
  \textbf{Proofs.} Determinism, totality, replay as a fold of the log,
  and atomicity of a kept step (an event that does not apply keeps the
  state) are proved in Lean 4 on a core calculus. The Rust program
  emitter is proved on a core as a verified translation into MiniRust, a
  deep embedding of the emitted Rust subset; that the Rust compiler
  implements MiniRust is trusted and tested. This proof covers the
  program emitter, not the Form projection.
\item
  \textbf{Gate.} In the engine's design, a candidate passes replay of
  recorded inputs, a shadow run with a coverage bound, an output schema
  check and a Kani proof certificate, and a proposed rewrite also needs
  a Z3 proof. In the public record of one run, a local
  1.5-billion-parameter model made 18 proposals, 13 of them distinct;
  the gate did not admit 15, and the 3 it admitted are one rewrite under
  three names that repeats a rule the engine already had.
\item
  \textbf{Energy.} In the engine's design every receipt carries
  \texttt{est\_j}; an unpriced model call is never read as zero.
  \texttt{measured\_\allowbreak{}j} (wall meter) is empty in every result so far.
\end{itemize}

\subsection{Sources for this section}\label{sources-for-this-section}

Labels as in the ledger: Full = the whole primary source (paper,
documentation page, README or project description) was read, including
references and appendices where present. Every source listed here was
read in full on or before 14 October 2026; {[}C46{]} points to sources
listed in the decision-head section.

\begin{enumerate}
\def\labelenumi{\arabic{enumi}.}
\tightlist
\item
  {[}C1{]} C. Simonyi, ``The Death of Computer Languages, the Birth of
  Intentional Programming'', Microsoft Research MSR-TR-95-52, September
  1995.
  \href{https://www.microsoft.com/en-us/research/publication/the-death-of-computer-languages-the-birth-of-intentional-programming/}{Microsoft
  Research};
  \href{https://uni.r-brause.de/seminare/SemSS07/simonyi95death.pdf}{PDF
  copy}. Full.
\item
  {[}C2{]} Unison, ``The big idea''.
  \href{https://www.unison-lang.org/docs/the-big-idea/}{unison-lang.org}.
  Full.
\item
  {[}C3{]} R. Tate, M. Stepp, Z. Tatlock and S. Lerner, ``Equality
  Saturation: A New Approach to Optimization'', Logical Methods in
  Computer Science 7(1:10), 2011 (extended version of POPL 2009).
  \href{https://arxiv.org/abs/1012.1802}{arXiv:1012.1802}. Full.
\item
  {[}C4{]} M. Willsey, C. Nandi, Y. R. Wang, O. Flatt, Z. Tatlock and P.
  Panchekha, ``egg: Fast and Extensible Equality Saturation'', POPL
  2021. \href{https://arxiv.org/abs/2004.03082}{arXiv:2004.03082}. Full.
\item
  {[}C5{]} Tacit.
  \href{https://github.com/weetster/tacit}{github.com/weetster/tacit}.
  Full (README).
\item
  {[}C6{]} X07.
  \href{https://github.com/x07lang/x07}{github.com/x07lang/x07}. Full
  (README).
\item
  {[}C7{]} Codifide.
  \href{https://github.com/codifide/codifide-programming-language}{github.com/codifide/codifide-programming-language}.
  Full (README).
\item
  {[}C8{]} NAIL.
  \href{https://pypi.org/project/nail-lang/}{pypi.org/project/nail-lang}.
  Full (project description).
\item
  {[}C9{]} M. Allamanis, P. Chanthirasegaran, P. Kohli and C. Sutton,
  ``Learning Continuous Semantic Representations of Symbolic
  Expressions'', ICML 2017 (PMLR 70).
  \href{https://arxiv.org/abs/1611.01423}{arXiv:1611.01423}. Full.
\item
  {[}C10{]} Y. Gurevich, ``Sequential Abstract State Machines Capture
  Sequential Algorithms'', ACM TOCL 1(1), 2000.
  \href{https://doi.org/10.1145/343369.343384}{doi:10.1145/343369.343384};
  \href{https://www.cs.tau.ac.il/~nachumd/models/gurevich.pdf}{PDF
  copy}. Full.
\item
  {[}C11{]} M. Fowler, ``Event Sourcing'', 2005.
  \href{https://martinfowler.com/eaaDev/EventSourcing.html}{martinfowler.com}.
  Full.
\item
  {[}C12{]} The Elm Architecture.
  \href{https://guide.elm-lang.org/architecture/}{guide.elm-lang.org}.
  Full.
\item
  {[}C13{]} P. Alvaro, W. R. Marczak, N. Conway, J. M. Hellerstein, D.
  Maier and R. Sears, ``Dedalus: Datalog in Time and Space'', UC
  Berkeley EECS-2009-173, December 2009.
  \href{https://www2.eecs.berkeley.edu/Pubs/TechRpts/2009/EECS-2009-173.pdf}{PDF}.
  Full.
\item
  {[}C14{]} J. M. Hellerstein and P. Alvaro, ``Keeping CALM: When
  Distributed Consistency is Easy'', 2019.
  \href{https://arxiv.org/abs/1901.01930}{arXiv:1901.01930}. Full. The
  formal proof is T. J. Ameloot, F. Neven and J. Van den Bussche,
  ``Relational transducers for declarative networking'', JACM 60(2),
  2013.
\item
  {[}C15{]} Temporal, ``Workflow Definition''.
  \href{https://docs.temporal.io/workflow-definition}{docs.temporal.io}.
  Full.
\item
  {[}C16{]} Restate, ``Key Concepts: Durable Execution''.
  \href{https://docs.restate.dev/foundations/key-concepts\#durable-execution}{docs.restate.dev}.
  Full.
\item
  {[}C17{]} C. Bhandari, ``Protocol-Aware Deterministic Simulation
  Testing'', TigerBeetle, 20 August 2026.
  \href{https://tigerbeetle.com/blog/2026-08-20-protocol-aware-dst/}{tigerbeetle.com}.
  Full.
\item
  {[}C18{]} FoundationDB, ``Simulation and Testing''.
  \href{https://apple.github.io/foundationdb/testing.html}{apple.github.io/foundationdb}.
  Full.
\item
  {[}C19{]} musil (project description, version 0.35.6).
  \href{https://pypi.org/project/musil/}{pypi.org/project/musil}. Full.
\item
  {[}C20{]} AILANG.
  \href{https://github.com/sunholo-data/ailang}{github.com/sunholo-data/ailang}.
  Full (README).
\item
  {[}C21{]} agent-scroll.
  \href{https://github.com/p-vbordei/agent-scroll}{github.com/p-vbordei/agent-scroll}.
  Full (README).
\item
  {[}C22{]} C. Lattner et al., ``MLIR: A Compiler Infrastructure for the
  End of Moore's Law'', 2020.
  \href{https://arxiv.org/abs/2002.11054}{arXiv:2002.11054}. Full.
\item
  {[}C23{]} OMG, MDA Guide Version 1.0, 2003.
  \href{https://www.omg.org/mda/mda_files/MDA_Guide_Version1-0.pdf}{PDF}.
  Full.
\item
  {[}C24{]} M. Petre, ``UML in practice'', ICSE 2013.
  \href{https://oro.open.ac.uk/35805/}{author copy}. Full.
\item
  {[}C25{]} Protocol Buffers overview.
  \href{https://protobuf.dev/overview/}{protobuf.dev}. Full.
\item
  {[}C26{]} Lumos introduction.
  \href{https://docs.lumos-lang.org/getting-started/introduction}{docs.lumos-lang.org}.
  Full.
\item
  {[}C27{]} Forge.
  \href{https://github.com/Sanjith-Shan/Forge}{github.com/Sanjith-Shan/Forge}.
  Full (README).
\item
  {[}C28{]} Faber Romanus.
  \href{https://github.com/faberlang/faber}{github.com/faberlang/faber}.
  Full (README).
\item
  {[}C29{]} Intent.
  \href{https://github.com/lhaig/intent}{github.com/lhaig/intent}. Full
  (README).
\item
  {[}C30{]} NanoLang.
  \href{https://github.com/jordanhubbard/nanolang}{github.com/jordanhubbard/nanolang}.
  Full (README).
\item
  {[}C31{]} J. W. Backus et al., ``The FORTRAN Automatic Coding
  System'', 1957.
  \href{https://archive.computerhistory.org/resources/text/Fortran/102663113.05.01.acc.pdf}{Computer
  History Museum copy}. Full.
\item
  {[}C32{]} Z. Manna and R. Waldinger, ``A Deductive Approach to Program
  Synthesis'', ACM TOPLAS 2(1), 1980.
  \href{https://doi.org/10.1145/357084.357090}{doi:10.1145/357084.357090}.
  Full.
\item
  {[}C33{]} B. Böckeler, ``Understanding Spec-Driven-Development: Kiro,
  spec-kit, and Tessl'', martinfowler.com, 15 October 2025.
  \href{https://martinfowler.com/articles/exploring-gen-ai/sdd-3-tools.html}{martinfowler.com}.
  Full.
\item
  {[}C34{]} S. Bursuc, T. Ehrenborg, S. Lin et al., ``A benchmark for
  vericoding: formally verified program synthesis'', 2025.
  \href{https://arxiv.org/abs/2509.22908}{arXiv:2509.22908}. Full.
\item
  {[}C35{]} X. Leroy, ``Formal verification of a realistic compiler'',
  CACM, 2009.
  \href{https://xavierleroy.org/publi/compcert-CACM.pdf}{PDF}. Full.
\item
  {[}C36{]} G. Klein et al., ``seL4: Formal Verification of an OS
  Kernel'', SOSP 2009.
  \href{https://www.sigops.org/s/conferences/sosp/2009/papers/klein-sosp09.pdf}{PDF}.
  Full.
\item
  {[}C37{]} A. Solar-Lezama, L. Tancau, R. Bodik, V. Saraswat and S.
  Seshia, ``Combinatorial Sketching for Finite Programs'', ASPLOS 2006.
  \href{https://people.csail.mit.edu/asolar/papers/asplos06-final.pdf}{PDF}.
  Full.
\item
  {[}C38{]} T. Hubert et al., ``Olympiad-level formal mathematical
  reasoning with reinforcement learning'', Nature, 2025.
  \href{https://www.nature.com/articles/s41586-025-09833-y}{Nature}.
  Full.
\item
  {[}C39{]} C. Sun et al., ``Clover: Closed-Loop Verifiable Code
  Generation''.
  \href{https://arxiv.org/abs/2310.17807}{arXiv:2310.17807}. Full.
\item
  {[}C40{]} C. Yang et al., ``AutoVerus''.
  \href{https://arxiv.org/abs/2409.13082}{arXiv:2409.13082}. Full.
\item
  {[}C41{]} Y. Sun, J. Liu, D. Kroening and J. Xue, ``Agentic Model
  Checking'', 2026.
  \href{https://arxiv.org/abs/2605.21434}{arXiv:2605.21434}. Full.
\item
  {[}C42{]} A. Cunha, ``Verified LLM-Driven Synthesis for Concept
  Design'', 2026.
  \href{https://arxiv.org/abs/2607.15718}{arXiv:2607.15718}. Full.
\item
  {[}C43{]} Vera.
  \href{https://github.com/aallan/vera}{github.com/aallan/vera}. Full
  (README).
\item
  {[}C44{]} Vow.
  \href{https://github.com/vow-lang/vow}{github.com/vow-lang/vow}. Full
  (README).
\item
  {[}C45{]} Thermite.
  \href{https://github.com/dollspace-gay/Thermite}{github.com/dollspace-gay/Thermite}.
  Full (README).
\item
  {[}C46{]} Decision models as a class: see the decision-head section
  above and its sources {[}B7; B8; B12{]}.
\item
  {[}C47{]} J. Schmidhuber, ``Goedel Machines: Self-Referential
  Universal Problem Solvers Making Provably Optimal Self-Improvements'',
  TR IDSIA-19-03 (v5, 2006).
  \href{https://arxiv.org/abs/cs/0309048}{arXiv:cs/0309048}. Full.
\item
  {[}C48{]} E. Yudkowsky and M. Herreshoff, ``Tiling Agents for
  Self-Modifying AI, and the Löbian Obstacle'', MIRI early draft, 7
  October 2013.
  \href{https://intelligence.org/files/TilingAgentsDraft.pdf}{PDF}.
  Full.
\item
  {[}C49{]} G. Wang et al., ``Voyager: An Open-Ended Embodied Agent with
  Large Language Models'', 2023.
  \href{https://arxiv.org/abs/2305.16291}{arXiv:2305.16291}. Full.
\item
  {[}C50{]} E. Zelikman, E. Lorch, L. Mackey and A. T. Kalai,
  ``Self-Taught Optimizer (STOP): Recursively Self-Improving Code
  Generation'', COLM 2024.
  \href{https://arxiv.org/abs/2310.02304}{arXiv:2310.02304}. Full.
\item
  {[}C51{]} J. Zhang, S. Hu, C. Lu, R. Lange and J. Clune, ``Darwin
  Gödel Machine: Open-Ended Evolution of Self-Improving Agents'', ICLR
  2026. \href{https://arxiv.org/abs/2505.22954}{arXiv:2505.22954}. Full.
\item
  {[}C52{]} A. Novikov et al., ``AlphaEvolve: A coding agent for
  scientific and algorithmic discovery'', Google DeepMind white paper,
  2025. \href{https://arxiv.org/abs/2506.13131}{arXiv:2506.13131}. Full.
\item
  {[}C53{]} J. Bader, A. Scott, M. Pradel and S. Chandra, ``Getafix:
  Learning to Fix Bugs Automatically'', OOPSLA 2019.
  \href{https://arxiv.org/abs/1902.06111}{arXiv:1902.06111}. Full.
\item
  {[}C54{]} K. Etemadi et al., ``Sorald: Automatic Patch Suggestions for
  SonarQube Static Analysis Violations'', 2021 (v2 2022).
  \href{https://arxiv.org/abs/2103.12033}{arXiv:2103.12033}. Full.
\item
  {[}C55{]} K. Liu et al., ``TBar: Revisiting Template-based Automated
  Program Repair'', 2019.
  \href{https://arxiv.org/abs/1903.08409}{arXiv:1903.08409}. Full.
\item
  {[}C56{]} R. Landauer, ``Irreversibility and heat generation in the
  computing process'', 1961.
  \href{https://doi.org/10.1147/rd.53.0183}{doi:10.1147/rd.53.0183}.
  Full.
\item
  {[}C57{]} C. H. Bennett, ``Logical Reversibility of Computation'',
  1973.
  \href{https://doi.org/10.1147/rd.176.0525}{doi:10.1147/rd.176.0525}.
  Full.
\item
  {[}C58{]} L. Chen, M. Zaharia and J. Zou, ``FrugalGPT'', 2023.
  \href{https://arxiv.org/abs/2305.05176}{arXiv:2305.05176}. Full.
\item
  {[}C59{]} M. H. Erol, B. El, M. Suzgun, M. Yüksekgönül and J. Zou,
  ``Cost-of-Pass: An Economic Framework for Evaluating Language
  Models'', ICLR 2026.
  \href{https://arxiv.org/abs/2504.13359}{arXiv:2504.13359}. Full.
\item
  {[}C60{]} S. Kapoor, B. Stroebl, Z. S. Siegel, N. Nadgir and A.
  Narayanan, ``AI Agents That Matter'', 2024.
  \href{https://arxiv.org/abs/2407.01502}{arXiv:2407.01502}. Full.
\item
  {[}C61{]} H. Pan, L. Huang, X. Zhou, M. Xing and Y. Wu, ``Toward
  Compiler World Models: Learning Latent Dynamics for Efficient Tensor
  Program Search'', 2026.
  \href{https://arxiv.org/abs/2606.09312}{arXiv:2606.09312}. Full.
\end{enumerate}

\textbf{Not carried over from the survey}, because the primary source
was not obtained or did not support the claim: Smalltalk image claims,
Darklang hashing, Pantograph and Grove, Codong, Lume, Magpie, Codex
(Tedrow), Spec IR, Boruna, d-raft, Sockudo, DBOS, P language compilation
to C, UNCOL and ANDF, Thrift and OpenAPI, Pact, Aver, Prove, reqlan,
translation validation (Pnueli), SyGuS, Futamura and Turchin, PAR,
DéjàVu and Gabel and Su, Simon and Kadane, and the comparative
scoreboard.

\section{1. History: seventy years of ``say what you want, get the
program''}\label{history-seventy-years-of-say-what-you-want-get-the-program}

The wish behind this topic is old. What has changed is which parts of it
are now cheap. The timeline below lists the main attempts; the text
after it says what each one achieved and where it stopped.


\begin{center}
\openiefig{fig/rce-jepa-fig09-timeline.png}{0.94}{https://research.openie.dev/figures/rce-jepa/timeline.svg}
\end{center}


\emph{Figure 9. Attempts to turn a statement of intent into a working
program. CASE means computer-aided software engineering; UML means
Unified Modeling Language; Z, B and VDM (Vienna Development Method) are
mathematical specification methods; Coq is a proof assistant. Green:
machine-checked systems. Amber: learned models. Blue: the question this
topic studies. .}

\begin{itemize}
\tightlist
\item
  \textbf{Automatic programming, 1950s.} FORTRAN (FORmula TRANslation)
  was presented as ``The FORTRAN Automatic Coding System''. Its authors
  wrote that its purpose was to let the IBM 704 ``code problems for
  itself and produce as good programs as human coders (but without the
  errors)'' (Backus et al., 1957, p.~188; PDF p.~2). Example: a
  scientist wrote a formula; the system wrote the machine code. Today
  that is called a compiler, not automatic programming.
\item
  \textbf{Deductive synthesis, 1960s-80s.} Simon's heuristic compiler
  (1963) treated programming as heuristic problem solving: a means-end
  search, built on the General Problem Solver (GPS), that writes
  routines from descriptions of the state before and after, by
  ``selective trial-and-error search'' (Simon, RAND RM-3588-PR, 1963,
  Summary, p.~vii). Manna and Waldinger (1980) treated a program as
  something extracted from a constructive proof that a specification can
  be met: ``From such a proof, a program to compute z can be extracted''
  (Manna and Waldinger, STAN-CS-78-690, 1978, Specification section; the
  same sentence appears in the TOPLAS article, 1980, p.~92); both
  versions note that ``no implementation exists'' (report, Motivation,
  p.~2; TOPLAS, p.~91). The JACM article states Simon's conclusion
  directly: ``compiling tasks--at least simple compiling tasks--can
  fruitfully be viewed as problem-solving tasks'' (Simon, 1963, Summary,
  p.~505), and it makes no practical claim: ``no claims will be made for
  the heuristic programming procedures described here as practical
  approaches to the construction of compilers'' (Introduction, p.~493).
  Manna and Waldinger observed that theorem proving had been ``abandoned
  as an approach to program synthesis when the development of
  sufficiently powerful automatic theorem provers appeared to flounder''
  (TOPLAS, 1980, p.~120). All four texts were read in full.
\item
  \textbf{Program derivation, 1975-86.} Dijkstra introduced guarded
  commands ``as a building block'' for programs and presented ``the
  formal derivation of programs'' from conditions on the final state
  (Dijkstra, EWD 472, abstract). Bird's calculus of lists (the
  Bird-Meertens formalism) is ``used to derive efficient solutions'' by
  rewriting with algebraic identities (Bird, 1986, PDF p.~3). Example:
  for Gries's problem ``to compute the minimum of the sums of all
  segments'' of a list, the direct program takes O(n\^{}3) steps, and
  rewriting with the calculus yields ``a linear time algorithm'' (Sec.
  4.2, PDF p.~34).
\item
  \textbf{Fourth-generation languages, CASE tools and expert systems,
  1980s.} Report generators, form builders and diagram tools promised
  programs from screens and pictures. Brooks reviewed graphical
  programming, automatic programming and expert systems in 1986 and
  argued that none of them ``promises even one order-of-magnitude
  improvement'' (Brooks, 1986, PDF pp.~1, 8-11). Kemerer reported,
  citing an earlier study, that ``one year after introduction, 70
  percent of CASE tools and techniques are never used'' and that ``the
  first project written with an integrated CASE tool typically fails to
  deliver improved results''; he attributes part of the abandonment to a
  learning curve that dips below the old technology before it rises
  (Kemerer, 1992, pp.~23-24, Fig. 1).
\item
  \textbf{Japan's Fifth Generation Computer Systems project, 1982-92.} A
  national effort built on logic programming in the Prolog family; ICOT
  spent ``about 54 billion Yen, that is some 450 million US\$'' over ten
  years (van de Riet, 1993, Sec. 1, p.~83). One goal in the 1981 plan
  was ``to lessen the burden of software generation in order that a high
  level requirements specification is sufficient for automatic
  processing'' (Sec. 3, p.~85). The parallel inference machine met its
  100 MLIPS target, but van de Riet scores ``break-through in
  architecture'' and ``break-through in software'' as failures and
  judges technical impact ``small'' (Sec. 6, p.~103); on the effort
  needed to program the system, ``there is almost no evidence'' (Sec.
  4.1, p.~86).
\item
  \textbf{Refinement and model-driven methods, 1980s-2000s.} Z, B and
  VDM refine a mathematical specification step by step toward code. The
  Model Driven Architecture guide of the Object Management Group
  ``specifies three viewpoints on a system, a computation independent
  viewpoint, a platform independent viewpoint and a platform specific
  viewpoint'' (OMG, 2003, Sec. 2.2.5, PDF p.~13). Example: draw a class
  diagram once and generate Java and SQL (Structured Query Language)
  from it. Hutchinson et al.~note that MDE ``has been developed largely
  without the support of empirical data'' (2011, Sec. 1, p.~471); among
  250+ practitioners surveyed, 74\% said MDE ``requires them to carry
  out significant extra training'' (Sec. 4.3.1, p.~475), and one
  interviewee reported 4 to 8 times productivity internally while
  telling outsiders ``1.5 to 2'' (Sec. 5.2.6, p.~478). The authors call
  their results ``interesting and suggestive rather than conclusive''
  (Sec. 3.1, p.~473).
\item
  \textbf{Machine-checked systems, 2008-09.} CompCert is a C compiler
  whose ``proof of semantic preservation'' was carried out in the Coq
  proof assistant (Leroy, 2009, abstract, PDF p.~1). seL4 is a
  microkernel of ``8,700 lines of C code'' with a machine-checked
  ``proof of functional correctness'' (abstract, PDF p.~1); the proof
  cost ``about 20 py'' (person-years) (Klein et al., 2009, Sec. 5.2, PDF
  p.~13). Coq can extract a proved function into ordinary code in
  ``functional languages such as Ocaml, Haskell or Scheme'' (Letouzey,
  2008, Sec. 1, PDF p.~1); the extraction step itself was not yet
  mechanically verified, and its correctness rested on paper proofs
  (Sec. 6, PDF p.~10).
\end{itemize}

\section{2. Why the earlier attempts
stalled}\label{why-the-earlier-attempts-stalled}

\begin{itemize}
\tightlist
\item
  \textbf{Essential versus accidental complexity.} Brooks separated
  ``essential tasks, the fashioning of the complex conceptual structures
  that compose the abstract software entity'', from ``accidental tasks,
  the representation of these abstract entities in programming
  languages'' (abstract, PDF p.~1). He quoted Parnas: ``automatic
  programming always has been a euphemism for programming with a
  higher-level language than was presently available to the programmer''
  (PDF p.~9), and wrote that ``even perfect program verification can
  only establish that a program meets its specification'' (PDF p.~11)
  (Brooks, 1986). Example: a better notation for a payroll rule does not
  decide what the rule should be.
\item
  \textbf{Natural language is ambiguous.} Dijkstra argued that formal
  texts are, ``when you come to think of it, an amazingly effective tool
  for ruling out all sorts of nonsense that, when we use our native
  tongues, are almost impossible to avoid'' (Dijkstra, EWD 667, para.
  6). Example: ``sort the users by name and date'' has at least two
  readings.
\item
  \textbf{Specification costs as much as the program.} In seL4 the
  abstract specification took ``about 4 person months'' and the proof
  about 20 person-years, against ``a total cost of 2.2 py including the
  Haskell effort'' to build the kernel itself (Klein et al., 2009, Sec.
  5.2, PDF pp.~12-13). The 20 py includes about 9 py of research and
  reusable proof tooling; the seL4-specific proof ``was 11 py'' (PDF
  p.~13).
\item
  \textbf{Combinatorial search.} Searching all programs grows
  exponentially with program length. DreamCoder addresses this by
  learning a library of reusable pieces and a neural network ``to guide
  the search for programs'' in a ``wake-sleep'' loop (Ellis et al.,
  2021, abstract, PDF p.~1); the paper calls such search ``intractable
  for general-purpose algorithms, because of the combinatorial nature of
  the search space'' (introduction, PDF p.~2); the supplement repeats
  that finding programs ``is difficult because of the infinitely large,
  combinatorial search landscape'' (supplement S4.1, PDF p.~42 of the
  version with supplement), the same kind of search problem that limited
  deductive synthesis.
\item
  \textbf{Brittleness and adoption.} Diagram-to-code tools generated
  code that was then hand-edited, and the model and the code drifted
  apart. Petre reports from earlier surveys that ``the biggest perceived
  problem of model-centric approaches is keeping the model up-to-date
  with the code (68\% agreement)'' (Forward and Lethbridge, as reported
  by Petre, 2013, author copy p.~3); among Petre's own 50 informants, 35
  did not use UML (Unified Modeling Language) and 3 used automated code
  generation (pp.~4-5). Adoption studies of CASE and of model-driven
  engineering are listed in the ledger as not read, so no figures from
  them are quoted.
\item
  \textbf{Opinion.} The common thread is that each approach fixed one
  cost (notation, proof, generation) while another stayed manual:
  writing the specification, searching for the program, or checking the
  result.
\end{itemize}

\section{3. Three lanes toward a World of
Code}\label{three-lanes-toward-a-world-of-code}

The wish in Section 1 is old. What changed by November 2026 is that
several research programmes, working for their own reasons, now put a
learned proposer next to something that can say yes or no. This section
records three lanes. A lane is a convergence observation: agreement in
direction, not a ranking, and not a claim that one programme replaces
another. Every figure below is as reported by the cited source. None of
them is a measurement of the Reality Code Engine except the final row of
Section 3.3.

\subsection{3.1 Training on popular
architectures}\label{training-on-popular-architectures}

This lane keeps a widely used model shape, usually a decoder that
predicts the next token (a word piece), and changes the training data or
the reward.

{\footnotesize\setlength{\tabcolsep}{3.8pt}\renewcommand{\arraystretch}{1.15}
\begin{longtable}[]{@{}
  >{\raggedright\arraybackslash}p{(\linewidth - 4\tabcolsep) * \real{0.2687}}
  >{\raggedright\arraybackslash}p{(\linewidth - 4\tabcolsep) * \real{0.3881}}
  >{\raggedright\arraybackslash}p{(\linewidth - 4\tabcolsep) * \real{0.3433}}@{}}
\toprule\noalign{}
\begin{minipage}[b]{\linewidth}\raggedright
\textbf{Who}
\end{minipage} & \begin{minipage}[b]{\linewidth}\raggedright
\textbf{As reported}
\end{minipage} & \begin{minipage}[b]{\linewidth}\raggedright
\textbf{Unfinished in that source}
\end{minipage} \\
\midrule\noalign{}
\endhead
\bottomrule\noalign{}
\endlastfoot
Meta FAIR CodeGen. CWM, a 32 billion parameter dense decoder {[}B23{]} &
Code is learned better when the model sees execution, not only text.
Mid-training uses Python traces and agent trajectories. SWE-bench
Verified \href{mailto:pass@1}{\nolinkurl{pass@1}} 53.9\% for the base
model, 65.8\% with a best of 16 samples at test time. LiveCodeBench
68.6\%. CruxEval-O 94.3\%. 120 million traced functions. 3 million
ForagerAgent trajectories from 10.2 thousand images. An 8 billion
parameter ablation: traces lift CruxEval-O from 44.6 to 73.9, and
Forager lifts SWE-bench Verified from 18.4 to 22.1. & The authors call
trace prediction a ``first step''. A neural debugger, and uses in
verification, are future work in that report. \\
SWE-RL, Wei et al. & Reinforcement learning on pull-request history. The
reward is similarity to an oracle patch, not semantic equality.
Llama3-SWE-RL-70B is reported at 41.0\% on SWE-bench Verified. & The
reward is not semantic equality. \\
SWE-smith, Yang et al. & Synthetic tasks per repository: 50,000
instances, 128 repositories. SWE-agent-LM-32B at 40.2\%
\href{mailto:pass@1}{\nolinkurl{pass@1}}. & Supervised fine-tuning only
in the paper. Reinforcement learning is an external system named
SkyRL. \\
DiffuCoder, Gong et al.~Mercury, Inception Labs & Non-autoregressive
(masked diffusion) code generation. DiffuCoder is a 7 billion parameter
model on 130 billion tokens; coupled-GRPO is reported at an additional
4.4\% on EvalPlus. Mercury is reported at 1,109 and 737 tokens per
second on an H100. & Quality parity with autoregressive models is not
established in the sources read here. \\
PatchDiff, Wang, Pradel and Liu. METR, Becker et al. & Benchmarks can
overstate success, and field productivity can fall. PatchDiff: 7.8\% of
patches counted correct fail developer tests, and resolve rates are
inflated by about 6.4 points. METR, on early-2025 tools: 16 developers,
246 tasks, measured completion time 19\% higher. A forecast of 24\%
lower, and a post-hoc estimate of 20\% lower, did not match that
measurement. & Both checks are dated to early or mid 2025 tools. \\
SWE-Bench Pro, Scale AI. Public leaderboard fetched 6 November 2026 &
Built to resist contamination (copyleft and private repositories) and to
ask for long changes: average 107.4 lines across 4.1 files. 1,865 tasks,
41 repositories (731 public, 276 private, 858 held out). Public board
that day: Muse Spark 1.1 at 61.50 +/- 3.10; gpt-5.4 (xHigh) at 59.10;
Muse Spark at 55.00; claude-opus-4-6 (thinking) at 51.90; gemini-3.1-pro
at 46.10. The top four of those ran with mini-swe-agent. Launch text on
the same page: GPT-5 at 23.3\% and Opus 4.1 at 23.1\%, against more than
70\% on SWE-bench Verified. Private set: Opus 4.1 at 17.8\%, GPT-5 at
14.9\%. The page prints no update date. & The oracle is still
fail-to-pass and pass-to-pass tests. \\
Multi-SWE-bench, ByteDance Seed & Issue resolving across Java,
TypeScript, JavaScript, Go, Rust, C and C++. 1,632 instances from 2,456
candidates, with 68 annotators. & A November 2026 leaderboard figure was
not retrieved. None is stated on this page. \\
\end{longtable}
}

\textbf{Lane reading.} The training reports above move toward semantics:
execution traces, environment dynamics, and rewards a checker can see.
The test suite is still the oracle, and PatchDiff is evidence that a
test suite is incomplete. SWE-Bench Pro's board is higher than its
launch text and still grades patches by tests. That is the opening the
Laws of Coding name. An exact checker replaces ``tests passed'' with an
accepted Form. Acceptance is part of the world, not a score fitted
afterward.

\subsection{3.2 Creative hybrids}\label{creative-hybrids}

This lane pairs a learned proposer with a gate that already exists: an
evaluator, a proof kernel, a type system, or a set of examples.

{\footnotesize\setlength{\tabcolsep}{3.8pt}\renewcommand{\arraystretch}{1.15}
\begin{longtable}[]{@{}
  >{\raggedright\arraybackslash}p{(\linewidth - 4\tabcolsep) * \real{0.2373}}
  >{\raggedright\arraybackslash}p{(\linewidth - 4\tabcolsep) * \real{0.4407}}
  >{\raggedright\arraybackslash}p{(\linewidth - 4\tabcolsep) * \real{0.3220}}@{}}
\toprule\noalign{}
\begin{minipage}[b]{\linewidth}\raggedright
\textbf{Who}
\end{minipage} & \begin{minipage}[b]{\linewidth}\raggedright
\textbf{As reported}
\end{minipage} & \begin{minipage}[b]{\linewidth}\raggedright
\textbf{Unfinished in that source}
\end{minipage} \\
\midrule\noalign{}
\endhead
\bottomrule\noalign{}
\endlastfoot
AlphaEvolve, Google DeepMind {[}C52{]} & A large language model proposes
code edits. An evolutionary loop wraps it. Automatic evaluators decide.
A 4 by 4 complex matrix multiplication in 48 scalar multiplications.
Also datacenter scheduling, circuit simplification, and a speedup of its
own training, with no further figure used here. & It works where an
evaluator exists. \\
AlphaProof, Google DeepMind {[}C38{]} & Reinforcement learning in the
style of AlphaZero, inside Lean. The formal environment grounds the
reasoning. 80 million autoformalized statements. At the 2024
International Mathematical Olympiad, 3 of 5 non-geometry problems, a
silver-medal result with AlphaGeometry 2, at multi-day compute. &
Compute cost. The report is mathematics, not software. \\
DeepSeek-Prover-V2-671B & Subgoal decomposition by DeepSeek-V3, proofs
of the subgoals stitched into a chain-of-thought cold start, then
reinforcement learning. The Lean 4 kernel is the gate. miniF2F-test pass
ratio 88.9\%. PutnamBench 49 of 658. A new ProverBench of 325 problems.
6 of 15 AIME (American Invitational Mathematics Examination) 2024 and
2025 problems in the formal setting, against 8 for DeepSeek-V3 by an
informal majority vote. & Mathematics, not software. The formal and
informal gap is still nonzero. \\
Kimina-Prover Preview, Moonshot and Numina & From Qwen2.5-72B.
Large-scale reinforcement learning with a formal reasoning pattern in
Lean 4. miniF2F 80.7\% at
\href{mailto:pass@8192}{\nolinkurl{pass@8192}}. Distilled models at 1.5
billion and 7 billion parameters were released. & A preview release. \\
Goedel-Prover, Princeton & Natural-language problems autoformalized to
Lean (1.64 million statements). Expert iteration builds more than
800,000 proofs. Supervised fine-tuning only. miniF2F 57.6\% at
\href{mailto:pass@32}{\nolinkurl{pass@32}}. PutnamBench 7 at
\href{mailto:pass@512}{\nolinkurl{pass@512}}. & No reinforcement
learning in version 1. \\
Vericoding benchmark, Beneficial AI Foundation {[}C34{]} & Generate code
and a proof from a formal specification. 12,504 specifications. Success
82.2\% Dafny, 44.2\% Verus, 26.8\% Lean. The abstract rounds those
shares to 82\%, 44\% and 27\%, which is the rounding in the commitment
section below. The split is 3,029 Dafny, 2,334 Verus and 7,141 Lean.
Dafny verification on that line of work moved from 68\% to 96\% in a
year. Natural-language descriptions add little. & The cost of writing
the specification remains. \\
Type-constrained decoding, Mündler et al. & The type system constrains
token sampling. Authors affiliated with the Swiss Federal Institute of
Technology Zurich and the University of California, Berkeley.
Compilation errors down 75.3\% on HumanEval and 52.1\% on MBPP (the
short Python set named in that report). A syntax-only constraint gives
9.0\% and 4.9\% on those two sets. & TypeScript scale only. Types are
not full acceptance of behaviour. \\
DreamCoder, Ellis et al.~LILO, Grand et al. & Library learning by
compression. LILO adds a large language model and symbolic refactoring
(Stitch). DreamCoder is quoted from the full text in Section 2. LILO is
abstract-level on this pass. & Small domain-specific languages in the
sources read for the lane. \\
egg, Willsey et al.~{[}C4{]} & Equality saturation: a set of equivalent
programs stored as one e-graph. Principles of Programming Languages,
2021. & Someone still has to write the rewrite rules. \\
ARC Prize 2024, Chollet, Knoop, Kamradt and Landers. Akyurek et al. &
Program synthesis guided by deep learning, plus test-time training. The
ARC (Abstraction and Reasoning Corpus) private score on ARC-AGI-1 moved
from 33\% to 55.5\% during 2024. An 85\% target was not met. Akyurek et
al.~update weights at inference on the task's examples: 53.0\% on the
ARC public validation set with an 8 billion parameter model, and 61.9\%
when ensembled with program synthesis. & The 85\% target was not met.
The Akyurek regime is few-shot. \\
ARC Prize 2025 results page & Winners on ARC-AGI-2 private, at a cost
cap of 0.20 dollars per task. NVARC (the winning entry named on that
page) 24.0\%, with synthetic data, test-time training and a tiny
recursive net. ARChitects 16.5\%, a two-dimensional masked diffusion
language model with recursive refinement. MindsAI 12.6\%. Paper awards
on that page: a tiny recursive net of about 7 million parameters at
about 45\% on ARC-AGI-1 and about 8\% on ARC-AGI-2; SOAR (evolutionary
program synthesis with self-fine-tuning, as that page names it) up to
52\% on ARC-AGI-1; CompressARC, a minimum-description-length method with
no pretraining, about 20\% to 34\% on ARC-AGI-1 and about 4\% on
ARC-AGI-2. & The grand prize was unclaimed. ARC-AGI-3, an interactive
track, is in development. \\
Large Concept Models, Meta FAIR & Autoregressive prediction in a
sentence-embedding space, SONAR (the sentence-embedding space named in
that report, 200 languages), instead of tokens. Mean squared error,
diffusion, and quantized variants. 1.6 billion parameter models on about
1.3 trillion tokens. One architecture scaled to 7 billion parameters on
about 2.7 trillion tokens. & In the abstract, a concept is a sentence, a
proxy. No code-specific result is given there. \\
\end{longtable}
}

\textbf{Lane reading.} Each of these hybrids has the shape stated in the
principles above: the model proposes, and a deterministic gate decides.
AlphaEvolve and AlphaProof are clear points of contact. The Lean provers
are the same shape with the kernel as the gate, and the gains they
report came from data generated against that kernel. The ARC Prize 2025
page shows program search, test-time training, diffusion and small
recursive nets on one loop: propose, check against examples, refine. The
shared limit is also shared. They need an evaluator, a kernel, or a set
of examples. The Laws of Coding are the claim that an environment of
that kind exists for coding in general, not only for mathematics and for
benchmark kernels. Large Concept Models are the closest token-free
relative of JEPA in language: prediction in a concept space, still
without an exact acceptance step.

\subsection{3.3 Reconceptions of computation, programming and
coding}\label{reconceptions-of-computation-programming-and-coding}

This lane does not start from a benchmark score. It starts from a
question about what a program is.

{\footnotesize\setlength{\tabcolsep}{3.8pt}\renewcommand{\arraystretch}{1.15}
\begin{longtable}[]{@{}
  >{\raggedright\arraybackslash}p{(\linewidth - 4\tabcolsep) * \real{0.2803}}
  >{\raggedright\arraybackslash}p{(\linewidth - 4\tabcolsep) * \real{0.3939}}
  >{\raggedright\arraybackslash}p{(\linewidth - 4\tabcolsep) * \real{0.3258}}@{}}
\toprule\noalign{}
\begin{minipage}[b]{\linewidth}\raggedright
\textbf{Who}
\end{minipage} & \begin{minipage}[b]{\linewidth}\raggedright
\textbf{As reported}
\end{minipage} & \begin{minipage}[b]{\linewidth}\raggedright
\textbf{Unfinished in that source}
\end{minipage} \\
\midrule\noalign{}
\endhead
\bottomrule\noalign{}
\endlastfoot
HVM2, Higher Order Company & Interaction combinators as a model of
computation, with parallelism inherent in the model. This page cites the
model. & Figures in that paper are not used here. Production adoption is
open. \\
Unison {[}C2{]} & A definition is the hash of its syntax tree. Names are
views over that hash. The language is in production use. & Ecosystem
size. \\
Lean, Dafny and Verus communities, via the vericoding benchmark
{[}C34{]} & A proof is a program the kernel can check. Correctness is a
checkable object. The measured floor used here is the vericoding row
above. & The cost of writing the specification remains. \\
egg {[}C4{]} & Program equivalence classes as data. Programs live in a
quotient. & Completeness of the rewrite rules. \\
Lean FRO (Focused Research Organization), Year 1 roadmap, August 2023 to
July 2024 & Lean as a language and a proof assistant for mathematics and
for software and hardware verification. The deliverables named on that
page include a standard library, a module system, Lake and Reservoir,
proof automation, machine-learning integration and a new code generator.
& Later-year roadmaps are not used on this page. \\
Conal Elliott, denotational design & ``The instance's meaning is the
meaning's instance.'' A semantic model defines correctness and carries
it to the implementation. The note is from 2009, revised 2016. & A
manual design discipline. \\
Radul and Sussman, propagators & Computation as autonomous machines that
add information to shared cells. The account unifies expression-style
programming and constraint programming, including incremental
distributed search. A 2009 technical report, with a Scheme
implementation. & Adoption. \\
Chollet, on the measure of intelligence {[}T37{]} & Intelligence as
skill-acquisition efficiency. ARC is the measure. Program synthesis is
the framing. The 2019 paper is the citation. Later talks on synthesis
done at the moment of the task are not cited on their own. Where the ARC
Prize 2024 report names deep-learning-guided program synthesis, that
report is the citation. & The formal measure is not operationalized
beyond ARC in the abstract read for this lane. \\
This project, the Laws of Coding & A canonical Form is the object.
Surface text is a projection. Cost is counted. Measured results: H1
median rho 9.34 over 185 programs; H2
\texttt{C\_check\ /\ E\_model\ =\ 0.000582\ \textless{}=\ 0.05}
(\texttt{1,718×} cheaper on M5 Max); H5 \& H11 on unseen disjoint Form
shapes via \texttt{FormEBM-\allowbreak{}v2} (AUROC \texttt{0.8922}, \texttt{93.0\%}
5-step monotonicity) and span-normalized \texttt{FormEBM-\allowbreak{}v3}
(\texttt{6.05M} params: AUROC \texttt{0.9437}, \texttt{94.0\%}
\texttt{E\_depth} / \texttt{92.0\%} \texttt{E\_total} monotonicity,
zero-shot OOD H15 pooled AUROC lifted from \texttt{0.2585} to
\texttt{0.5278}); H7, 1,275 mutants and 0 violations on the fuzzed arm;
H8, 10,500 streams, six targets, 0 mismatches; H14 run 3, 111 of 151
(73.5\%) closed at 0 model calls; H15 across the 40 residual tasks,
17/40 to 34/40 single-shot, \texttt{27/40} (\texttt{k\ =\ 16}) for
\texttt{FormJEPA-\allowbreak{}v1} (\texttt{8.42M} params, \texttt{474.\allowbreak{}76\ mJ} RAPL
core \texttt{reported\_\allowbreak{}j}), \texttt{40/40} (\texttt{0} contradictions)
in \texttt{H15-\allowbreak{}Iterative} on \texttt{Opus\ 5.5} and
\texttt{Sonnet\ 5.\allowbreak{}5}, \texttt{40/40} 0-call Gear-0 replay
(\texttt{0.325\ µs}, \texttt{1.77\ µJ}/\allowbreak{}probe), \texttt{16.67×} 3-Gear
episode reduction on \texttt{Zipf(\allowbreak{}0.\allowbreak{}85)} (\texttt{200/200} admitted),
and 151-task \texttt{rce\ repair} rung \texttt{f} audit
(\texttt{3/119\ =\ 2.5\%\ \textgreater{}\ 2.0\%}, \texttt{0}
contradictions, stays in \texttt{observe}). & Full intent-to-Form
synthesis (H3), cross-module signature-and-trace OOD energy calibration
beyond \texttt{0.5278}, H4, H6, H9, H10, H12, H13, the diffusion and
evolutionary arms of H15, and wall-meter \texttt{measured\_\allowbreak{}j} remain
open. \\
\end{longtable}
}

\textbf{Lane reading.} Unison (identity is a hash of structure), egg
(programs as equivalence classes), interaction nets (computation as
local rewriting) and proof assistants (correctness as an object) each
hold one piece of what the Laws of Coding state together: one canonical
Form, exact replay, projection to targets, and checkable acceptance.
Elliott's denotational design and the propagator account are earlier
statements of the same demand: meaning defines the implementation, and
computation adds information rather than only emitting text. They are
independent convergence, each from a different starting point. Chollet's
measure says the thing to count is how efficiently a system acquires a
skill, which is why proposer comparisons track both call budgets and
physical energy (Section 8).

\section{4. Why now: the case for closed-form, deterministic,
AI-proposed
code}\label{why-now-the-case-for-closed-form-deterministic-ai-proposed-code}

\textbf{This section is opinion, built from the history above and from
the three lanes in Section 3. It is a case to be tested, not a result.}

\begin{itemize}
\tightlist
\item
  \textbf{Checkers are cheap and trustworthy where the earlier attempts
  lacked them.} Using Csmith, a random program generator, the authors
  found that CompCert was ``the only compiler we have tested for which
  Csmith cannot find wrong-code errors'' after ``about six CPU-years''
  of trying (Yang et al., 2011, Sec. 3.1, PDF p.~6). Machine-checked
  rules can carry the trust that a generator cannot. CompCert itself
  uses this pattern: its register allocator runs an unverified
  graph-coloring heuristic and checks each result with a verifier proved
  correct in Coq, because graph coloring is ``easier to validate a
  posteriori than to directly prove correct'' (Leroy, 2009, Sec. 4.2,
  PDF p.~7).
\item
  \textbf{Deterministic semantics answer the nondeterminism problem.} If
  the generator is allowed to vary but the accepted artifact is checked
  and replayed exactly, run-to-run variation in the model stops
  mattering for the result.
\item
  \textbf{Verifier-in-the-loop evidence exists.} AlphaCode's behavioral
  filtering, Codex-S's oracle test selection (77.5\% versus 44.5\% for
  log-probability ranking), Clover, AutoVerus and AlphaProof all show
  large gains when a check selects among proposals. Section 3 records
  the same shape in later systems: a model proposes, and an evaluator, a
  proof kernel or a type system decides.
\item
  \textbf{Learned predictors now work in latent space.} I-JEPA and
  V-JEPA 2 show that predicting a summary of the target, rather than
  every detail, can work well for those tasks: V-JEPA 2 had higher
  success on 4 of 5 robot tasks (tied on one) and planned about 15x
  faster than a generative video model in its comparison (Table 3, PDF
  p.~15). In a program, details such as identifier spelling, formatting
  and target language can be filled in deterministically after the
  structure is chosen. Section 5 states the premise and where it stops
  once Forms and the Laws of Coding exist.
\item
  \textbf{Token-by-token generation is expensive.} LLaMA 65B, sharded
  across 8 to 32 V100 GPUs with a 512-token output limit, used ``about
  3-4 Joules for a output token'', counting GPU energy only (Samsi et
  al., 2023, Secs. III-C and IV-C, PDF pp.~4-6). In a study of 88
  models, text-generation models averaged 0.047 kWh (kilowatt-hours) per
  1,000 queries, each query generating only 10 new tokens (Table 2), and
  ``multi-purpose, generative architectures are orders of magnitude more
  expensive than task-specific systems'', partly because a generator
  chooses among ``an entire vocabulary'' (Luccioni et al., 2024,
  abstract, PDF p.~1; Sec. 4.1 and Table 2, PDF p.~6; Sec. 4.2.1, PDF
  p.~8).
\item
  \textbf{Same-host check-to-model comparison (Section 8, H2).} Rather
  than mixing external V100 numbers with local laptop timings, Section 8
  measures both on the same Apple M5 Max host: \texttt{0.\allowbreak{}04904\ J} per
  Form check versus \texttt{84.25\ J} per candidate on local
  \texttt{Qwen\ 3.\allowbreak{}6\ 35B-\allowbreak{}A3B} (\texttt{124.81\ J} per admitted residual
  repair).
\end{itemize}

\subsection{What could still fail}\label{what-could-still-fail}

\begin{itemize}
\tightlist
\item
  \textbf{Expressiveness ceiling (H6).} If real programs need effects or
  nondeterminism the Form cannot state, the ``all code'' premise fails.
\item
  \textbf{Intent may be the bottleneck.} A smaller output space does not
  fix underspecified intent. Signal: H1 and H2 hold while H3 fails.
\item
  \textbf{Spec gaming.} The checker proves a Form is well-formed and
  meets stated properties, not that it matches what the user meant.
  Benchmark tasks carry behavioral traces checked by replay; the
  residual risk is labeled, not solved.
\item
  \textbf{Corpus scarcity.} Signal: H3 flat as the corpus grows.
\item
  \textbf{Collapse.} JEPAs need anti-collapse measures, such as the EMA
  (exponential moving average) target encoder of I-JEPA and V-JEPA 2;
  LeCun (2022) proposes VICReg-style information maximization instead
  (Sec. 4.5). A collapsed Form embedding makes H5 fail. Section 5 states
  this limit beside the others.
\item
  \textbf{Measurement error.} Rootless counters (macmon
  \texttt{IOReport} L2, kernel \texttt{ri\_\allowbreak{}energy\_\allowbreak{}nj} L3, RAPL) cover
  SoC or process counters (\texttt{reported\_\allowbreak{}j}), not
  alternating-current (AC) wall power (\texttt{measured\_\allowbreak{}j}).
\item
  \textbf{Baseline drift.} Co-tenant processes on shared hosts drift L2
  idle baselines (Falsifier 6), requiring per-process-identifier (PID)
  L3 counters or dedicated idle hosts.
\end{itemize}

\section{5. JEPA as premise, and where it
stops}\label{jepa-as-premise-and-where-it-stops}

The premise stays. A JEPA predicts in an embedding space, not in the
space of the raw input. The thesis cites LeCun's 2022 statement of that
idea. V-JEPA 2 reports 62 hours of robot video and zero-shot
pick-and-place success from 65\% to 80\% {[}B1{]}. The
action-conditioned predictor, as the proposer section states from the
same report, was trained on fewer than 62 hours of unlabeled robot
video. LLM-JEPA reports that an embedding-space loss beats a token loss
on a natural-language-to-regular-expression task, on GSM8K (a
grade-school math set) and on Spider (a text-to-database query set), at
1 billion parameter scale (Huang, LeCun and Balestriero, 2025). For
code, this matches the intuition that the useful object is meaning, not
tokens.

Once Forms exist, and once the Laws of Coding fix equivalence and
acceptance, that premise is not the whole method. Six limits follow.
They are reasons to widen the proposer, not reasons to drop the premise.

\begin{enumerate}
\def\labelenumi{\arabic{enumi}.}
\tightlist
\item
  \textbf{An approximate geometry, where an exact one can be computed.}
  JEPA learns an approximate latent geometry. The Laws of Coding give an
  exact one: the canonical Form and the equivalences the checker already
  decides. Learning those by self-supervision spends data on a map that
  can be computed. H10 is the test of whether learning on the quotient
  saves data. H10 has not been run.
\item
  \textbf{Collapse.} Representation collapse, and the need for
  regularizers against it, are open problems in JEPA training. An exact
  canonicalizer cannot collapse. It maps each surface program to one
  Form by a rule, not by a loss.
\item
  \textbf{No native acceptance.} JEPA has no native decoder and no
  acceptance step. Code needs a discrete, checkable output. The energy
  score ranks candidates. It does not decide. The checker decides.
\item
  \textbf{Points, not a quotient.} Alternate Forms are a quotient. Many
  surface programs can be one Form, and one Form can be projected to
  many targets. JEPA embeds points. It does not certify that two points
  are the same Form. An e-graph does. A hash of structure does. In
  Section 8, \texttt{rce\ fmt\ nameless} maps \texttt{7} of \texttt{9}
  alpha-renamed local-variable mutants onto the identical memory-mapped
  Gear-0 (\texttt{DNULLMM1}) slot at zero model calls.
\item
  \textbf{Scale of the evidence.} LLM-JEPA's evidence is at 1 billion to
  2 billion parameter scale, on small datasets. No source read for this
  page shows JEPA training for repository-scale code.
\item
  \textbf{The cost of learned planning.} V-JEPA 2's model predictive
  control (MPC) took about 16 seconds per action {[}B1{]}. Search over
  Forms with a cheap exact checker (\texttt{7.32\ ms},
  \texttt{49.04\ mJ} L3 \texttt{reported\_\allowbreak{}j} on M5 Max) and Gear-0
  lookup (\texttt{0.325\ µs}, \texttt{1.77\ µJ}) bounds repeated rollout
  cost before any learned planner is trained.
\end{enumerate}

\textbf{Consequence.} JEPA stays as the premise: predict meaning, not
text. The proposer slot generalizes to any learned proposer over Forms.
The instances named in Section 3 are autoregressive models with
execution mid-training (CWM), diffusion code models, reinforcement
learning in a formal environment (AlphaProof and the Lean provers),
evolutionary search (AlphaEvolve), library learning (DreamCoder and
LILO), and constrained decoding. Large Concept Models sit beside the
premise in language, not yet in code. The gate is the same for all of
them. Section 6 is the taxonomy. Section 7 says what the world is that
they propose into.

\section{6. The proposer slot is
architecture-neutral}\label{the-proposer-slot-is-architecture-neutral}

Figure 1 draws one pipeline. The learned box is a slot, not a single
architecture. JEPA is the instance in the figure. The gate drawn outside
the box does not change if a different proposer is plugged in. The table
says, for each kind of proposer read in Section 3, what it learns, what
it needs from the Laws of Coding (a canonicalizer, a checker, a cost
meter), the failure mode the sources already name, and which hypothesis
it moves or would move. Section 8 now reports both the first
purpose-trained latent Form JEPA repair proposer (\texttt{FormJEPA-\allowbreak{}v1},
\texttt{8.42M} trainable parameters) and six off-the-shelf
autoregressive, reasoning, and hybrid Gated-DeltaNet / MoE models on the
40 residual tasks.

{\footnotesize\setlength{\tabcolsep}{3.8pt}\renewcommand{\arraystretch}{1.15}
\begin{longtable}[]{@{}
  >{\raggedright\arraybackslash}p{(\linewidth - 8\tabcolsep) * \real{0.1441}}
  >{\raggedright\arraybackslash}p{(\linewidth - 8\tabcolsep) * \real{0.2009}}
  >{\raggedright\arraybackslash}p{(\linewidth - 8\tabcolsep) * \real{0.2009}}
  >{\raggedright\arraybackslash}p{(\linewidth - 8\tabcolsep) * \real{0.2271}}
  >{\raggedright\arraybackslash}p{(\linewidth - 8\tabcolsep) * \real{0.2271}}@{}}
\toprule\noalign{}
\begin{minipage}[b]{\linewidth}\raggedright
\textbf{Proposer}
\end{minipage} & \begin{minipage}[b]{\linewidth}\raggedright
\textbf{What it learns}
\end{minipage} & \begin{minipage}[b]{\linewidth}\raggedright
\textbf{What it needs from the checker}
\end{minipage} & \begin{minipage}[b]{\linewidth}\raggedright
\textbf{Known failure mode}
\end{minipage} & \begin{minipage}[b]{\linewidth}\raggedright
\textbf{Hypothesis it moves / would move}
\end{minipage} \\
\midrule\noalign{}
\endhead
\bottomrule\noalign{}
\endlastfoot
JEPA, an embedding predictor (\texttt{FormJEPA-\allowbreak{}v1} in Section 8) & An
approximate latent geometry of Form structure and edits (EMA target
encoder + VICReg + latent-conditioned typed edit decoder) & A
canonicalizer, a checker, a cost meter, and a discrete edit head. The
energy score does not accept & No native acceptance. Uncalibrated latent
distance alone gives AUROC \texttt{0.618}, and compact non-pretrained
models miss open multi-module API migrations (\texttt{0/9} on H15)
without Gear-2 fallback & \texttt{FormJEPA-\allowbreak{}v1} measured on H15
(\texttt{22/40} at \texttt{k\ =\ 1}, \texttt{27/40} at
\texttt{k\ =\ 16}) and disjoint synthetic defects (\texttt{177/200});
full intent-to-Form H3, calibrated H5/H11, H10, and H13 remain open \\
Autoregressive / reasoning / trace-trained decoders (Claude, Grok,
generative pretrained transformer open-weights (GPT-OSS), DeepSeek, Qwen
3.6, CWM class) & Next-token prediction, chain-of-thought reasoning,
hybrid Gated-DeltaNet / MoE, or execution traces & A checker as
acceptance, plus Gear-0 crystallization so recurring shapes do not
re-invoke the model & Token-space retries spend \texttt{84.25\ J} per
local candidate (\texttt{Qwen\ 3.6}), and long reasoning traces can
exhaust fixed completion caps (\texttt{DeepSeek\ V4.\allowbreak{}1} \texttt{17/40}) &
First single-shot (\texttt{k\ =\ 1}) H15 arms measured in Section 8
(\texttt{17/40} to \texttt{34/40}); RJ-F6 tests whether trace-trained
token models match Form-quotient training at equal energy \\
Diffusion code model & A masked or denoised program, not a left-to-right
token stream & A checker at acceptance. A constraint can also act while
sampling & Quality parity with autoregressive models is not established
in the sources read here & The diffusion arm of H15 (open) \\
Reinforcement learning in a formal environment & Steps that a kernel
accepts, as in AlphaProof and the Lean provers of Section 3 & The kernel
is the environment. For software, that role is the checker & In the
reports read here the domain is mathematics, not software. Compute cost
is part of the AlphaProof result. Goedel-Prover version 1 has no
reinforcement learning & H3 and H12, if the same loop is run on Forms
(open) \\
Evolutionary search & Edits, kept or dropped by an evaluator, as in
AlphaEvolve and in SOAR on the ARC Prize 2025 page & An evaluator. Here
the evaluator is the checker, not a test suite alone & The loop works
where an evaluator exists. Example checks, as on ARC, are not a general
acceptance rule for software & The evolutionary arm of H15, on the
residual of H14 (open) \\
Library learning (DreamCoder, LILO) & A reusable library, by
compression, with or without a large language model in the loop & A
checker on the programs the library emits & The domains in the sources
read for this lane are small & H10 (open) \\
Constrained decoding & Nothing new in the weights. The constraint masks
illegal tokens, as in the type-constrained report & A grammar or a type
system at sampling time, and the checker for behaviour & The headline
result is TypeScript scale. Syntax and types are not full acceptance.
Compilation errors fell 75.3\% on HumanEval and 52.1\% on MBPP, against
9.0\% and 4.9\% for syntax only & H3, only for the share of failures
that are type failures \\
Large Concept Models & The next sentence embedding, not the next token &
An exact acceptance step, which the abstract does not have & A concept
there is a sentence. No code-specific result in the abstract & None of
H1 to H15 until a code result exists \\
\end{longtable}
}

The deterministic ladder in the proposer section runs before any of
these (H14: \texttt{111/151} closed at \texttt{0} calls, leaving
\texttt{40} residual tasks). Section 8 reports both the six-model
off-the-shelf H15 baseline plus Shape D (\texttt{DNULLMM1}) Gear-0
crystallization and the purpose-trained \texttt{FormJEPA-\allowbreak{}v1} arm plus
3-Gear Cascade on those 40 tasks, while keeping the diffusion and
evolutionary arms, intent-to-Form H3, and wall-meter server joule parity
explicitly open.

\section{7. A World of Code}\label{a-world-of-code}

\textbf{Definition.} A World of Code is the state space of Forms, plus
the exact transition and the exact acceptance rule. In the notation of
the thesis, the state space is S, the transition is the function δ, and
acceptance is the checker's verdict on a Form. ``World model'', for
code, means a model of this space. Unlike a physical world model, the
physics is known exactly. Part 4 of the thesis says that for a
JEPA-style predictor. The same physics bounds every proposer in Section
6. CWM learns an interpreter from traces. V-JEPA 2 learns physics from
video. In a Form, the transition is given. There is no irreducible noise
in the next state.

\textbf{Where learning is removed, and where it remains.}
Canonicalization, equivalence and replay are computations. Once the Laws
of Coding fix them, a model does not have to learn them. What remains to
learn is the map from a stated intent to a Form, guidance of the search,
and ranking of whatever the deterministic ladder leaves open
(\texttt{40} real-history tasks after H14 run 3). Once any proposer
solves one of those residual tasks and passes the gate, Gear-0
crystallization turns every recurrence (and \texttt{7/9} alpha-renamings
under \texttt{rce\ fmt\ nameless}) into a \texttt{0.325\ µs}
deterministic lookup.

\textbf{Evaluation.} A test suite is not the Laws of Coding. PatchDiff
reports that 7.8\% of patches counted correct fail developer tests, and
that resolve rates are inflated by about 6.4 points. The METR trial
reports 16 developers, 246 tasks, and a measured completion time 19\%
higher with early-2025 tools. SWE-Bench Pro, on the board fetched 6
November 2026, still uses fail-to-pass and pass-to-pass tests as the
oracle. Multi-SWE-bench widens the set of languages to Java, TypeScript,
JavaScript, Go, Rust, C and C++. Acceptance on this page is the gate
(\texttt{rce\ check} exit \texttt{0} plus obligation preservation),
verified further against historical originals by \texttt{rce\ hash} and
\texttt{rce\ equiv}.

\textbf{Falsifier of the widened claim.} RJ-F6, stated with the others
in Section 8, names the observation that would defeat the Form quotient
as a help to proposers: if a token-space model closes the residual of 40
tasks at equal energy and equal gate acceptance without the Form
quotient. Single-shot frontier token models close up to \texttt{34/40}
(\texttt{85.0\%}) on first encounter, whereas \texttt{FormJEPA-\allowbreak{}v1}
(\texttt{8.42M} parameters over quotient-normalized Form S-expressions)
matches local 35B token-space repair at \texttt{27/40}
(\texttt{k\ =\ 16}) at over an order of magnitude lower
\texttt{reported\_\allowbreak{}j} and \texttt{rce\ fmt\ nameless} supplies
\texttt{7/9} zero-call quotient hits under alpha-renaming; controlled
equal-size quotient-versus-surface sample-efficiency tests (H10) and
intent-to-Form H3 remain open.

\section{8. One concrete instance: the Reality Code
Engine}\label{one-concrete-instance-the-reality-code-engine}

RCE is one closed program form among possible others; the questions
above apply to any language with a small, checkable program shape. In
RCE every program is a Form P = (S, E, δ, D, λ):

\begin{itemize}
\tightlist
\item
  \textbf{S, state:} the typed data the program holds, for example the
  cursor position and text buffer of an editor.
\item
  \textbf{E, events:} the inputs it can receive, for example ``key
  pressed'' or ``file opened''.
\item
  \textbf{δ (delta), transition function:} the rule giving the next
  state from the current state and an event.
\item
  \textbf{D, derivations:} the recorded, hash-chained history of steps,
  which makes every run replayable bit for bit.
\item
  \textbf{λ (lambda), output:} what the program shows or emits for a
  given state, for example the visible lines of text.
\end{itemize}

A checker accepts a Form or names the rule it breaks. At commit
\texttt{1762fb4b} (20 November 2026), the public reference checker
replays an accepted Form exactly, passes \texttt{2,221} of
\texttt{2,236} conformance tests with \texttt{0} failures (\texttt{837}
Rust unit tests), emits Rust, Python, JavaScript and Scala, and is
backed by \texttt{5,537} Lean 4 kernel theorem constants across
\texttt{64} modules (theorems T1 to T63). On Apple M5 Max, dual-channel
\texttt{reported\_\allowbreak{}j} receipts record \texttt{49.04\ mJ} per Form check
(L3 \texttt{ri\_\allowbreak{}energy\_\allowbreak{}nj}, \texttt{7.32\ ms}/\allowbreak{}Form) and
\texttt{15.44\ J} per full 62-Form check, trace-run and five-target emit
suite (\texttt{2.32\ s}/\allowbreak{}pass; \texttt{measured\_\allowbreak{}j\ =\ None},
\texttt{board\_\allowbreak{}metered\ =\ false}). The full description is on the
\href{https://research.openie.dev/about/rce/}{RCE topic page} and at
\href{https://github.com/openIE-dev/rce}{github.com/openIE-dev/rce}.

\textbf{RCE-JEPA} is the working name, kept so the thread stays
findable. The subject is World Coding. JEPA, a joint embedding
predictor, is the premise: predict meaning, not surface text. The
proposer slot is any learned proposer over Forms. Checking, proofs,
derivation, emission and placement stay with RCE. The aim is the best
performance per clock cycle and per joule. \emph{Working hypothesis:} if
all code can be encoded in the Form, all code can be produced by joint
embedding or energy-based models over it. The hypotheses below are
designed so that the evidence can refute this.

\subsection{Defined quantities}\label{defined-quantities}

\begin{itemize}
\tightlist
\item
  \textbf{Output space.} Free-form generation over vocabulary V with
  length L has raw space \textbar V\textbar{}\textsuperscript{L}.
  Ω\textsubscript{n} (omega) is the set of valid Forms of size n (Form
  nodes); Ω\textsubscript{n}\textsuperscript{✓} ⊆ Ω\textsubscript{n} is
  the subset the checker accepts. Reported as
  log\textsubscript{2}\textbar Ω\textsubscript{n}\textbar{} bits,
  estimated by grammar counting and sampling.
\item
  \textbf{Description length.} b(P) is the bits to encode Form P
  canonically; b(code) = tokens ×
  log\textsubscript{2}\textbar V\textbar{} for the equivalent emitted
  source. Ratio ρ (rho) = b(code) / b(P).
\item
  \textbf{Verification cost.} C\textsubscript{check}(P): wall time,
  cycles and joules for the checker to decide P. Unit tests give
  C\textsubscript{test}; a passing test is evidence, not proof.
\item
  \textbf{Acceptance probability.} a = P(checker accepts and task
  specification satisfied \textbar{} intent), measured as
  \href{mailto:pass@1}{\nolinkurl{pass@1}} over the benchmark.
\item
  \textbf{Joules per accepted program.} J\textsubscript{acc} =
  (E\textsubscript{model} + E\textsubscript{check} +
  E\textsubscript{retry}) / N\textsubscript{accepted}, summed over a
  benchmark run (labelled \texttt{reported\_\allowbreak{}j} for L2/L3 counters and
  \texttt{measured\_\allowbreak{}j} only when an external hardware power meter is
  attached).
\item
  \textbf{Cycles per accepted program.} K\textsubscript{acc}, the same
  in CPU (central processing unit) or GPU cycles where counters are
  readable.
\item
  \textbf{AUROC} (area under the receiver operating characteristic
  curve): the probability that a randomly chosen accepted Form gets a
  lower energy than a randomly chosen one that requires change; 0.5 is
  chance, 1.0 is perfect.
\end{itemize}

\textbf{Argument (opinion, to be tested).} A Form has five typed slots
with fixed composition rules. Constraining output to Ω\textsubscript{n}
removes most of \textbar V\textbar{}\textsuperscript{L} before the model
is consulted, so fewer bits are predicted per program (ρ \textgreater{}
1), the checker gives a definite verdict, accepted or requires change,
in place of a probabilistic test, and a requires-change verdict is cheap
and structured, naming the rule, so retries are guided rather than
blind.

\subsection{Falsifiable hypotheses H1 to
H15}\label{falsifiable-hypotheses-h1-to-h15}

All fifteen hypotheses in numerical order. H1 to H7 test the RCE-JEPA
thesis directly, H8 to H10 test the ``Laws of Coding'' thesis above, H11
to H13 test the proposer, H14 measures the deterministic ladder beneath
it, and H15 compares proposer classes on the 40-task residual alongside
Shape D (\texttt{DNULLMM1}) Gear-0 crystallization.

{\footnotesize\setlength{\tabcolsep}{3.8pt}\renewcommand{\arraystretch}{1.15}
\begin{longtable}[]{@{}
  >{\raggedright\arraybackslash}p{(\linewidth - 8\tabcolsep) * \real{0.0357}}
  >{\raggedright\arraybackslash}p{(\linewidth - 8\tabcolsep) * \real{0.3095}}
  >{\raggedright\arraybackslash}p{(\linewidth - 8\tabcolsep) * \real{0.3095}}
  >{\raggedright\arraybackslash}p{(\linewidth - 8\tabcolsep) * \real{0.1667}}
  >{\raggedright\arraybackslash}p{(\linewidth - 8\tabcolsep) * \real{0.1786}}@{}}
\toprule\noalign{}
\begin{minipage}[b]{\linewidth}\raggedright
\textbf{ID}
\end{minipage} & \begin{minipage}[b]{\linewidth}\raggedright
\textbf{Hypothesis}
\end{minipage} & \begin{minipage}[b]{\linewidth}\raggedright
\textbf{Metric}
\end{minipage} & \begin{minipage}[b]{\linewidth}\raggedright
\textbf{Supported if}
\end{minipage} & \begin{minipage}[b]{\linewidth}\raggedright
\textbf{Disproved if}
\end{minipage} \\
\midrule\noalign{}
\endhead
\bottomrule\noalign{}
\endlastfoot
H1 & Form space is much smaller than source space for the same programs.
\textbf{Supported on hand-written corpus} (median ρ = 9.34 over 185
programs) & ρ = b(code)/b(P), median over corpus & ρ ≥ 5 & ρ \textless{}
2 \\
H2 & Checker is cheap relative to generation. \textbf{Supported in
\texttt{reported\_\allowbreak{}j} on Apple M5 Max}:
\texttt{C\_\allowbreak{}check\ /\allowbreak{}\ E\_\allowbreak{}model\ =\ 0.\allowbreak{}04904\ J\ /\allowbreak{}\ 84.\allowbreak{}25\ J\ =\ 0.\allowbreak{}000582}
(\texttt{1,718×} cheaper; \texttt{measured\_\allowbreak{}j\ =\ None}) &
C\textsubscript{check} / E\textsubscript{model} per candidate (joules) &
≤ 0.05 & ≥ 0.5 \\
H3 & Representation-space prediction reaches useful acceptance &
\href{mailto:pass@1}{\nolinkurl{pass@1}} accepted Forms on held-out
intent-to-Form tasks & ≥ baseline
\href{mailto:pass@1}{\nolinkurl{pass@1}} on the same tasks & \textless{}
0.5 × baseline \\
H4 & Energy per accepted program falls & J\textsubscript{acc}(RCE-JEPA)
/ J\textsubscript{acc}(baseline) & ≤ 0.1 & ≥ 1 \\
H5 & Energy ranking concentrates mass on valid Forms. \textbf{AUROC
threshold met on unseen test Form shapes}: contrastive
\texttt{FormEBM-\allowbreak{}v2} (\texttt{5.85M} params) reaches AUROC
\textbf{\texttt{0.8922\ \textgreater{}=\ 0.85}} (\texttt{84.05\%}
solvable top-1 accuracy), and span-normalized decomposed
\texttt{FormEBM-\allowbreak{}v3} (\texttt{6.05M} params) raises unseen-shape AUROC to
\textbf{\texttt{0.9437\ \textgreater{}=\ 0.85}} (\texttt{0.9386} on
\texttt{E\_def} alone) across \texttt{3,064} candidates on \texttt{200}
unseen test beams (vs \texttt{0.618} for uncalibrated
\texttt{FormJEPA-\allowbreak{}v1}), with \texttt{83.44\%} \texttt{E\_def} and
\texttt{80.98\%} \texttt{E\_total} solvable top-1 accuracy
(\texttt{68.0\%} / \texttt{66.0\%} raw across all \texttt{200} beams)
and \texttt{2.23×} / \texttt{2.14×} fewer \texttt{rce\ check} calls;
zero-shot OOD H15 pooled AUROC rises from \texttt{0.2585} to
\texttt{0.5278} (\texttt{0.6018} on large-span valid vs short invalid),
while cross-module signature/trace OOD calibration remains open & top-1
type-check rate; AUROC of E for accepted vs requires change & ≥ 0.8
top-1 valid; AUROC ≥ 0.85 & AUROC ≤ 0.6 \\
H6 & The reference apps are encodable & share of the 6 reference apps
plus RCEdb fully expressed as Forms that pass the checker and replay
bit-identically & 7 of 7 & any app requires an escape hatch outside the
Form \\
H7 & Security is model-independent (fuzzed arm: 0 violations in 1,275
mutants; model arm open) & adversarial Forms (fuzzed and
model-generated) accepted by the checker that violate a stated invariant
& 0 & ≥ 1 confirmed \\
H8 & Determinism: every accepted Form replays to a bit-identical
trajectory and output. \textbf{Supported} (0 mismatches on 10,500
streams across 6 targets) & mismatches over at least 10,000 recorded
event streams on at least 3 different machines and operating systems & 0
& a single mismatch \\
H9 & Cost is predicted before execution & median error of the cost-table
prediction against measured energy (J, named meter) for a stated
machine, over the test families & ≤ 20\% & \textgreater{} 50\% \\
H10 & Learning on the quotient saves data & training examples a
predictor of canonical Forms needs to reach the same test acceptance
rate as one of equal size trained on surface code & at most half (a
saving of 2 times or more) & a saving below 1.25 times \\
H11 & Energy is monotone in edit distance. \textbf{Supported on unseen
5-step cumulative defect chains}: contrastive \texttt{FormEBM-\allowbreak{}v2}
reaches \texttt{93.0\%\ \textgreater{}=\ 80\%} strict 5-step
monotonicity (Spearman \texttt{ρ\ =\ 0.\allowbreak{}9931}), and two-head
\texttt{FormEBM-\allowbreak{}v3} reaches
\textbf{\texttt{94.0\%\ \textgreater{}=\ 80\%}} on \texttt{E\_depth}
(\texttt{94/100}, Spearman \texttt{ρ\ =\ 0.\allowbreak{}9966}) and
\textbf{\texttt{92.0\%\ \textgreater{}=\ 80\%}} on composite
\texttt{E\_total} (\texttt{92/100}, Spearman \texttt{ρ\ =\ 0.\allowbreak{}9846}, vs
\texttt{36.0\%} for uncalibrated \texttt{FormJEPA-\allowbreak{}v1}); unconstrained
same-line random-walk slices reach \texttt{38.0\%} (\texttt{v2}) and
\texttt{32.0\%} / \texttt{18.0\%} (\texttt{v3} \texttt{E\_depth} /
\texttt{E\_total}) & Share of test families in which energy E rises with
edit distance 1 to k from an accepted fit & ≥ 80\% of families &
\textless{} 80\% of families \\
H12 & The decision head is calibrated on gate outcomes & ECE per
decision type, after 10\textsuperscript{4} receipts & ECE ≤ 0.1 and
below a frequency table over shape hashes & ECE \textgreater{} 0.2, or
no better than the frequency table \\
H13 & A discrete latent plans better than a continuous one over long
edit sequences & Planning success at edit-sequence lengths 1, 5 and 20,
discrete arm against continuous SIGReg arm & Discrete arm higher at
length 20 & Neither arm reaches the H3 threshold \\
H14 & A deterministic ladder closes a share of tasks with no model call.
\textbf{Supported by its registered rule}: three registered runs {[}B30,
B31, B32{]}, reported below the table & Share of a registered task set
that the public deterministic ladder (reuse, lookup by type, determined
rewrites, bounded enumeration) closes with zero model calls, every
closure admitted by the checker, at a frozen ladder commit & ≥ 10\% of
the registered tasks, reported with the rung that closed each
(registered threshold, unchanged). Measured: run 1, 118 of 120 (98.3\%),
on mutations the ladder inverts by construction; run 2, 31 of 151
(20.5\%), 27 of them guard negations; run 3, plain arm, 111 of 151
(73.5\%), 110 of them guard negations and renames & \textless{} 10\% of
the registered tasks, the same floor as falsifier RCE-F1 of the RCE
paper. Not met in any run. On the 41 tasks from real history, which the
rule does not single out, runs 2 and 3 each close 1 of 41 (2.4\%) \\
H15 & Proposer classes differ on the 40-task residual. \textbf{Measured
across single-shot proposers {[}B33{]}, \texttt{FormJEPA-\allowbreak{}v1} {[}B34{]},
bounded checker-feedback \texttt{H15-\allowbreak{}Iterative} + \texttt{FormEBM-\allowbreak{}v2}
{[}B35{]}, and span-normalized \texttt{FormEBM-\allowbreak{}v3} {[}B36{]}}:
single-shot models admit \texttt{17/40} to \texttt{34/40},
\texttt{FormJEPA-\allowbreak{}v1} admits \texttt{22/40} (\texttt{k\ =\ 1};
\texttt{25/40} with \texttt{FormEBM-\allowbreak{}v2}, \texttt{24/40} typed /
\texttt{27/40} multi-proposer with \texttt{FormEBM-\allowbreak{}v3} \texttt{E\_def})
and \texttt{27/40} (\texttt{k\ =\ 16}), \texttt{H15-\allowbreak{}Iterative} closes
\textbf{\texttt{40/40} (\texttt{100.0\%})} on \texttt{Opus\ 5.5} and
\texttt{Sonnet\ 5.\allowbreak{}5} with \texttt{0} contradictions (\texttt{39/40} on
\texttt{Grok\ 4.\allowbreak{}20}), warm \texttt{DNULLMM1} Gear-0 replays
\textbf{\texttt{40/40} at \texttt{0} calls}, and the 3-Gear Cascade
admits \texttt{200/200} \texttt{Zipf(\allowbreak{}0.\allowbreak{}85)} arrivals with
\texttt{16.67×} fewer Gear-2 episodes; diffusion, evolutionary, and
wall-meter arms remain open & Gate-accepted share (\texttt{C(\allowbreak{}z)\ =\ 1})
of the 40 residual tasks from the plain arm of H14 run 3, audited by
\texttt{rce\ check}, obligation preservation, \texttt{rce\ hash}, and
\texttt{rce\ equiv} & Registered runs report per-arm shares and local
\texttt{reported\_\allowbreak{}j} {[}B33, B34, B35, B36{]}; diffusion/evolutionary
arms and server wall-meter parity remain open & Classes scored on
different tasks or unverified outputs \\
\end{longtable}
}

\textbf{H14, measured.} The ladder is \texttt{rce\ repair} of
openIE-dev/rce, with the rungs of the proposer section above and an
enumeration bound of 2000 candidates. It has three runs, each registered
before it ran, and together they are the current result. Run 1 {[}B30{]}
closed 118 of 120 tasks with zero model calls, but 119 of its 120 tasks
were operator and integer-literal mutations, which bounded enumeration
inverts by construction, and 9 of the 118 closures pass the examples and
checks without being the original program. Run 2 {[}B31{]} has no
operator or literal mutation: 41 changes from the git history of
openIE-dev/rce, 27 guard negations and 83 renames, 151 tasks in all. It
closed 31 of them (20.5\%) with zero model calls: 1 of the 41 history
tasks; all 27 guard negations, every one by the enumeration edit that
removes a \texttt{not}, which inverts that mutation by construction; and
3 of the 83 renames, none of them restored to the original. Of the 41
history tasks, 39 are programs updated because the language changed, not
defects; only 2 are defects of a program, and the ladder closed one of
them. Without the guard negations, 4 of 124 tasks close (3.2\%). Run 2's
prediction P3, that 10\% to 40\% of the renames close, failed, and its
falsifier F1 was met: the closure rate of run 1 does not carry over to
changes from real history. Run 3 {[}B32{]} ran run 2's 151 tasks
unchanged, at openIE-dev/rce \texttt{8a5d90ea}, with rung 5, renames,
after rung 4; its registration records that renames are the class rung 5
was built for after run 2. Its plain arm, the ladder given only the
task, closed 111 of 151 (73.5\%) with zero model calls, 103 of them
restored to the original: all 83 renames (80 at rung 5, 3 at rung 4), 77
of them restored; all 27 guard negations, each with its outcome in run
2; and 1 of the 41 history tasks, the same one as in run 2. Without the
guard negations, 84 of 124 tasks close (67.7\%). Its known-good arm gave
the ladder each task's original program as well: it closed 116 of 151,
114 restored, including all 83 renames restored and 6 of the 41 history
tasks. Run 3's prediction P6, that in the known-good arm every guard and
history task keeps its plain-arm outcome, failed: 5 more history tasks
close, and the one closed in the plain arm closes at rung 5 instead of
rung 3. None of run 3's six falsifiers was met, and the ladder reported
zero model calls in both arms. H14's registered rule is support at 10\%
or more of the registered tasks, and falsification below 10\%, for the
share the public deterministic ladder closes with zero model calls,
every closure admitted by the checker, at a frozen ladder commit. The
plain arm meets those conditions. By the rule, H14 is supported: the
plain arm of run 3 closes 111 of 151 (73.5\%), leaving 40 open residual
tasks from git history.

\textbf{H2, single-shot H15, and Shape D (\texttt{DNULLMM1}) Gear-0 on
the 40 open residual tasks, measured {[}B33{]}.} Registered prior to
execution (\texttt{rce-\allowbreak{}jepa-\allowbreak{}h15-\allowbreak{}and-\allowbreak{}shape-\allowbreak{}d-\allowbreak{}gear0-\allowbreak{}2026-\allowbreak{}11-\allowbreak{}22},
\texttt{sha256:\ b64c4c67\allowbreak{}1ae85755.\allowbreak{}.\allowbreak{}.\allowbreak{}}) against \texttt{openIE-\allowbreak{}dev/\allowbreak{}rce}
commit \texttt{1762fb4b} (\texttt{2026-\allowbreak{}11-\allowbreak{}20}, binary
\texttt{sha256:\ 1d3b0d07\allowbreak{}ff5b2e8f.\allowbreak{}.\allowbreak{}.\allowbreak{}}), the 40 residual tasks
(\texttt{sha256:\ 3ca3a327\allowbreak{}69145e9b.\allowbreak{}.\allowbreak{}.\allowbreak{}}; \texttt{25} \texttt{W-SYN-1},
\texttt{5} \texttt{W-LEN-1}, \texttt{4} \texttt{W-\allowbreak{}NAME-\allowbreak{}1,\allowbreak{}W-\allowbreak{}TY-\allowbreak{}1},
\texttt{3} \texttt{W-NAME-1}, \texttt{1} \texttt{W-WORLD-1}, \texttt{1}
\texttt{W-\allowbreak{}CAP-\allowbreak{}3,\allowbreak{}W-\allowbreak{}WORLD-\allowbreak{}1}, \texttt{1} \texttt{W-INT-1}) were evaluated
under strict gate predicate \texttt{C(\allowbreak{}z)\ =\ 1}: (1)
\texttt{rce\ check\ candidate.\allowbreak{}rce} exits \texttt{0} (enforcing
\texttt{W-UI-1} and \texttt{W-MUT-1} non-vacuity), (2)
\texttt{no\_\allowbreak{}spec\_\allowbreak{}deletion} preserves all declared \texttt{state},
\texttt{event}, \texttt{error}, \texttt{struct}, \texttt{kind},
\texttt{fn}, \texttt{extern}, and \texttt{step} headers plus every
\texttt{example} left-hand-side trace and \texttt{check} count of the
mutant, and (3) secondary equivalence audit via \texttt{rce\ hash}
(\texttt{restored}), \texttt{rce\ equiv\ fuel\ 2000}
(\texttt{behaviour-\allowbreak{}equal}), and \texttt{rce\ repair} rung \texttt{f}
\texttt{observe} mode.

{\scriptsize\setlength{\tabcolsep}{2.4pt}\renewcommand{\arraystretch}{1.12}
\begin{longtable}[]{@{}
  >{\raggedright\arraybackslash}p{(\linewidth - 14\tabcolsep) * \real{0.1250}}
  >{\raggedright\arraybackslash}p{(\linewidth - 14\tabcolsep) * \real{0.1250}}
  >{\raggedright\arraybackslash}p{(\linewidth - 14\tabcolsep) * \real{0.1250}}
  >{\raggedright\arraybackslash}p{(\linewidth - 14\tabcolsep) * \real{0.1250}}
  >{\raggedright\arraybackslash}p{(\linewidth - 14\tabcolsep) * \real{0.1250}}
  >{\raggedright\arraybackslash}p{(\linewidth - 14\tabcolsep) * \real{0.1250}}
  >{\raggedright\arraybackslash}p{(\linewidth - 14\tabcolsep) * \real{0.1250}}
  >{\raggedright\arraybackslash}p{(\linewidth - 14\tabcolsep) * \real{0.1250}}@{}}
\toprule\noalign{}
\begin{minipage}[b]{\linewidth}\raggedright
Arm / Architecture Family
\end{minipage} & \begin{minipage}[b]{\linewidth}\raggedright
Model
\end{minipage} & \begin{minipage}[b]{\linewidth}\raggedright
Calls
\end{minipage} & \begin{minipage}[b]{\linewidth}\raggedright
Admitted (\texttt{C(z)=1})
\end{minipage} & \begin{minipage}[b]{\linewidth}\raggedright
Restored + Beh-Eq
\end{minipage} & \begin{minipage}[b]{\linewidth}\raggedright
In / Out Tokens
\end{minipage} & \begin{minipage}[b]{\linewidth}\raggedright
Cost (\$)
\end{minipage} & \begin{minipage}[b]{\linewidth}\raggedright
On-Chip \texttt{reported\_\allowbreak{}j}
\end{minipage} \\
\midrule\noalign{}
\endhead
\bottomrule\noalign{}
\endlastfoot
\textbf{Arm 0}: Shape D Cold Static Gear-0 & \texttt{DNULLMM1} (0
records) + \texttt{a..e} & 0 & 0 / 40 (0.0\%) & 0 & 0 / 0 & \$0.0000 &
\textless{} 0.001 J (37.7 µs) \\
\textbf{Arm 1a}: Cloud Autoregressive + Adaptive Thinking &
\texttt{claude-\allowbreak{}opus-\allowbreak{}5-\allowbreak{}5} & 40 & \textbf{34 / 40 (85.0\%)} & 33 (30 + 3)
& 114,855 / 32,159 & \$1.3783 & \texttt{not\_\allowbreak{}available} (N) \\
\textbf{Arm 1b}: Cloud Autoregressive + Adaptive Thinking &
\texttt{claude-\allowbreak{}sonnet-\allowbreak{}5-\allowbreak{}5} & 40 & \textbf{32 / 40 (80.0\%)} & 32 (30 +
2) & 114,855 / 30,764 & \$0.8060 & \texttt{not\_\allowbreak{}available} (N) \\
\textbf{Arm 2c}: Cloud Reasoning Transformer &
\texttt{xai:grok-\allowbreak{}4.\allowbreak{}20-\allowbreak{}0309-\allowbreak{}reasoning} & 40 & \textbf{31 / 40 (77.5\%)} &
30 (28 + 2) & 76,546 / 18,474 & \$0.2747 & \texttt{not\_\allowbreak{}available}
(N) \\
\textbf{Arm 2a}: Cloud Reasoning MoE (Cerebras) &
\texttt{cerebras:gpt-\allowbreak{}oss-\allowbreak{}120b} & 40 & \textbf{29 / 40 (72.5\%)} & 27 (25
+ 2) & 78,825 / 43,596 & \$0.1340 & \texttt{not\_\allowbreak{}available} (N) \\
\textbf{Arm 3}: Local Hybrid Gated-DeltaNet + Sparse MoE &
\texttt{Qwen3.\allowbreak{}6-\allowbreak{}35B-\allowbreak{}A3B} (quantized \texttt{Q5\_K\_M}) & 40 & \textbf{27
/ 40 (67.5\%)} & 26 (25 + 1) & 82,330 / 18,411 & \$0.0000 &
\textbf{3,369.89 J} system-under-test (\texttt{124.\allowbreak{}81\ J/\allowbreak{}adm}) \\
\textbf{Arm 2b}: Cloud multi-head latent attention (MLA) Sparse MoE
(Fireworks) & \texttt{deepseek-\allowbreak{}v4p1-\allowbreak{}flash} & 40 & \textbf{17 / 40
(42.5\%)} & 16 (15 + 1) & 63,446 / 64,215 & \$0.0930 &
\texttt{not\_\allowbreak{}available} (N) \\
\textbf{Arm 5a}: Purpose-Trained Latent Form JEPA (\texttt{k\ =\ 1})
{[}B34{]} & \texttt{FormJEPA-\allowbreak{}v1} (\texttt{8.42M} params, greedy) & 40 &
\textbf{22 / 40 (55.0\%)} & 22 (20 + 2) & Form S-expr & \$0.0000 &
\textbf{58.65 mJ} core / \textbf{1.19 J} pkg RAPL/prop \\
\textbf{Arm 5b}: Purpose-Trained Latent Form JEPA (\texttt{k\ =\ 16})
{[}B34{]} & \texttt{FormJEPA-\allowbreak{}v1} (\texttt{8.42M} params, latent beam) &
40 & \textbf{27 / 40 (67.5\%)} & 27 (25 + 2) & Form S-expr & \$0.0000 &
\textbf{474.76 mJ} core / \textbf{6.74 J} pkg RAPL/prop \\
\textbf{Arm 6}: 3-Gear Cascade v1 (\texttt{t\ =\ 1} Cold Gear-0 →
\texttt{FormJEPA-\allowbreak{}v1} → Single-Shot \texttt{Opus\ 5.5}) {[}B34{]} &
\texttt{DNULLMM1} + \texttt{FormJEPA-\allowbreak{}v1} + \texttt{Opus\ 5.5} & 13 G2 &
\textbf{34 / 40 (85.0\%)} & 33 & Form + 13 cloud & \textbf{\$0.8108} &
\texttt{FormJEPA} RAPL + 13 cloud (N) \\
\textbf{Arm 4a}: Shape D Warm Gear-0 v1 (\texttt{t\ =\ 2} Replay)
{[}B33{]} & \texttt{DNULLMM1} (34 crystallized) & \textbf{0} &
\textbf{34 / 40 (85.0\%)} & 33 & 0 / 0 & \textbf{\$0.0000} &
\textbf{0.325 µs / 1.77 µJ} probe + \texttt{rce\ check} \\
\textbf{Arm 7a}: Bounded Checker-Feedback + Stdlib Headers
(\texttt{\textless{}=\ 5} turns) {[}B35{]} & \texttt{claude-\allowbreak{}opus-\allowbreak{}5-\allowbreak{}5}
(\texttt{H15-\allowbreak{}Iterative}) & 42 & \textbf{40 / 40 (100.0\%)} & 39 (34 + 5)
& 133,527 / 45,711 & \$1.8104 & \texttt{not\_\allowbreak{}available} (N) \\
\textbf{Arm 7b}: Bounded Checker-Feedback + Stdlib Headers
(\texttt{\textless{}=\ 5} turns) {[}B35{]} & \texttt{claude-\allowbreak{}sonnet-\allowbreak{}5-\allowbreak{}5}
(\texttt{H15-\allowbreak{}Iterative}) & 43 & \textbf{40 / 40 (100.0\%)} & 39 (35 + 4)
& 137,385 / 46,668 & \textbf{\$1.1122} & \texttt{not\_\allowbreak{}available} (N) \\
\textbf{Arm 7c}: Bounded Checker-Feedback + Stdlib Headers
(\texttt{\textless{}=\ 5} turns) {[}B35{]} &
\texttt{xai:grok-\allowbreak{}4.\allowbreak{}20-\allowbreak{}0309-\allowbreak{}reasoning} (\texttt{H15-\allowbreak{}Iterative}) & 50 &
\textbf{39 / 40 (97.5\%)} & 37 (30 + 7) & 269,930 / 28,903 & \$4.0555 &
\texttt{not\_\allowbreak{}available} (N) \\
\textbf{Arm 4b}: Shape D Warm Gear-0 v2 (\texttt{t\ =\ 2} Full Replay)
{[}B35{]} & \texttt{DNULLMM1} (40 crystallized) & \textbf{0} &
\textbf{40 / 40 (100.0\%)} & 39 (34 + 5) & 0 / 0 & \textbf{\$0.0000} &
\textbf{0.325 µs / 1.77 µJ} probe + \texttt{rce\ check} \\
\end{longtable}
}

\begin{itemize}
\tightlist
\item
  \textbf{Hypothesis H2 supported (\texttt{1,718×} local check-to-model
  energy advantage):} On Apple M5 Max, per-Form checking costs
  \texttt{C\_\allowbreak{}check\ =\ 0.\allowbreak{}049038\ J} (L3 \texttt{ri\_\allowbreak{}energy\_\allowbreak{}nj},
  \texttt{7.32\ ms}) while one local \texttt{Qwen\ 3.\allowbreak{}6\ 35B-\allowbreak{}A3B}
  candidate generation costs \texttt{E\_\allowbreak{}model\ =\ 84.\allowbreak{}247\ J}
  (\texttt{3,\allowbreak{}303.\allowbreak{}62\ J} L2 GPU + \texttt{66.26\ J} L3
  \texttt{llama-\allowbreak{}server} PID across 40 candidates; \texttt{14,\allowbreak{}476.\allowbreak{}18\ J}
  gross L2 SoC, with \texttt{derived\_\allowbreak{}marginal\_\allowbreak{}j} invalidated by
  Falsifier 6 at \texttt{68.31\%} idle drift), yielding
  \textbf{\texttt{C\_check\ /\ E\_model\ =\ 0.000582\ \textless{}=\ 0.05}}
  (\texttt{measured\_\allowbreak{}j\ =\ None}, \texttt{board\_\allowbreak{}metered\ =\ false}).
\item
  \textbf{Shape D (\texttt{DNULLMM1}) Gear-0 latency, energy, nameless
  invariance, and \texttt{Zipf(\allowbreak{}0.\allowbreak{}85)} amortization:} Over
  \texttt{40,000} portable operating-system interface (POSIX)
  \texttt{mmap} lookups against the frozen 34-record store
  (\texttt{sha256:\ d92a57b4\allowbreak{}0de0b0ec.\allowbreak{}.\allowbreak{}.\allowbreak{}}, verified byte-identical
  pre/post replay), mean probe latency is \textbf{\texttt{0.\allowbreak{}3246\ µs}}
  (\texttt{324.6\ ns}, \texttt{9,516.5} instructions) and
  \textbf{\texttt{1.771\ µJ}} (\texttt{1,\allowbreak{}771.\allowbreak{}3\ nJ} L3
  \texttt{ri\_\allowbreak{}energy\_\allowbreak{}nj}). Keying slots by \texttt{rce\ fmt\ nameless}
  (theorem T7) yields \textbf{\texttt{7\ /\ 9} (\texttt{77.8\%})}
  zero-call hits when local variables are alpha-renamed. On a
  \texttt{200}-arrival \texttt{Zipf(\allowbreak{}s\ =\ 0.\allowbreak{}85)} recurring stream under
  verified-only crystallization, 2-Gear \texttt{DNULLMM1} cuts model
  calls from \texttt{200} to \textbf{\texttt{48} (\texttt{4.17×}
  reduction)} for \texttt{Opus\ 5.5} (\texttt{\$6.81} to
  \texttt{\$2.63}; \texttt{5.41×} / \texttt{37} calls with miss
  deduplication), \texttt{56} (\texttt{3.57×}) for \texttt{Sonnet\ 5.\allowbreak{}5},
  \texttt{57} (\texttt{3.51×}) for \texttt{Grok\ 4.\allowbreak{}20}, \texttt{66}
  (\texttt{3.03×}) for \texttt{GPT-\allowbreak{}OSS-\allowbreak{}120B}, and \texttt{94}
  (\texttt{2.13×} calls, \textbf{\texttt{1.56×} local energy reduction}:
  \texttt{16,\allowbreak{}457.\allowbreak{}9\ J} to \texttt{10,\allowbreak{}553.\allowbreak{}5\ J}) for local
  \texttt{Qwen\ 3.6}.
\item
  \textbf{Scope disclosure on H15:} Two of the 40 historical originals
  predate the \texttt{W-MUT-1} vacuous-example rule
  (\texttt{spec/\allowbreak{}examples/\allowbreak{}57\_\allowbreak{}float\_\allowbreak{}vectors.\allowbreak{}rce@b2b932da} and
  \texttt{programs/\allowbreak{}multi/\allowbreak{}washer.\allowbreak{}rce@2cc82de7}), accounting for part of
  the ceiling near \texttt{34/40}. Diffusion code models, evolutionary
  Form search, and hardware wall-meter \texttt{measured\_\allowbreak{}j} remain open.
\end{itemize}

\textbf{\texttt{FormJEPA-\allowbreak{}v1} latent Form repair proposer, 3-Gear
Cascade, and 151-task \texttt{rce\ repair} rung \texttt{f} promotion
audit, measured {[}B34{]}.} Pre-registered before dataset generation,
weight initialization, and rung-\texttt{f} execution
(\texttt{rce-\allowbreak{}jepa-\allowbreak{}form-\allowbreak{}jepa-\allowbreak{}v1-\allowbreak{}2026-\allowbreak{}11-\allowbreak{}22},
\texttt{sha256:\ 22682e6e\allowbreak{}936ec120.\allowbreak{}.\allowbreak{}.\allowbreak{}};
\texttt{rce-\allowbreak{}repair-\allowbreak{}rung-\allowbreak{}f-\allowbreak{}wave3-\allowbreak{}eval-\allowbreak{}2026-\allowbreak{}11-\allowbreak{}22},
\texttt{sha256:\ 04aff4b2\allowbreak{}b44cdbea.\allowbreak{}.\allowbreak{}.\allowbreak{}}) against \texttt{openIE-\allowbreak{}dev/\allowbreak{}rce}
commit \texttt{1762fb4b} (Linux x86\_64 binary
\texttt{sha256:\ 5ef9931f\allowbreak{}f23a310f.\allowbreak{}.\allowbreak{}.\allowbreak{}}), \texttt{FormJEPA-\allowbreak{}v1}
(\texttt{8,424,073} trainable parameters, \texttt{13,\allowbreak{}917,\allowbreak{}577} with EMA
target buffers; checkpoint \texttt{sha256:\ cbfbf227\allowbreak{}4175c737.\allowbreak{}.\allowbreak{}.\allowbreak{}}) was
trained from random initialization (\texttt{seed\ =\ 20261122}) on
\texttt{6,329} checker-verified defect pairs across \texttt{1,106}
disjoint valid Form shapes (extracted from \texttt{1,387} unique Form
shapes after four-layer path, slug, declaration-name, and
\texttt{458}-digest \texttt{rce\ fmt\ nameless} exclusion with
\texttt{0} overlap against the 151 ladder or 40 residual tasks). It
pairs a 6-layer Pre-LN Transformer context encoder (\texttt{d\ =\ 256})
over quotient-normalized Form AST S-expression tokens with an EMA target
encoder (\texttt{τ\ =\ 0.\allowbreak{}996}), a 3-layer residual MLP latent predictor
trained under Smooth-L\textsubscript{1} plus VICReg variance-covariance
anti-collapse regularization (raising context effective rank from
\texttt{32.90} at epoch 0 to \textbf{\texttt{85.24}} post-training,
coordinate standard deviation \texttt{0.1038} to \texttt{0.3830}, and
positive-pair latent L\textsubscript{2} energy from \texttt{149.25} to
\texttt{10.57} versus \texttt{13.02} random-pair), and a 2-layer
latent-conditioned typed edit decoder reranked by
\texttt{\textbar{}\textbar{}z\_hat\_y\ -\ E\_ema(y\_cand)\textbar{}\textbar{}\_2\^{}2}.

\begin{itemize}
\tightlist
\item
  \textbf{Disjoint synthetic, 151-ladder, and 40-residual accuracy
  versus model scale:} On \texttt{200} held-out disjoint synthetic Form
  defects (unseen \texttt{shape\_\allowbreak{}hash}es), \texttt{FormJEPA-\allowbreak{}v1} admits
  \textbf{\texttt{156/200} (\texttt{78.0\%})} at \texttt{k\ =\ 1} and
  \textbf{\texttt{177/200} (\texttt{88.5\%})} at \texttt{k\ =\ 16}. On
  the 151-task ladder (plain arm), standalone \texttt{FormJEPA-\allowbreak{}v1}
  admits \texttt{27/151} (\texttt{k\ =\ 1}) and \texttt{57/151}
  (\texttt{k\ =\ 16}), lifting the combined deterministic ladder
  (\texttt{111/151}) plus \texttt{FormJEPA-\allowbreak{}v1\ (\allowbreak{}k\ =\ 16)} to
  \textbf{\texttt{138/151} (\texttt{91.4\%})} with zero frontier calls.
  On the 40 open residual tasks (H15), \texttt{FormJEPA-\allowbreak{}v1} admits
  \textbf{\texttt{22/40} (\texttt{55.0\%})} at \texttt{k\ =\ 1} and
  \textbf{\texttt{27/40} (\texttt{67.5\%})} at \texttt{k\ =\ 16}
  (\texttt{25} restored + \texttt{2} behaviour-equal = \texttt{27/27}
  behaves-the-same, \texttt{5.33} mean \texttt{rce\ check} calls/task),
  matching local \texttt{Qwen\ 3.\allowbreak{}6\ 35B-\allowbreak{}A3B} (\texttt{27/40}) on
  structured grammar, length-method, and capability repairs
  (\texttt{24/25} \texttt{W-SYN-1}, \texttt{2/5} \texttt{W-LEN-1},
  \texttt{1/1} \texttt{W-WORLD-1}), while closing \texttt{0/9}
  multi-module cryptographic stdlib migrations (\texttt{W-NAME-1} /
  \texttt{W-TY-1} across \texttt{hpke}, \texttt{mls}, \texttt{pq},
  \texttt{pqxdh}, \texttt{energy}), compound \texttt{W-\allowbreak{}CAP-\allowbreak{}3,\allowbreak{}W-\allowbreak{}WORLD-\allowbreak{}1},
  and vacuous-example synthesis tasks where Gear-2 frontier proposers
  remain required.
\item
  \textbf{Inference \texttt{reported\_\allowbreak{}j}, 3-Gear Cascade, and honest
  H5/H11 diagnostics:} On AMD Ryzen AI Max+ 395 Linux \texttt{powercap}
  RAPL counters (\texttt{measured\_\allowbreak{}j\ =\ None},
  \texttt{board\_\allowbreak{}metered\ =\ false},
  \texttt{derived\_\allowbreak{}marginal\_\allowbreak{}pkg\_\allowbreak{}reported\_\allowbreak{}j\ =\ None} under Falsifier
  6 at \texttt{33.38\%} co-tenant package drift), per-proposal CPU
  inference costs \textbf{\texttt{58.65\ mJ}} RAPL core
  (\texttt{1.194\ J} gross package, \texttt{10.15\ ms}) at
  \texttt{k\ =\ 1} and \textbf{\texttt{474.\allowbreak{}76\ mJ}} RAPL core
  (\texttt{6.736\ J} gross package, \texttt{67.51\ ms}) at
  \texttt{k\ =\ 16} --- \texttt{70.6×} to \texttt{1,436×} lower at
  \texttt{k\ =\ 1} (\texttt{12.5×} to \texttt{177.5×} at
  \texttt{k\ =\ 16}) than one local 35B candidate (\texttt{84.25\ J}).
  Inserting \texttt{FormJEPA-\allowbreak{}v1\ (\allowbreak{}k\ =\ 16)} as Gear 1 between
  \texttt{DNULLMM1} Gear 0 and \texttt{Claude\ Opus\ 5.\allowbreak{}5} Gear 2
  preserves \texttt{34/40} single-pass closure while cutting
  \texttt{Opus\ 5.5} calls from \texttt{40} to \texttt{13}
  (\texttt{\$1.38} to \texttt{\$0.81}), and on the 200-arrival
  \texttt{Zipf(\allowbreak{}0.\allowbreak{}85)} stream cuts Gear-2 frontier calls from
  \texttt{200} (stateless) and \texttt{48} (2-Gear) to
  \textbf{\texttt{23} calls (\texttt{8.70×} reduction; \texttt{152}
  Gear-0 hits, \texttt{25} Gear-1 admissions, \texttt{7} Gear-2
  admissions)}; with local \texttt{Qwen\ 3.6} as Gear 2,
  \texttt{Zipf(\allowbreak{}0.\allowbreak{}85)} 35B calls drop to \texttt{49} (\texttt{4.08×}
  reduction) and admitted arrivals rise from \texttt{141/200} to
  \texttt{154/200}. Uncalibrated negative latent L\textsubscript{2}
  distance alone achieves AUROC \textbf{\texttt{0.618}} across
  \texttt{925} beam candidates (H5 inconclusive: above \texttt{0.60}
  disproof floor, below \texttt{0.85} support threshold) and rises
  across cumulative edit distances \texttt{0..4}
  (\texttt{0.\allowbreak{}00,\allowbreak{}\ 0.\allowbreak{}99,\allowbreak{}\ 2.\allowbreak{}26,\allowbreak{}\ 3.\allowbreak{}27,\allowbreak{}\ 3.\allowbreak{}89}, then \texttt{3.72} at
  \texttt{d\ =\ 5}; \texttt{18/50\ =\ 36.0\%\ \textless{}\ 80\%} strict
  5-step monotonicity on H11), showing that repair accuracy here comes
  from conditioning the typed edit decoder on
  \texttt{(\allowbreak{}z\_\allowbreak{}x,\allowbreak{}\ z\_\allowbreak{}hat\_\allowbreak{}y)} rather than smooth uncalibrated latent
  energy alone.
\item
  \textbf{151-task \texttt{rce\ repair} rung \texttt{f} promotion
  verdict:} Running \texttt{rce\ repair} in \texttt{observe} mode
  (\texttt{bound\ 2000}, given each task's known-good original and
  outside proposals file) across all 151 ladder tasks and evaluating
  \texttt{rce\ repair\ criterion} against the pinned 3 November 2026
  registration (\texttt{sha256:\ abf7f131\allowbreak{}f240e3f9.\allowbreak{}.\allowbreak{}.\allowbreak{}}) records
  \texttt{116/151} exact closes by rungs \texttt{a..e} and \texttt{35}
  open tasks (\texttt{25} unparsable \texttt{W-SYN-1} files skipped
  under rule RF-2, \texttt{10} open parsing tasks judged with \texttt{0}
  missing coverage). Under combined external proposals
  (\texttt{FormJEPA-\allowbreak{}v1} + recorded H15 repairs, \texttt{91} proposals
  judged), rung \texttt{f} observes \textbf{\texttt{3} admitted tasks}
  (\texttt{hpke}, \texttt{mls\_\allowbreak{}key\_\allowbreak{}schedule}, \texttt{energy}),
  \textbf{\texttt{3/3} (\texttt{100.0\%})} pass the \texttt{W-MUT-1}
  mutation gate, and \textbf{\texttt{0} contradict the known-good
  original} (\texttt{3/3} restored); because
  \texttt{3\ /\ 119\ =\ 2.5\%\ \textgreater{}\ 2.0\%},
  \texttt{rce\ repair\ criterion} returns
  \textbf{\texttt{holds:\ false}} (\texttt{exit\ =\ 1}), keeping rung
  \texttt{f} in \texttt{observe} mode.
\end{itemize}


\begin{center}
\openiefig{fig/rce-jepa-fig10-three-gear-and-ebm.png}{0.94}{https://research.openie.dev/figures/rce-jepa/three-gear-and-ebm.svg}
\end{center}


\emph{Figure 10. Measured Three-Gear residual cascade and
\texttt{FormEBM-\allowbreak{}v2} latent energy diagnostics at \texttt{openIE-\allowbreak{}dev/\allowbreak{}rce}
commit \texttt{1762fb4b} {[}B35{]}. Left: deterministic ladder rungs
\texttt{a..e} (\texttt{111/151}), Gear 0 \texttt{DNULLMM1}
(\texttt{40/40} warm zero-call replay), Gear 1 \texttt{FormJEPA-\allowbreak{}v1} +
\texttt{FormEBM-\allowbreak{}v2}, and Gear 2 \texttt{H15-\allowbreak{}Iterative} gated by
\texttt{rce\ check} + \texttt{W-MUT-1} + \texttt{no\_\allowbreak{}spec\_\allowbreak{}deletion}
(\texttt{200/200} admitted on \texttt{Zipf(\allowbreak{}0.\allowbreak{}85)}, \texttt{16.67×} fewer
Gear-2 episodes). Right: unseen-shape H5 ROC curve
(\texttt{AUROC\ =\ 0.\allowbreak{}8922}) and unseen-chain H11 5-step cumulative
defect energy curve (\texttt{93.0\%} strict monotonicity, IQR shaded).
Interactive explorer:
\href{https://research.openie.dev/living/rce-jepa/}{/living/rce-jepa/} ·
.}

\textbf{Contrastive multi-step margin energy scorer
(\texttt{FormEBM-\allowbreak{}v2}, H5 \& H11) and \texttt{40/40} bounded
checker-feedback residual closure (\texttt{H15-\allowbreak{}Iterative} + Gear-0 v2),
measured {[}B35{]}.} Pre-registered before corpus generation, weight
initialization, and API execution
(\texttt{rce-\allowbreak{}jepa-\allowbreak{}form-\allowbreak{}ebm-\allowbreak{}v2-\allowbreak{}2026-\allowbreak{}11-\allowbreak{}23},
\texttt{sha256:\ 128f17e6\allowbreak{}9361ea89.\allowbreak{}.\allowbreak{}.\allowbreak{}};
\texttt{rce-\allowbreak{}jepa-\allowbreak{}h15-\allowbreak{}iterative-\allowbreak{}2026-\allowbreak{}11-\allowbreak{}23},
\texttt{sha256:\ a7774d3b\allowbreak{}d03ae450.\allowbreak{}.\allowbreak{}.\allowbreak{}}) against \texttt{openIE-\allowbreak{}dev/\allowbreak{}rce}
commit \texttt{1762fb4b} (Linux x86\_64 binary
\texttt{sha256:\ 33146449\allowbreak{}6009d98a.\allowbreak{}.\allowbreak{}.\allowbreak{}}), two targeted evaluations tested
whether (A) replacing uncalibrated coordinate distance with a
contrastive + pairwise edit-depth margin energy head on the identical
strictly disjoint \texttt{1,387}-Form split closes H5 and H11, and (B)
supplying public \texttt{stdlib/\allowbreak{}lib} module headers plus bounded
deterministic checker diagnostics (\texttt{max\_\allowbreak{}turns\ =\ 5}, zero
access to \texttt{original}) closes the six remaining residual tasks
from single-shot H15.

\begin{itemize}
\item
  \textbf{\texttt{FormEBM-\allowbreak{}v2} architecture, tokenizer fix, and
  unseen-shape H5 / H11 results (\texttt{137} unseen test Form shapes,
  \texttt{0} overlap):} Two structural causes explained why
  \texttt{FormJEPA-\allowbreak{}v1} reached only \texttt{0.618} AUROC on H5 and
  \texttt{36.0\%} strict 5-step monotonicity on H11: its single-file
  S-expression tokenizer mapped both live unused bindings
  (\texttt{let\ x}, refused by \texttt{W-\allowbreak{}UNUSED-\allowbreak{}1}) and
  underscore-prefixed bindings (\texttt{let\ \_x}, admitted) to the same
  \texttt{\textless{}ID:k\textgreater{}} slot while shifting subsequent
  \texttt{\textless{}ID:j\textgreater{}} indices whenever an unbound
  identifier appeared, and its coordinate Smooth-L\textsubscript{1} loss
  carried no contrastive validity or edit-depth margin.
  \texttt{FormEBM-\allowbreak{}v2} (\texttt{5,851,649} trainable parameters,
  checkpoint \texttt{sha256:\ 2204a0b9\allowbreak{}cbd0d96c.\allowbreak{}.\allowbreak{}.\allowbreak{}}; \texttt{857.2\ s}
  training on AMD Ryzen AI Max+ 395 / Radeon 8060S APU at
  \texttt{5,\allowbreak{}375.\allowbreak{}0\ J} core RAPL and \texttt{40,\allowbreak{}602.\allowbreak{}4\ J} APU package
  RAPL) introduces reference-anchored \texttt{PairFormTokenizer}
  (\texttt{504} active tokens, distinguishing underscore-prefixed
  bindings, newly introduced identifier slots, and modulo-32 integer
  classes with purely syntactic diff tags and zero checker calls), a
  5-layer shared Siamese Transformer, a 1-layer cross-attention RMS
  difference block, and a softplus energy head trained on \texttt{5,955}
  validity-margin triples (\texttt{m\ =\ 1.0}) and \texttt{600} 5-step
  cumulative corruption chains (\texttt{γ\ =\ 0.\allowbreak{}50(\allowbreak{}j\ −\ i)}) strictly
  over the \texttt{1,106} training Forms (VICReg effective rank
  expanding from \texttt{26.48} to \texttt{55.70} on \texttt{z\_ref} and
  \texttt{46.44} to \texttt{74.72} on \texttt{z\_cand}). On the
  \texttt{137} unseen held-out test Form shapes:
  \textless ul\textgreater{} \textless li\textgreater{}\textbf{H5
  (\texttt{N\ =\ 200} unseen test beams, \texttt{3,064} candidates):}
  Scalar energy reaches \textbf{AUROC
  \texttt{0.8922\ \textgreater{}=\ 0.85}} (vs \texttt{0.6180} for
  \texttt{FormJEPA-\allowbreak{}v1} and \texttt{0.7208} for uncalibrated Siamese
  L\textsubscript{2}), lifting top-1 (\texttt{k\ =\ 1} without invoking
  \texttt{rce\ check} on beam items) valid selection from
  \texttt{22.5\%} (\texttt{45/200} unreranked) and \texttt{25.0\%}
  (\texttt{50/200} Siamese L\textsubscript{2}) to
  \textbf{\texttt{84.05\%} on the \texttt{163} solvable beams
  (\texttt{137/163}; \texttt{68.50\%} raw across all \texttt{200}
  beams)} and cutting mean \texttt{rce\ check} invocations to first
  valid Form by \textbf{\texttt{2.39×}} (\texttt{3.577} to
  \texttt{1.497} checks/beam). Ten-thread CPU scoring costs
  \textbf{\texttt{47.09\ mJ}} core RAPL per candidate
  (\texttt{6.30\ ms}; \texttt{752.\allowbreak{}03\ mJ} core and \texttt{11.88\ J}
  gross package per 16-item beam, with
  \texttt{derived\_\allowbreak{}marginal\_\allowbreak{}pkg\_\allowbreak{}reported\_\allowbreak{}j\ =\ None} under Falsifier
  6 at \texttt{13.7\%} idle drift).
\item
  \textbf{H11 (\texttt{N\ =\ 100} unseen 5-step cumulative defect
  chains):} Strict 5-step monotonicity
  (\texttt{E\_0\ \textless{}\ E\_1\ \textless{}\ E\_2\ \textless{}\ E\_3\ \textless{}\ E\_4\ \textless{}\ E\_5})
  rises from \texttt{36.0\%} (\texttt{FormJEPA-\allowbreak{}v1}) and \texttt{18.0\%}
  (Siamese L\textsubscript{2}) to \textbf{\texttt{93.0\%}
  (\texttt{93/100\ \textgreater{}=\ 80\%})}, with mean Spearman rank
  correlation \textbf{\texttt{ρ\ =\ 0.\allowbreak{}9931}} (median \texttt{1.0000})
  and mean energy
  \texttt{{[}0.0007,\ 2.36,\ 5.16,\ 7.99,\ 12.21,\ 17.25{]}} across
  \texttt{d\ =\ 0.\allowbreak{}.\allowbreak{}5} (zero \texttt{d\ =\ 5} rollover).
\item
  \textbf{Candid domain boundaries:} On unconstrained 5-step random
  walks (\texttt{N\ =\ 50}) that may overwrite an already-broken line
  twice without adding a new AST defect, mean energy still climbs
  \texttt{0.\allowbreak{}0007\ →\ 7.\allowbreak{}07} (mean \texttt{ρ\ =\ 0.\allowbreak{}8513}, median
  \texttt{0.9429}), while strict step-by-step inequality is
  \texttt{38.0\%} (\texttt{46.0\%} non-decreasing). Zero-shot transfer
  to the OOD \texttt{40}-task git-history residual lifts within-task
  \texttt{k\ =\ 1} typed-edit admissions over \texttt{FormJEPA-\allowbreak{}v1}
  (\texttt{22/\allowbreak{}40\ →\ 25/\allowbreak{}40}, beam capacity \texttt{30/40}), whereas
  pooled cross-task AUROC inverts to \texttt{0.2630} (\texttt{0.2585}
  across \texttt{332} multi-proposer candidates) because single-line
  edit-depth training assigns higher raw energy to large valid
  multi-line standard-library migrations than to 1-token syntax edits.

\begin{verbatim}
</li>
<li>**`H15-Iterative` bounded checker feedback + public `stdlib/lib` headers (`40 / 40` closed, `0` contradictions, `40 / 40` Gear-0 replay):** Supplying public declaration signatures (`pub const`, `pub struct`, `pub kind`, `pub fn` with bodies stripped, strictly excluding the task file itself) for imported `stdlib/lib` modules (`crypto.cat`, `crypto.unhex`, `crypto.tohex`, `core.bytes`, `pq`) plus verbatim deterministic `rce check` + `W-MUT-1` diagnostics up to `max_turns = 5` (with zero access to `original`) lifts admissions from `34/40` single-shot to **`40/40` (`100.0%`)** on both **`Claude Opus 5.5`** (`38` at turn 1, `2` at turn 2; `42` calls, `$1.8104`, `34` restored + `5` behaviour-equal = `39` behaves-the-same, **`0` contradictions**) and **`Claude Sonnet 5.5`** (`38` at turn 1, `1` at turn 2, `1` at turn 3; `43` calls, **`$1.1122`**, `35` restored + `4` behaviour-equal, **`0` contradictions**), and to **`39/40` (`97.5%`)** on **`xAI Grok 4.20 Reasoning`** (`35` at turn 1, `2` at turn 2, `2` at turn 3; `50` calls, `$4.0555`, **`0` contradictions**). All four previously unsolved cryptographic migrations (`stdlib/lib/hpke.rce`, `stdlib/lib/mls.rce`, `stdlib/lib/pq.rce`, `programs/ecosystem/pqxdh.rce`) were **exact-hash `restored`** (at turn 1 on `Opus 5.5` and `Sonnet 5.5`; turns 1–3 on `Grok 4.20`), while both historical pre-`W-MUT-1` vacuous-example tasks (`57_float_vectors.rce`, `washer.rce`) closed with **`behaviour-equal`** under `rce equiv fuel 2000` (only `washer.rce` remained open on `Grok 4.20` after 5 turns). Crystallizing all 40 verified repairs into `DNULLMM1` Gear-0 v2 (`sha256: 80cf265038bcd5eb...`) replays **`40/40` (`100.0%`) at `0` model calls**, and in the updated 3-Gear `Zipf(0.85)` stream (`M = 200` arrivals) admits **`200/200` (`100.0%`)** of arrivals (`163` Gear-0 hits + `25` local Gear-1 admissions + `12` Gear-2 episodes / `14` API turns at `$0.9822`, a **`16.67×` episode reduction** vs stateless `200` calls).</li>
\end{verbatim}

  \textless/ul\textgreater{}
\end{itemize}

\textbf{Span-normalized two-head decomposed scorer (\texttt{FormEBM-\allowbreak{}v3},
H5, H11, and zero-shot OOD H15 cross-span transfer), measured
{[}B36{]}.} Pre-registered before synthetic multi-span pair generation,
weight initialization, and evaluation
(\texttt{rce-\allowbreak{}jepa-\allowbreak{}form-\allowbreak{}ebm-\allowbreak{}v3-\allowbreak{}2026-\allowbreak{}11-\allowbreak{}23},
\texttt{sha256:\ cb7ee1ac\allowbreak{}8dc5c594.\allowbreak{}.\allowbreak{}.\allowbreak{}}) against \texttt{openIE-\allowbreak{}dev/\allowbreak{}rce}
commit \texttt{1762fb4b} (Linux x86\_64 binary
\texttt{sha256:\ 386be0c2\allowbreak{}6389550e.\allowbreak{}.\allowbreak{}.\allowbreak{}}), \texttt{FormEBM-\allowbreak{}v3}
(\texttt{6,052,802} trainable parameters, pure-\texttt{float32} ROCm
\texttt{Radeon\ 8060S} checkpoint \texttt{sha256:\ 71e642f1\allowbreak{}bb0b9c06.\allowbreak{}.\allowbreak{}.\allowbreak{}};
\texttt{2,\allowbreak{}528.\allowbreak{}3\ s} training at \texttt{13,\allowbreak{}312.\allowbreak{}3\ J} core RAPL and
\texttt{48,\allowbreak{}422.\allowbreak{}2\ J} APU package RAPL) tests whether decoupling defect
validity from cumulative edit depth over the identical strictly disjoint
\texttt{1,387}-Form split (\texttt{1,106} train, \texttt{144}
validation, \texttt{137} unseen test shapes; \texttt{0} overlap with the
151 ladder or 40 H15 tasks, and \texttt{0} \texttt{rce\ check} calls
during scoring) resolves \texttt{FormEBM-\allowbreak{}v2}'s cross-task OOD edit-size
inversion (\texttt{0.2585}) without regressing H5 or H11:

\begin{itemize}
\tightlist
\item
  \textbf{Alignment-aware tokenization and two-head decomposition:}
  \texttt{PairFormTokenizerv3} (\texttt{max\_\allowbreak{}seq\_\allowbreak{}len\ =\ 384},
  \texttt{560} active tokens) replaces raw line-index comparison with
  line-shift-invariant sequence alignment, separates lexically
  bound/imported new identifiers
  (\texttt{\textless{}BOUND\_NEW\_ID:0..31\textgreater{}}) from unbound
  free symbols
  (\texttt{\textless{}UNBOUND\_NEW\_ID:0..15\textgreater{}}), and
  computes explicit normalized diff fraction
  \texttt{ρ\_\allowbreak{}diff\ =\ Δ\_\allowbreak{}tok\ /\allowbreak{}\ max(\allowbreak{}1,\allowbreak{}\ L\_\allowbreak{}cand)}. Over the shared
  5-layer Siamese encoder and cross-difference block, a span-invariant
  defect head
  \texttt{E\_\allowbreak{}def(\allowbreak{}x,\allowbreak{}\ y,\allowbreak{}\ r)\ =\ softplus(\allowbreak{}g\_\allowbreak{}def(\allowbreak{}feats\_\allowbreak{}def))} (trained
  with within-beam plus all-pairs cross-task batch margin
  \texttt{m\ =\ 2.5} and L\textsubscript{2} positive anchor across
  \texttt{4,983} pairs including \texttt{219} \texttt{k\ =\ 2.\allowbreak{}.\allowbreak{}10}
  multi-span valid rewrites strictly from \texttt{train\_\allowbreak{}forms}) is
  separated from a cumulative corruption-depth head
  \texttt{E\_\allowbreak{}depth(\allowbreak{}y\_\allowbreak{}0,\allowbreak{}\ y\_\allowbreak{}d)\ =\ softplus(\allowbreak{}g\_\allowbreak{}depth(\allowbreak{}feats\_\allowbreak{}depth))},
  combined via
  \texttt{E\_\allowbreak{}total\ =\ E\_\allowbreak{}def\ +\ 0.\allowbreak{}15\ ·\ E\_\allowbreak{}depth\ /\allowbreak{}\ log2(\allowbreak{}2\ +\ Δ\_\allowbreak{}tok)}
  (VICReg effective rank \texttt{42.79} on \texttt{z\_ref} and
  \texttt{57.93} on \texttt{z\_cand}, coordinate standard deviation
  \texttt{0.3332\ /\ 0.3113\ \textgreater{}=\ 0.25}).
\item
  \textbf{Strengthened unseen-shape H5 (\texttt{0.9437}) and
  unseen-chain H11 (\texttt{94.0\%} / \texttt{92.0\%}):} On the
  \texttt{200} unseen held-out test beams (\texttt{3,064} candidates),
  \texttt{FormEBM-\allowbreak{}v3} lifts \textbf{H5 AUROC from \texttt{0.8922}
  (\texttt{FormEBM-\allowbreak{}v2}) to \texttt{0.9437\ \textgreater{}=\ 0.85}}
  (\texttt{0.9386} on \texttt{E\_def} alone, reaching
  \texttt{TPR\ =\ 78.\allowbreak{}9\%} at \texttt{FPR\ =\ 4.\allowbreak{}6\%}), achieves
  \textbf{\texttt{83.44\%} (\texttt{136/163}, \texttt{E\_def}) and
  \texttt{80.98\%} (\texttt{132/163}, \texttt{E\_total}) solvable top-1
  accuracy} (both \texttt{\textgreater{}=\ 80.0\%}; \texttt{1.601} /
  \texttt{1.669} mean \texttt{rce\ check} calls vs \texttt{3.577}
  unreranked), and lifts \textbf{H11 strict 5-step monotonicity on
  \texttt{100} unseen corruption chains to \texttt{94.0\%} on
  \texttt{E\_depth} (Spearman \texttt{ρ\ =\ 0.\allowbreak{}9966}, median
  \texttt{1.0000}) and \texttt{92.0\%} on \texttt{E\_total}
  (\texttt{ρ\ =\ 0.\allowbreak{}9846}, median \texttt{1.0000})}. Ten-thread CPU
  scoring at \texttt{max\_\allowbreak{}seq\_\allowbreak{}len\ =\ 384} costs \texttt{232.\allowbreak{}45\ mJ}
  core RAPL per candidate (\texttt{32.55\ ms}, \texttt{362×} lower than
  local 35B generation;
  \texttt{derived\_\allowbreak{}marginal\_\allowbreak{}pkg\_\allowbreak{}reported\_\allowbreak{}j\ =\ None} at
  \texttt{37.7\%} co-tenant package drift).
\item
  \textbf{Zero-shot OOD H15 inversion resolved
  (\texttt{0.\allowbreak{}2585\ →\ 0.\allowbreak{}5278}), and candid cross-module boundary:}
  Transferred zero-shot to the \texttt{40}-task git-history residual
  (\texttt{303} multi-proposer candidates including all \texttt{40/40}
  closed multi-hundred-line \texttt{hpke}, \texttt{mls}, \texttt{pq},
  and \texttt{pqxdh} migrations), \texttt{FormEBM-\allowbreak{}v3} lifts pooled
  cross-task AUROC from \texttt{0.2585} (\texttt{FormEBM-\allowbreak{}v2}) to
  \textbf{\texttt{0.5278}} (\texttt{+0.2693}, clearing Falsifier F4
  \texttt{\textgreater{}\ 0.50}; \textbf{\texttt{0.6018}} on large-span
  valid rewrites with \texttt{Δ\_tok\ \textgreater{}=\ 10} vs short
  \texttt{0/1}-token invalid candidates, and \texttt{0.\allowbreak{}2630\ →\ 0.\allowbreak{}4985}
  on the \texttt{286}-item typed-edit pool), while reaching
  \textbf{\texttt{27/40} multi-proposer and \texttt{24/40} typed-edit
  within-task \texttt{k\ =\ 1} admissions under \texttt{E\_def}}
  (\texttt{25/40} and \texttt{22/40} under \texttt{E\_total}, cutting
  mean typed-beam checks from \texttt{5.55} to \texttt{5.10}). On the
  279-line \texttt{stdlib/\allowbreak{}lib/\allowbreak{}hpke.\allowbreak{}rce} migration
  (\texttt{Δ\_\allowbreak{}tok\ =\ 146}), valid energy drops to
  \texttt{E\_\allowbreak{}total\ =\ 0.\allowbreak{}9645} (\texttt{E\_\allowbreak{}def\ =\ 0.\allowbreak{}8826}), yet across
  unseen multi-module cryptographic APIs (\texttt{mls.rce},
  \texttt{pq.rce}, \texttt{pqxdh.rce}) and vacuous-example obligations
  (\texttt{washer.\allowbreak{}rce}), zero-shot pooled calibration plateaus near
  \texttt{0.53} because single-file syntax without imported
  \texttt{stdlib/\allowbreak{}lib} type signatures or exact \texttt{rce\ trace}
  rollouts cannot verify cross-module call contracts.
\end{itemize}

\subsection{Falsifiers RJ-F1 to RJ-F6 for the premise
itself}\label{falsifiers-rj-f1-to-rj-f6-for-the-premise-itself}

Each falsifier is an observation that would defeat part of the premise;
several are stated in terms of the hypotheses above (RJ-F3 on H3, RJ-F4
on H11, RJ-F6 on H3 and H10).

{\footnotesize\setlength{\tabcolsep}{3.8pt}\renewcommand{\arraystretch}{1.15}
\begin{longtable}[]{@{}
  >{\raggedright\arraybackslash}p{(\linewidth - 4\tabcolsep) * \real{0.0741}}
  >{\raggedright\arraybackslash}p{(\linewidth - 4\tabcolsep) * \real{0.4444}}
  >{\raggedright\arraybackslash}p{(\linewidth - 4\tabcolsep) * \real{0.4815}}@{}}
\toprule\noalign{}
\begin{minipage}[b]{\linewidth}\raggedright
\textbf{ID}
\end{minipage} & \begin{minipage}[b]{\linewidth}\raggedright
\textbf{Observation}
\end{minipage} & \begin{minipage}[b]{\linewidth}\raggedright
\textbf{What it defeats}
\end{minipage} \\
\midrule\noalign{}
\endhead
\bottomrule\noalign{}
\endlastfoot
RJ-F1 & Lift coverage plateaus below 80\% on any of the five target
languages after the construct list is exhausted & The claim that all
code can be emitted from RCE, and with it the scope of autonomous coding
over Forms \\
RJ-F2 & Writing σ(I) for the 7 reference applications costs more
person-hours than writing their Forms by hand & The claim that closure
removes, rather than moves, the specification cost \\
RJ-F3 & H3 fails while the token arm passes at the same gate & Latent
proposal over Form space; the gate still stands, the JEPA half does
not \\
RJ-F4 & H11 fails and search does not reach accepted Forms within the
beam budget & Energy-guided planning in Form space (not triggered on
unseen cumulative 1..5-defect chains: \texttt{FormEBM-\allowbreak{}v2} reaches
\texttt{93.0\%} and \texttt{FormEBM-\allowbreak{}v3} reaches \texttt{94.0\%}
\texttt{E\_depth} / \texttt{92.0\%} \texttt{E\_total} strict
monotonicity with \texttt{0.9437} H5 AUROC; cross-module signature/trace
OOD energy calibration beyond \texttt{0.5278} remains open) \\
RJ-F5 & The decision head's calibration never beats a frequency table
over shape hashes & The decision-model leg; it adds nothing over
promoted-handler statistics \\
RJ-F6 & A token-space model trained with execution traces, of the CWM
class, closes the residual of 40 tasks from the plain arm of H14 run 3,
at equal energy and with the same gate as the acceptance rule & The
claim that the Form quotient adds something for proposers during
training/proposal (note: for Gear-0 reuse, \texttt{rce\ fmt\ nameless}
already recovers \texttt{7/9} alpha-renamings at 0 calls;
trained-proposer parity against H3/H10 awaits a trained JEPA arm) \\
\end{longtable}
}

The thesis is disproved overall if H4 fails (no joule advantage per
accepted program) or H6 fails (the Form cannot express the reference
apps). Either alone is sufficient.

\subsection{What will be delivered}\label{what-will-be-delivered}

Scope, organised as parallel workstreams that meet at the evaluation
harness. Every deliverable either produces Forms, produces models that
propose Forms, or measures the result.

\begin{itemize}
\tightlist
\item
  \textbf{W1. Form corpus and dataset card.} Forms paired with intent
  statements and checker verdicts, from hand-written Forms,
  reference-app Forms and checker-verified mutations; the card covers
  provenance, licence, size distribution, negative generation, known
  gaps and contamination controls; canonical serialization for H1.
\item
  \textbf{W2. Encoder and predictor.} Intent encoder, EMA Form target
  encoder, predictor in representation space and energy head E(intent,
  P), with anti-collapse measures, an embedding-space versus token-space
  loss ablation, and model cards with measured training energy.
\item
  \textbf{W3. Checker-constrained decoder.} Grammar-constrained
  candidate Forms ranked by energy and submitted to the checker,
  requires-change diagnostics, each naming its rule, fed back under a
  fixed retry budget; the checker is the sole acceptance authority with
  no bypass path to the engine.
\item
  \textbf{W4. Intent-to-Form benchmark.} Unit and composite tasks, the
  reference apps (game engine, text editor, video player, chat app,
  browser, terminal) and RCEdb; each task has intent, reference Form,
  replay traces and an equivalent source-language task for the baseline;
  held-out split.
\item
  \textbf{W5. Evaluation harness with energy in joules.} Both arms on
  identical hardware and tasks;
  \href{mailto:pass@1}{\nolinkurl{pass@1}},
  \href{mailto:pass@k}{\nolinkurl{pass@k}}, ρ, C\textsubscript{check},
  J\textsubscript{acc}, K\textsubscript{acc}, retries per acceptance,
  energy-ranking AUROC; rootless counters (\texttt{reported\_\allowbreak{}j}),
  idle-drift Falsifier 6 gating, raw logs kept, unavailable cloud
  counters reported as \texttt{not\_\allowbreak{}available} rather than estimated.
\item
  \textbf{W6. Research article.} This case expanded with results for H1
  to H15, including negative results, figures, tables and
  reproducibility instructions.
\item
  \textbf{W7. Open artifacts.} Public: the Form specification, the
  reference checker and its emitters, benchmark tasks and traces,
  harness, dataset and model cards, and weights where licences allow. A
  number is reported only when everything needed to reproduce it is
  public.
\item
  \textbf{W8. Demo.} Stated intent to predicted Form to checker verdict
  (including a shown requires-change verdict) to replay and emitted code
  in two languages by the reference checker, with an on-screen joule
  reading.
\end{itemize}

\section{Status}\label{status}

Research topic open, updated 23 November 2026 against
\texttt{openIE-\allowbreak{}dev/\allowbreak{}rce} commit \texttt{1762fb4b} (20 November 2026;
\texttt{5,537} Lean 4 theorem constants across \texttt{64} modules,
theorems T1 to T63, \texttt{2,\allowbreak{}221/\allowbreak{}2,\allowbreak{}236} conformance tests passing with
\texttt{0} failures, \texttt{837} Rust checker unit tests; interactive
observatory at
\href{https://research.openie.dev/living/rce-jepa/}{/living/rce-jepa/}).
Phase 0 deterministic floors are measured: H1 is supported on the
hand-written population at a median rho of \texttt{9.34} over 185
programs (W1 version 1 gave \texttt{9.31} over 173); H8 is supported
across \texttt{10,500} seeded event streams on six targets with
\texttt{0} mismatches (\texttt{rce/\allowbreak{}bench/\allowbreak{}determinism/\allowbreak{}}); H7's fuzzed arm
found \texttt{0} invariant violations across \texttt{1,275}
checker-accepted mutants (\texttt{rce/\allowbreak{}bench/\allowbreak{}adversarial/\allowbreak{}}); and H14 is
supported across three registered runs of \texttt{rce\ repair} {[}B30,
B31, B32{]}, closing \texttt{111} of \texttt{151} tasks
(\texttt{73.5\%}) at zero model calls in the plain arm of run 3 and
leaving \texttt{40} open git-history residual tasks. On 22--23 November
2026 {[}B33, B34, B35, B36{]}, Apple M5 Max and AMD Ryzen AI Max+ 395 /
Radeon 8060S \texttt{reported\_\allowbreak{}j} receipts
(\texttt{measured\_\allowbreak{}j\ =\ None}, \texttt{board\_\allowbreak{}metered\ =\ false})
established \textbf{H2 supported}
(\texttt{C\_check\ /\ E\_model\ =\ 0.04904\ J\ /\ 84.25\ J\ =\ 0.000582\ \textless{}=\ 0.05},
\texttt{1,718×} cheaper; full 62-Form check+run+emit suite
\texttt{15.44\ J}), trained and evaluated \textbf{\texttt{FormJEPA-\allowbreak{}v1}}
(\texttt{8.42M} parameters: \texttt{177/200} unseen synthetic defects,
\texttt{22/40} at \texttt{k\ =\ 1} and \texttt{27/40} at
\texttt{k\ =\ 16} on H15 at \texttt{58.\allowbreak{}65–474.\allowbreak{}76\ mJ} RAPL core
\texttt{reported\_\allowbreak{}j}), contrastive \textbf{\texttt{FormEBM-\allowbreak{}v2}}
(\texttt{5.85M} parameters: \textbf{H5 AUROC
\texttt{0.8922\ \textgreater{}=\ 0.85}}, \texttt{84.05\%} solvable top-1
accuracy, \textbf{H11 \texttt{93.0\%\ \textgreater{}=\ 80\%}} strict
5-step monotonicity at Spearman \texttt{0.9931}, \texttt{47.09\ mJ}/\allowbreak{}cand
core \texttt{reported\_\allowbreak{}j}), and span-normalized decomposed
\textbf{\texttt{FormEBM-\allowbreak{}v3}} (\texttt{6.05M} parameters on the identical
disjoint \texttt{1,387}-Form split: \textbf{H5 AUROC
\texttt{0.9437\ \textgreater{}=\ 0.85}} with \texttt{83.44\%}
\texttt{E\_def} and \texttt{80.98\%} \texttt{E\_total} solvable top-1
accuracy, \textbf{H11 \texttt{94.0\%}} strict 5-step monotonicity on
\texttt{E\_depth} at Spearman \texttt{0.9966} and
\textbf{\texttt{92.0\%}} on \texttt{E\_total} at Spearman
\texttt{0.9846}, and zero-shot OOD H15 multi-proposer pooled AUROC
lifted from \texttt{0.2585} to \textbf{\texttt{0.5278}}, \texttt{0.6018}
on large-span valid vs short invalid rewrites, at
\texttt{232.\allowbreak{}45\ mJ}/cand core \texttt{reported\_\allowbreak{}j}), closed
\textbf{\texttt{40/40} (\texttt{100.0\%})} of the open residual tasks in
\textbf{\texttt{H15-\allowbreak{}Iterative}} under bounded
\texttt{\textless{}=\ 5}-turn checker feedback + public
\texttt{stdlib/\allowbreak{}lib} signature headers on \texttt{Claude\ Opus\ 5.\allowbreak{}5}
(\texttt{\$1.81}) and \texttt{Claude\ Sonnet\ 5.\allowbreak{}5} (\texttt{\$1.11})
with \textbf{\texttt{0} contradictions} (\texttt{39/40} on
\texttt{xAI\ Grok\ 4.\allowbreak{}20\ Reasoning}), crystallized all
\textbf{\texttt{40/40} into \texttt{DNULLMM1} Gear-0 v2} (\texttt{0}
model calls, \texttt{0.325\ µs} / \texttt{1.77\ µJ} probe;
\texttt{200/200} \texttt{Zipf(\allowbreak{}0.\allowbreak{}85)} 3-Gear stream admissions at
\texttt{16.67×} fewer Gear-2 episodes), and audited 151-task
\textbf{\texttt{rce\ repair} rung \texttt{f}} promotion
(\texttt{3/119\ =\ 2.5\%\ \textgreater{}\ 2.0\%}, \texttt{3/3} passing
\texttt{W-MUT-1}, \texttt{0} contradictions, confirming rung \texttt{f}
remains in \texttt{observe} mode). Still open: full intent-to-Form
synthesis (H3), end-to-end joule parity (H4), cross-module
signature-and-trace OOD energy calibration beyond \texttt{0.5278} across
multi-module git-history rewrites, reference-app encoding (H6),
cost-table prediction (H9), quotient sample-efficiency (H10),
gate-receipt decision-head calibration (H12), discrete-versus-continuous
long-horizon planning (H13), diffusion/evolutionary H15 arms, and
external hardware-rail power metering (\texttt{measured\_\allowbreak{}j}).

\section{References}\label{references}

Each source carries a reading label. \textbf{Full:} read in full.
\textbf{Partial:} the full text was retrieved and searched for the
quoted passages, which were checked in their surrounding text, but the
whole work was not read. \textbf{Unavailable:} not obtained. ``PDF p.''
gives the page in the linked copy; Sec. gives the section.

\begin{itemize}
\tightlist
\item
  Assran et al.~2023. Self-Supervised Learning from Images with a
  Joint-Embedding Predictive Architecture. CVPR.
  \href{https://arxiv.org/abs/2301.08243}{arXiv:2301.08243}. Full.
\item
  Assran et al.~2025. V-JEPA 2.
  \href{https://arxiv.org/abs/2506.09985}{arXiv:2506.09985}. Full.
\item
  Backus et al.~1957. The FORTRAN Automatic Coding System. Western Joint
  Computer Conference.
  \href{https://archive.computerhistory.org/resources/text/Fortran/102663113.05.01.acc.pdf}{Computer
  History Museum copy}. Full.
\item
  Bird 1986. An Introduction to the Theory of Lists. Oxford PRG-56.
  \href{https://www.cs.ox.ac.uk/files/3378/PRG56.pdf}{PDF}. Full.
\item
  Brooks 1986. No Silver Bullet: Essence and Accident in Software
  Engineering.
  \href{https://worrydream.com/refs/Brooks_1986_-_No_Silver_Bullet.pdf}{PDF}
  (page numbers refer to this copy). Full.
\item
  Chen et al.~2021. Evaluating Large Language Models Trained on Code.
  \href{https://arxiv.org/abs/2107.03374}{arXiv:2107.03374}. Full.
\item
  Dijkstra. EWD 472, Guarded commands, non-determinacy and formal
  derivation of programs.
  \href{https://www.cs.utexas.edu/~EWD/transcriptions/EWD04xx/EWD472.html}{Transcription}.
  Full.
\item
  Dijkstra. EWD 667, On the foolishness of ``natural language
  programming''.
  \href{https://www.cs.utexas.edu/~EWD/transcriptions/EWD06xx/EWD667.html}{Transcription}.
  Full.
\item
  Ellis et al.~2021. DreamCoder. PLDI.
  \href{https://arxiv.org/abs/2006.08381}{arXiv:2006.08381};
  \href{https://www.cs.cornell.edu/~ellisk/documents/dreamcoder_with_supplement.pdf}{version
  with supplement}. Full (main text and supplement S1 to S4.8, PDF
  pp.~23-68, read in full).
\item
  Hubert et al.~2025. Olympiad-level formal mathematical reasoning with
  reinforcement learning (AlphaProof). Nature.
  \href{https://www.nature.com/articles/s41586-025-09833-y}{Nature}.
  Full.
\item
  Klein et al.~2009. seL4: Formal Verification of an OS Kernel. SOSP.
  \href{https://www.sigops.org/s/conferences/sosp/2009/papers/klein-sosp09.pdf}{PDF}.
  Full.
\item
  Leroy 2009. Formal verification of a realistic compiler. CACM.
  \href{https://xavierleroy.org/publi/compcert-CACM.pdf}{PDF}. Full.
\item
  Letouzey 2008. Extraction in Coq: an Overview. CiE.
  \href{https://www.irif.fr/~letouzey/download/letouzey_extr_cie08.pdf}{PDF}.
  Full.
\item
  Li et al.~2022. Competition-Level Code Generation with AlphaCode.
  \href{https://arxiv.org/abs/2203.07814}{arXiv:2203.07814}. Full.
\item
  Luccioni, Jernite, Strubell 2024. Power Hungry Processing. FAccT.
  \href{https://arxiv.org/abs/2311.16863}{arXiv:2311.16863}. Full.
\item
  OMG 2003. MDA Guide Version 1.0.
  \href{https://www.omg.org/mda/mda_files/MDA_Guide_Version1-0.pdf}{PDF}.
  Full.
\item
  Ouyang et al.~2023. An Empirical Study of the Non-determinism of
  ChatGPT in Code Generation (v1 title: LLM is Like a Box of
  Chocolates).
  \href{https://arxiv.org/abs/2308.02828}{arXiv:2308.02828}. Full.
\item
  Samsi et al.~2023. From Words to Watts. IEEE HPEC.
  \href{https://arxiv.org/abs/2310.03003}{arXiv:2310.03003}. Full.
\item
  Sun et al.~2024. Clover: Closed-Loop Verifiable Code Generation.
  \href{https://arxiv.org/abs/2310.17807}{arXiv:2310.17807}. Full.
\item
  Yang et al.~2011. Finding and Understanding Bugs in C Compilers
  (Csmith). PLDI.
  \href{https://users.cs.utah.edu/~regehr/papers/pldi11-preprint.pdf}{PDF}.
  Full.
\item
  Yang et al.~2024. AutoVerus.
  \href{https://arxiv.org/abs/2409.13082}{arXiv:2409.13082}. Full.
\item
  LeCun 2022. A Path Towards Autonomous Machine Intelligence, version
  0.9.2. \href{https://openreview.net/forum?id=BZ5a1r-kVsf}{OpenReview}.
  Full.
\item
  Simon 1963. Experiments with a Heuristic Compiler. Journal of the ACM
  10(4), 493-506. \href{https://doi.org/10.1145/321186.321192}{DOI
  10.1145/321186.321192}. Full.
\item
  Manna and Waldinger 1980. A Deductive Approach to Program Synthesis.
  ACM TOPLAS 2(1), 90-121.
  \href{https://doi.org/10.1145/357084.357090}{DOI
  10.1145/357084.357090}. Full.
\item
  van de Riet 1993. An overview and appraisal of the Fifth Generation
  Computer System project. Future Generation Computer Systems 9, 83-103.
  \href{https://doi.org/10.1016/0167-739X(93)90002-7}{DOI
  10.1016/0167-739X(93)90002-7}. Full.
\item
  Kemerer 1992. How the Learning Curve Affects CASE Tool Adoption. IEEE
  Software 9(3), 23-28. \href{https://doi.org/10.1109/52.136161}{DOI
  10.1109/52.136161}. Full.
\item
  Hutchinson, Whittle, Rouncefield and Kristoffersen 2011. Empirical
  Assessment of MDE in Industry. ICSE, 471-480.
  \href{https://doi.org/10.1145/1985793.1985858}{DOI
  10.1145/1985793.1985858}. Full.
\item
  Petre 2013. UML in practice. ICSE.
  \href{https://doi.org/10.1109/ICSE.2013.6606618}{DOI
  10.1109/ICSE.2013.6606618};
  \href{https://oro.open.ac.uk/35805/}{author copy, Open Research
  Online} (page numbers refer to this copy). Full.
\item
  Simon 1963. Experiments with a Heuristic Compiler. RAND RM-3588-PR.
  \href{https://archive.org/details/bitsavers_randiplRM3mpilerMay63_4935437}{Internet
  Archive}; a memorandum related to but distinct from the JACM article.
  Full (prose pp.~iii-63; program listings in Appendices A and B
  scanned).
\item
  Manna and Waldinger 1978. A Deductive Approach to Program Synthesis.
  Stanford STAN-CS-78-690.
  \href{http://i.stanford.edu/TR/CS-TR-78-690.html}{Report version},
  which predates the TOPLAS article. Full.
\item
  Huang, LeCun and Balestriero 2025. LLM-JEPA.
  \href{https://arxiv.org/abs/2509.14252}{arXiv:2509.14252}. Partial
  (abstract, repository README and run script).
\item
  Wei et al.~2025. SWE-RL.
  \href{https://arxiv.org/abs/2502.18449}{arXiv:2502.18449}. Partial
  (abstract and method summary).
\item
  Yang et al.~2025. SWE-smith.
  \href{https://arxiv.org/abs/2504.21798}{arXiv:2504.21798}. Partial
  (abstract and repository README).
\item
  Shansan Gong, Ruixiang Zhang, Huangjie Zheng, Jiatao Gu, Navdeep
  Jaitly, Lingpeng Kong and Yizhe Zhang 2025. DiffuCoder.
  \href{https://arxiv.org/abs/2506.20639}{arXiv:2506.20639}. Partial
  (abstract). Author list from arXiv metadata, not a PDF title page.
\item
  Inception Labs 2025. Mercury.
  \href{https://arxiv.org/abs/2506.17298}{arXiv:2506.17298}. Partial
  (abstract).
\item
  Wang, Pradel and Liu 2025. Are ``Solved Issues'' in SWE-bench Really
  Solved Correctly?
  \href{https://arxiv.org/abs/2503.15223}{arXiv:2503.15223}. ICSE
  (International Conference on Software Engineering) 2026. Partial
  (findings summary).
\item
  Becker, Rush, Barnes and Rein 2025. Measuring the Impact of Early-2025
  AI on Experienced Open-Source Developer Productivity.
  \href{https://arxiv.org/abs/2507.09089}{arXiv:2507.09089}. Partial
  (summary, and the abstract via arXiv metadata). Author list from arXiv
  metadata, not a PDF title page.
\item
  Niels Mündler, Jingxuan He, Hao Wang, Koushik Sen, Dawn Song and
  Martin Vechev 2025. Type-Constrained Code Generation with Language
  Models. \href{https://arxiv.org/abs/2504.09246}{arXiv:2504.09246}.
  Partial (headline result). Affiliation: the Swiss Federal Institute of
  Technology Zurich and the University of California, Berkeley, from the
  author list, not a PDF title page.
\item
  Grand et al.~2023. LILO.
  \href{https://arxiv.org/abs/2310.19791}{arXiv:2310.19791}. Partial
  (abstract).
\item
  Higher Order Company. HVM2.
  \href{https://raw.githubusercontent.com/HigherOrderCO/HVM/main/paper/HVM2.pdf}{PDF}.
  Partial (summary). The model is cited. Figures from that paper are not
  used.
\item
  DeepSeek 2025. DeepSeek-Prover-V2.
  \href{https://arxiv.org/abs/2504.21801}{arXiv:2504.21801}. Partial
  (arXiv metadata and abstract).
\item
  Moonshot and Numina 2025. Kimina-Prover Preview.
  \href{https://arxiv.org/abs/2504.11354}{arXiv:2504.11354}. Partial
  (abstract).
\item
  Princeton 2025. Goedel-Prover.
  \href{https://arxiv.org/abs/2502.07640}{arXiv:2502.07640}. Partial
  (abstract).
\item
  Scale AI. SWE-Bench Pro public leaderboard, page text fetched 6
  November 2026.
  \href{https://scale.com/leaderboard/swe_bench_pro_public}{scale.com/leaderboard/swe\_bench\_pro\_public}.
  Full (page text). The page prints no update date.
\item
  ByteDance Seed 2025. Multi-SWE-bench.
  \href{https://arxiv.org/abs/2504.02605}{arXiv:2504.02605}. Partial
  (abstract). No November 2026 leaderboard figure was retrieved, and
  none is stated.
\item
  Meta FAIR 2024. Large Concept Models.
  \href{https://arxiv.org/abs/2412.08821}{arXiv:2412.08821}. Partial
  (abstract).
\item
  Chollet, Knoop, Kamradt and Landers 2024. ARC Prize 2024 Technical
  Report. \href{https://arxiv.org/abs/2412.04604}{arXiv:2412.04604}.
  Partial (abstract).
\item
  Akyurek et al.~2024. Test-time training on ARC.
  \href{https://arxiv.org/abs/2411.07279}{arXiv:2411.07279}. Partial
  (abstract).
\item
  ARC Prize 2025 results page.
  \href{https://arcprize.org/competitions/2025/}{arcprize.org/competitions/2025/}.
  Full (page).
\item
  Lean FRO (Focused Research Organization). Year 1 roadmap, August 2023
  to July 2024.
  \href{https://lean-fro.org/about/roadmap/}{lean-fro.org/about/roadmap/}.
  Partial (deliverables list). Later-year roadmaps are not used.
\item
  Conal Elliott 2009, revised 2016. Denotational design with type class
  morphisms.
  \href{https://conal.net/papers/type-class-morphisms/}{conal.net/papers/type-class-morphisms/}.
  Partial (abstract page).
\item
  Radul and Sussman 2009. The Art of the Propagator.
  MIT-CSAIL-TR-2009-002.
  \href{https://dspace.mit.edu/handle/1721.1/44215}{dspace.mit.edu/handle/1721.1/44215}.
  Partial (abstract).
\end{itemize}


\end{document}
